% ============================================================
%  Number Theory — Full Course
%  Preamble + Chapters 1–2
% ============================================================

\documentclass[12pt,a4paper,openany]{book}
\usepackage[utf8]{inputenc}
\usepackage[T1]{fontenc}
\usepackage{lmodern}
\usepackage[margin=2.5cm]{geometry}
\usepackage{amsmath,amssymb,amsthm,mathrsfs,mathtools}
\usepackage{tikz}
\usetikzlibrary{arrows.meta,calc,positioning,patterns,decorations.markings,decorations.pathmorphing}
\usepackage{pgfplots}
\pgfplotsset{compat=1.18}
\usepackage[most]{tcolorbox}
\tcbuselibrary{theorems,skins,breakable}
\usepackage[colorlinks=true,linkcolor=blue!60!black,citecolor=green!50!black]{hyperref}
\usepackage{makeidx}
\makeindex
\usepackage{enumitem}
\usepackage{fancyhdr}
\usepackage{caption}
\usepackage{booktabs}
\usepackage{array}
\usepackage{microtype}
\usepackage{xcolor}

\pagestyle{fancy}
\fancyhf{}
\fancyhead[LE]{\slshape\leftmark}
\fancyhead[RO]{\slshape\rightmark}
\fancyfoot[C]{\thepage}

% ---- Theorem environments (amsthm) ----
\theoremstyle{definition}
\newtheorem{definition}{Definition}[chapter]
\newtheorem{example}[definition]{Example}
\newtheorem{exercise}{Exercise}[chapter]

\theoremstyle{plain}
\newtheorem{theorem}[definition]{Theorem}
\newtheorem{proposition}[definition]{Proposition}
\newtheorem{lemma}[definition]{Lemma}
\newtheorem{corollary}[definition]{Corollary}

\theoremstyle{remark}
\newtheorem{remark}[definition]{Remark}
\newtheorem{notation}[definition]{Notation}

% ---- tcolorbox visual wrapping ----
\tcolorboxenvironment{definition}{enhanced,breakable,colback=blue!5,colframe=blue!70!black,fonttitle=\bfseries,before skip=10pt,after skip=10pt,left=8pt,right=8pt}
\tcolorboxenvironment{theorem}{enhanced,breakable,colback=red!5,colframe=red!70!black,fonttitle=\bfseries,before skip=10pt,after skip=10pt,left=8pt,right=8pt}
\tcolorboxenvironment{proposition}{enhanced,breakable,colback=orange!5,colframe=orange!70!black,fonttitle=\bfseries,before skip=10pt,after skip=10pt,left=8pt,right=8pt}
\tcolorboxenvironment{lemma}{enhanced,breakable,colback=yellow!5,colframe=yellow!50!black,fonttitle=\bfseries,before skip=10pt,after skip=10pt,left=8pt,right=8pt}
\tcolorboxenvironment{corollary}{enhanced,breakable,colback=green!5,colframe=green!50!black,fonttitle=\bfseries,before skip=10pt,after skip=10pt,left=8pt,right=8pt}
\tcolorboxenvironment{remark}{enhanced,breakable,colback=gray!5,colframe=gray!50!black,before skip=8pt,after skip=8pt,left=6pt,right=6pt}
\tcolorboxenvironment{example}{enhanced,breakable,colback=purple!5,colframe=purple!50!black,before skip=8pt,after skip=8pt,left=6pt,right=6pt}

% ---- Custom commands ----
\newcommand{\N}{\mathbb{N}}
\newcommand{\Z}{\mathbb{Z}}
\newcommand{\Q}{\mathbb{Q}}
\newcommand{\R}{\mathbb{R}}
\newcommand{\C}{\mathbb{C}}
\newcommand{\F}{\mathbb{F}}
\newcommand{\abs}[1]{\left|#1\right|}
\newcommand{\norm}[1]{\left\|#1\right\|}
\newcommand{\floor}[1]{\left\lfloor #1 \right\rfloor}
\newcommand{\ceil}[1]{\left\lceil #1 \right\rceil}
\newcommand{\legendre}[2]{\left(\frac{#1}{#2}\right)}
\newcommand{\jacobi}[2]{\left(\frac{#1}{#2}\right)}
\DeclareMathOperator{\lcm}{lcm}
\DeclareMathOperator{\ord}{ord}

% ============================================================
\begin{document}

% ---- Title page ----
\begin{titlepage}
\begin{tikzpicture}[remember picture, overlay]
  \fill[left color=blue!15!white, right color=blue!5!white]
    (current page.south west) rectangle (current page.north east);
  \fill[color=orange!70!yellow]
    ([yshift=-2cm]current page.north west) rectangle ([yshift=-2.3cm]current page.north east);
  \fill[color=blue!80!black, rounded corners=8pt]
    ([xshift=3cm, yshift=-4cm]current page.north west) rectangle
    ([xshift=-3cm, yshift=-10cm]current page.north east);
  \node[anchor=center, text=white, font=\Huge\bfseries]
    at ([yshift=-6cm]current page.north) {Number Theory};
  \node[anchor=center, text=orange!80!yellow, font=\Large]
    at ([yshift=-7.5cm]current page.north) {Lecture Notes};
  \node[anchor=center, text=white!80, font=\large]
    at ([yshift=-8.8cm]current page.north) {Master M1 --- 2025--2026};
  \node[anchor=center, text=blue!80!black, font=\Large\itshape]
    at ([yshift=-12cm]current page.north) {Ya\'e Ulrich Gaba};
  \draw[orange!70!yellow, line width=2pt]
    ([xshift=5cm, yshift=-13.5cm]current page.north west) --
    ([xshift=-5cm, yshift=-13.5cm]current page.north east);
  \node[anchor=center, text width=12cm, align=center, text=blue!60!black, font=\itshape]
    at ([yshift=-16cm]current page.north) {``Mathematics is the queen of the sciences and number theory is the queen of mathematics.''\\[4pt]--- Carl Friedrich Gauss};
  \node[anchor=south, text=gray, font=\small]
    at ([yshift=1.5cm]current page.south) {\today};
\end{tikzpicture}
\end{titlepage}

\frontmatter

% ---- Preface ----
\chapter*{Preface}
\addcontentsline{toc}{chapter}{Preface}

Number theory, often called the \emph{queen of mathematics}, is one of the oldest
and most beautiful branches of the discipline. Its central objects---the integers---are
deceptively simple, yet the questions they raise have challenged mathematicians
for millennia and continue to drive deep research today.

This text aims to provide a rigorous, self-contained introduction to classical
number theory, beginning with the most elementary notions of divisibility and
building steadily toward the distribution of prime numbers, modular arithmetic,
quadratic reciprocity, and selected modern applications including cryptography.

\medskip
\noindent\textbf{Prerequisites.}
The reader is expected to be familiar with basic logic, proof techniques
(induction, contradiction, contrapositive), and elementary properties of
the integers as covered in a first course in discrete mathematics or algebra.
No prior exposure to number theory is assumed.

\medskip
\noindent\textbf{How to use this book.}
Definitions, theorems, and propositions are stated precisely and proved in full
unless explicitly noted otherwise. Numerous examples illustrate the theory, and
each chapter concludes with a graded set of exercises:
\begin{itemize}[nosep]
  \item[$\star$] --- routine applications of definitions and theorems;
  \item[$\star\star$] --- problems requiring some ingenuity or the combination
        of several results;
  \item[$\star\star\star$] --- challenging problems, often connected to
        competition mathematics or research.
\end{itemize}
The reader is strongly encouraged to attempt exercises before continuing to the
next chapter: number theory is learned by \emph{doing}.

\medskip
\noindent\textbf{Acknowledgements.}
The exposition owes a debt to the classic texts of Hardy and Wright,
Niven--Zuckerman--Montgomery, Ireland and Rosen, and Apostol, as well as to
generations of students whose questions have shaped the presentation.

% ---- Notation ----
\chapter*{Notation and Conventions}
\addcontentsline{toc}{chapter}{Notation and Conventions}

\index{notation}

Throughout this text we adopt the following conventions.

\begin{center}
\renewcommand{\arraystretch}{1.35}
\begin{tabular}{>{\raggedright}p{3.8cm} p{9cm}}
\toprule
\textbf{Symbol} & \textbf{Meaning} \\
\midrule
$\N$ & The set of natural numbers $\{0,1,2,\ldots\}$ \\
$\N^{*}$ & The set of positive integers $\{1,2,3,\ldots\}$ \\
$\Z$ & The ring of integers \\
$\Q,\;\R,\;\C$ & The fields of rational, real, and complex numbers \\
$\F_p$ & The finite field with $p$ elements ($p$ prime) \\
$a \mid b$ & $a$ divides $b$ \\
$a \nmid b$ & $a$ does not divide $b$ \\
$\gcd(a,b)$ & Greatest common divisor of $a$ and $b$ \\
$\lcm(a,b)$ & Least common multiple of $a$ and $b$ \\
$(a,b)$ & Shorthand for $\gcd(a,b)$ when context is clear \\
$a \equiv b \pmod{n}$ & $n$ divides $a-b$ \\
$\floor{x}$ & Floor function: greatest integer $\leq x$ \\
$\ceil{x}$ & Ceiling function: least integer $\geq x$ \\
$\pi(x)$ & Number of primes $\leq x$ \\
$\legendre{a}{p}$ & Legendre symbol \\
$\varphi(n)$ & Euler's totient function \\
$\ord_n(a)$ & Multiplicative order of $a$ modulo $n$ \\
$\log$ & Natural logarithm (unless otherwise stated) \\
\bottomrule
\end{tabular}
\end{center}

\medskip
\noindent
All rings are commutative with unity unless stated otherwise.
When we write $p$ without further qualification, it denotes a prime number.
The notation $\sum_{p \le x}$ means the sum over all primes $p$ not exceeding $x$.

% ---- Table of contents ----
\tableofcontents

\mainmatter

% ############################################################
% CHAPTER 1
% ############################################################
\chapter{Divisibility and the Euclidean Algorithm}
\label{ch:divisibility}
\index{divisibility}

\begin{quote}
\itshape
``The laws of nature are but the mathematical thoughts of God.''\\
--- attributed to Euclid
\end{quote}

\section{Historical context: Euclid's \textit{Elements}}
\label{sec:euclid-history}
\index{Euclid}\index{Elements@\textit{Elements}}

The systematic study of divisibility properties of integers traces back at least
to Books~VII--IX of Euclid's \textit{Elements} (c.~300 BCE). In these books,
Euclid established the foundation upon which all of number theory rests: the
notion of divisibility, the Euclidean algorithm for computing greatest common
divisors, and the celebrated proof that there are infinitely many prime numbers.

Euclid's approach was geometric in flavour---integers were represented as line
segments, and ``$a$ divides $b$'' meant that $a$ measures $b$ exactly. Despite
the geometric language, the logical structure is purely arithmetic and translates
directly into modern algebraic notation. The \emph{Euclidean algorithm},
Proposition~2 of Book~VII, is arguably the oldest non-trivial algorithm still
in everyday use: it underlies the computation of modular inverses in
cryptographic protocols executed billions of times per day.

In this chapter we develop divisibility theory from first principles, prove the
division algorithm and the correctness of the Euclidean algorithm, establish
B\'{e}zout's identity, and introduce the least common multiple and the notion
of coprimality.

% ============================================================
\section{Divisibility: definition and elementary properties}
\label{sec:divisibility-def}
\index{divisibility!definition}

\begin{definition}[Divisibility]
\label{def:divisibility}
Let $a,b \in \Z$. We say that $a$ \emph{divides}\index{divides} $b$, written
$a \mid b$, if there exists $k \in \Z$ such that $b = ak$. When $a \mid b$ we
also say that $a$ is a \emph{divisor}\index{divisor} (or \emph{factor}) of $b$,
and that $b$ is a \emph{multiple}\index{multiple} of $a$. If $a$ does not
divide $b$ we write $a \nmid b$.
\end{definition}

\begin{example}
\label{ex:div-basic}
We have $3 \mid 12$ since $12 = 3 \cdot 4$, and $7 \mid (-21)$ since
$-21 = 7 \cdot (-3)$. On the other hand, $5 \nmid 13$ because there is no
integer $k$ with $13 = 5k$.
\end{example}

\begin{proposition}[Basic properties of divisibility]
\label{prop:div-properties}
\index{divisibility!properties}
Let $a,b,c \in \Z$.
\begin{enumerate}[label=(\roman*)]
  \item \textbf{Reflexivity:} $a \mid a$ for all $a \in \Z$.
  \item \textbf{Transitivity:} If $a \mid b$ and $b \mid c$, then $a \mid c$.
  \item \textbf{Linearity:} If $a \mid b$ and $a \mid c$, then
        $a \mid (bx + cy)$ for all $x,y \in \Z$.
  \item \textbf{Comparison:} If $a \mid b$ and $b \neq 0$, then $\abs{a} \le \abs{b}$.
  \item \textbf{Antisymmetry on $\N^{*}$:} If $a,b \in \N^{*}$ with $a \mid b$
        and $b \mid a$, then $a = b$.
  \item $a \mid 0$ for every $a \in \Z$, and $1 \mid a$ for every $a \in \Z$.
  \item If $a \mid 1$, then $a \in \{-1, 1\}$.
\end{enumerate}
\end{proposition}

\begin{proof}
\begin{enumerate}[label=(\roman*)]
  \item $a = a \cdot 1$.
  \item Write $b = ak_1$ and $c = bk_2$. Then $c = a(k_1 k_2)$.
  \item Write $b = ak_1$ and $c = ak_2$. Then
        $bx + cy = a(k_1 x + k_2 y)$.
  \item Write $b = ak$ with $k \in \Z \setminus \{0\}$
        (since $b \ne 0$). Then $\abs{b} = \abs{a}\abs{k} \ge \abs{a}$.
  \item By~(iv), $a \le b$ and $b \le a$, so $a = b$.
  \item $0 = a \cdot 0$ and $a = 1 \cdot a$.
  \item Write $1 = ak$. Then $\abs{a}\abs{k} = 1$ with $\abs{a},\abs{k} \in \N^{*}$,
        so $\abs{a} = \abs{k} = 1$. \qedhere
\end{enumerate}
\end{proof}

\begin{remark}
\label{rem:divides-partial-order}
Properties~(i), (ii), and~(v) show that divisibility is a partial order on
$\N^{*}$. We shall visualise this order with a \emph{Hasse diagram} at the end
of the chapter (see Figure~\ref{fig:divisibility-lattice}).
\end{remark}

% ============================================================
\section{The division algorithm}
\label{sec:division-algorithm}
\index{division algorithm}

The following result is the cornerstone of integer arithmetic. Despite the
traditional name, it is a \emph{theorem of existence and uniqueness}, not a
constructive procedure (though its proof is readily made constructive).

\begin{theorem}[Division algorithm]
\label{thm:division-algorithm}
Let $a \in \Z$ and $b \in \Z$ with $b > 0$. There exist unique integers
$q$\index{quotient} (the \emph{quotient}) and $r$\index{remainder}
(the \emph{remainder}) such that
\[
  a = bq + r, \qquad 0 \le r < b.
\]
\end{theorem}

\begin{proof}
\textbf{Existence.}\index{division algorithm!existence}
Consider the set
\[
  S = \{ a - bk : k \in \Z \} \cap \N = \{ a - bk \ge 0 : k \in \Z \}.
\]
We claim $S \ne \varnothing$. Indeed, if $a \ge 0$ take $k = 0$; if $a < 0$
take $k = a$ (then $a - ba = a(1-b) \ge 0$ since $a < 0$ and $1 - b \le 0$).
By the well-ordering principle\index{well-ordering principle}, $S$ has a
least element $r = a - bq$ for some $q \in \Z$. By construction $r \ge 0$.

Suppose for contradiction that $r \ge b$. Then
\[
  a - b(q+1) = r - b \ge 0,
\]
so $r - b \in S$ and $r - b < r$, contradicting the minimality of $r$.
Therefore $0 \le r < b$.

\medskip
\textbf{Uniqueness.}\index{division algorithm!uniqueness}
Suppose $a = bq_1 + r_1 = bq_2 + r_2$ with $0 \le r_1, r_2 < b$. Then
\[
  b(q_1 - q_2) = r_2 - r_1.
\]
Since $\abs{r_2 - r_1} < b$ and $b \mid (r_2 - r_1)$, we must have
$r_2 - r_1 = 0$, i.e.\ $r_1 = r_2$. It follows immediately that $q_1 = q_2$.
\end{proof}

\begin{example}
\label{ex:division-algorithm}
Divide $a = -17$ by $b = 5$: we seek $q,r$ with $-17 = 5q + r$ and
$0 \le r < 5$. We find $q = -4$ and $r = 3$ since $-17 = 5(-4) + 3$.
\end{example}

\begin{remark}
\label{rem:division-alg-notation}
We write $r = a \bmod b$\index{modular reduction} or equivalently
$r = a - b\floor{a/b}$. This extends naturally to $b < 0$; some authors
require $0 \le r < \abs{b}$ for the general case.
\end{remark}

% ============================================================
\section{Greatest common divisor and the Euclidean algorithm}
\label{sec:gcd}
\index{greatest common divisor}\index{gcd}

\begin{definition}[Greatest common divisor]
\label{def:gcd}
Let $a,b \in \Z$, not both zero. The \emph{greatest common divisor} of $a$
and $b$, denoted $\gcd(a,b)$, is the largest positive integer $d$ such that
$d \mid a$ and $d \mid b$.
\end{definition}

\begin{remark}
\label{rem:gcd-convention}
By convention, $\gcd(0,0) = 0$. For all $a$, $\gcd(a,0) = \abs{a}$.
\end{remark}

\begin{proposition}[Characterisation of the gcd]
\label{prop:gcd-char}
Let $d = \gcd(a,b)$. Then:
\begin{enumerate}[label=(\roman*)]
  \item $d \mid a$ and $d \mid b$.
  \item If $c \mid a$ and $c \mid b$, then $c \mid d$.
\end{enumerate}
Conversely, any positive integer satisfying (i) and (ii) equals $\gcd(a,b)$.
\end{proposition}

\begin{proof}
Property~(i) is immediate from the definition. For~(ii), suppose $c \mid a$
and $c \mid b$. We shall prove $c \mid d$ after establishing B\'{e}zout's
identity (Theorem~\ref{thm:bezout}), which gives $d = ax + by$ for some
$x, y \in \Z$. Since $c \mid a$ and $c \mid b$, linearity of divisibility
(Proposition~\ref{prop:div-properties}(iii)) yields $c \mid d$.
\end{proof}

\begin{lemma}[Key reduction for the Euclidean algorithm]
\label{lem:gcd-reduction}
\index{Euclidean algorithm!key reduction}
For any $a \in \Z$ and $b \in \N^{*}$,
\[
  \gcd(a, b) = \gcd(b, a \bmod b).
\]
\end{lemma}

\begin{proof}
Write $a = bq + r$ with $0 \le r < b$ (Theorem~\ref{thm:division-algorithm}).
If $d \mid a$ and $d \mid b$, then $d \mid (a - bq) = r$. Conversely, if
$d \mid b$ and $d \mid r$, then $d \mid (bq + r) = a$. Hence the set of
common divisors of $(a,b)$ equals the set of common divisors of $(b,r)$,
so their greatest elements coincide.
\end{proof}

\begin{theorem}[Euclidean algorithm]
\label{thm:euclidean-algorithm}
\index{Euclidean algorithm}
Let $a, b \in \N$ with $b > 0$. Define the sequence
\begin{align*}
  r_0 &= a, \quad r_1 = b, \\
  r_{i+1} &= r_{i-1} \bmod r_i, \qquad i = 1, 2, \ldots
\end{align*}
There exists $n \ge 1$ such that $r_{n+1} = 0$, and then
$\gcd(a,b) = r_n$.
\end{theorem}

\begin{proof}
\textbf{Termination.}
The sequence $(r_i)$ satisfies $b = r_1 > r_2 > \cdots \ge 0$.
Since a strictly decreasing sequence of non-negative integers must be finite,
there exists a smallest $n$ with $r_{n+1} = 0$.

\medskip
\textbf{Correctness.}
By repeated application of Lemma~\ref{lem:gcd-reduction}:
\[
  \gcd(a,b) = \gcd(r_0, r_1) = \gcd(r_1, r_2) = \cdots
  = \gcd(r_n, r_{n+1}) = \gcd(r_n, 0) = r_n. \qedhere
\]
\end{proof}

\begin{example}
\label{ex:euclidean-algorithm}
Compute $\gcd(252, 198)$:
\begin{align*}
  252 &= 198 \cdot 1 + 54, \\
  198 &= 54 \cdot 3 + 36, \\
  54 &= 36 \cdot 1 + 18, \\
  36 &= 18 \cdot 2 + 0.
\end{align*}
Hence $\gcd(252, 198) = 18$.
\end{example}

\begin{remark}[Complexity of the Euclidean algorithm]
\label{rem:euclidean-complexity}
\index{Euclidean algorithm!complexity}\index{Fibonacci numbers}
Lam\'{e} proved in 1844 that the number of division steps is at most
$5$ times the number of digits of the smaller input. More precisely,
the worst case is achieved by consecutive Fibonacci numbers: computing
$\gcd(F_{n+1}, F_n)$ requires exactly $n-1$ steps. This gives a time
complexity of $O((\log \min(a,b))^2)$ when using standard integer arithmetic.
\end{remark}

% ============================================================
\section{B\'{e}zout's identity and the extended Euclidean algorithm}
\label{sec:bezout}
\index{Bezout's identity@B\'{e}zout's identity}

\begin{theorem}[B\'{e}zout's identity]
\label{thm:bezout}
Let $a,b \in \Z$, not both zero, and let $d = \gcd(a,b)$. There exist
$x, y \in \Z$ such that
\[
  ax + by = d.
\]
Moreover, $d$ is the smallest positive element of the set
$\{ax + by : x,y \in \Z\}$.
\end{theorem}

\begin{proof}
\index{Bezout's identity@B\'{e}zout's identity!proof}
Consider the set
\[
  S = \{ax + by : x,y \in \Z\} \cap \N^{*}.
\]
The set $S$ is non-empty (it contains $\abs{a}$ or $\abs{b}$, at least one of
which is positive). By the well-ordering principle, $S$ has a least element
$d_0 = ax_0 + by_0$.

We claim $d_0 \mid a$. Apply the division algorithm: $a = d_0 q + r$ with
$0 \le r < d_0$. Then
\[
  r = a - d_0 q = a - (ax_0 + by_0)q = a(1 - x_0 q) + b(-y_0 q).
\]
So $r \in \{ax + by : x,y \in \Z\}$ and $0 \le r < d_0$. Since $d_0$ is the
smallest \emph{positive} element of this set, we must have $r = 0$. Therefore
$d_0 \mid a$. By the same argument, $d_0 \mid b$.

Since $d_0$ is a common divisor of $a$ and $b$, we have $d_0 \le d$.
Conversely, $d \mid a$ and $d \mid b$ imply $d \mid (ax_0 + by_0) = d_0$,
so $d \le d_0$. Thus $d = d_0$.
\end{proof}

\begin{corollary}
\label{cor:gcd-linear-comb}
An integer $c$ is a $\Z$-linear combination of $a$ and $b$ if and only if
$\gcd(a,b) \mid c$.
\end{corollary}

\begin{proof}
If $c = ax + by$ and $d = \gcd(a,b)$, then $d \mid c$ by linearity.
Conversely, if $d \mid c$, write $c = dk$ and $d = ax_0 + by_0$; then
$c = a(kx_0) + b(ky_0)$.
\end{proof}

\begin{definition}[Extended Euclidean algorithm]
\label{def:extended-euclidean}
\index{Euclidean algorithm!extended}
The \emph{extended Euclidean algorithm} computes, alongside $\gcd(a,b)$,
integers $x$ and $y$ satisfying $ax + by = \gcd(a,b)$. It does so by
tracking B\'{e}zout coefficients through the recurrence:
\[
  \begin{pmatrix} x_i \\ y_i \end{pmatrix}
  = \begin{pmatrix} x_{i-2} \\ y_{i-2} \end{pmatrix}
  - q_i \begin{pmatrix} x_{i-1} \\ y_{i-1} \end{pmatrix},
\]
starting from $(x_0, y_0) = (1, 0)$ and $(x_1, y_1) = (0, 1)$, where
$q_i = \floor{r_{i-1}/r_i}$ at step $i$.
\end{definition}

\begin{example}
\label{ex:extended-euclidean}
Find $x,y$ with $252x + 198y = \gcd(252,198) = 18$.

\medskip
\begin{center}
\begin{tabular}{c c c c c}
\toprule
$i$ & $r_i$ & $q_i$ & $x_i$ & $y_i$ \\
\midrule
$0$ & $252$ & --- & $1$ & $0$ \\
$1$ & $198$ & $1$ & $0$ & $1$ \\
$2$ & $54$ & $3$ & $1$ & $-1$ \\
$3$ & $36$ & $1$ & $-3$ & $4$ \\
$4$ & $18$ & $2$ & $4$ & $-5$ \\
$5$ & $0$ & --- & --- & --- \\
\bottomrule
\end{tabular}
\end{center}

\medskip
We read off $x = 4$, $y = -5$, and verify: $252 \cdot 4 + 198 \cdot(-5) = 1008 - 990 = 18$. \checkmark
\end{example}

% ============================================================
\section{Coprime integers, Gauss's lemma, and the LCM}
\label{sec:coprime-lcm}

\begin{definition}[Coprime integers]
\label{def:coprime}
\index{coprime}
Integers $a$ and $b$ are \emph{coprime} (or \emph{relatively prime}) if
$\gcd(a,b) = 1$.
\end{definition}

\begin{proposition}
\label{prop:coprime-bezout}
$a$ and $b$ are coprime if and only if there exist $x, y \in \Z$ with
$ax + by = 1$.
\end{proposition}

\begin{proof}
If $\gcd(a,b) = 1$, B\'{e}zout gives the desired $x,y$.
Conversely, if $ax + by = 1$, any common divisor $d$ of $a$ and $b$ divides~$1$,
so $d = \pm 1$, hence $\gcd(a,b) = 1$.
\end{proof}

\begin{theorem}[Gauss's lemma]
\label{thm:gauss-lemma}
\index{Gauss's lemma}
If $a \mid bc$ and $\gcd(a,b) = 1$, then $a \mid c$.
\end{theorem}

\begin{proof}
By B\'{e}zout, write $ax + by = 1$. Multiply by $c$:
\[
  acx + bcy = c.
\]
Since $a \mid acx$ and $a \mid bc$ (hence $a \mid bcy$), we obtain $a \mid c$.
\end{proof}

\begin{corollary}[Euclid's lemma]
\label{cor:euclid-lemma}
\index{Euclid's lemma}
If $p$ is prime and $p \mid ab$, then $p \mid a$ or $p \mid b$.
\end{corollary}

\begin{proof}
If $p \nmid a$ then $\gcd(p,a) = 1$ (the only positive divisors of $p$ are $1$
and $p$). By Gauss's lemma, $p \mid b$.
\end{proof}

\begin{corollary}
\label{cor:euclid-lemma-general}
If $p$ is prime and $p \mid a_1 a_2 \cdots a_n$, then $p \mid a_i$ for some $i$.
\end{corollary}

\begin{proof}
Induction on $n$, using Corollary~\ref{cor:euclid-lemma} as the base case.
\end{proof}

\begin{definition}[Least common multiple]
\label{def:lcm}
\index{least common multiple}\index{lcm}
The \emph{least common multiple} of $a,b \in \Z \setminus \{0\}$, denoted
$\lcm(a,b)$, is the smallest positive integer $m$ such that $a \mid m$ and
$b \mid m$.
\end{definition}

\begin{proposition}
\label{prop:lcm-properties}
For $a, b \in \N^{*}$:
\begin{enumerate}[label=(\roman*)]
  \item $\lcm(a,b) = \dfrac{ab}{\gcd(a,b)}$.
  \item If $a \mid n$ and $b \mid n$, then $\lcm(a,b) \mid n$.
\end{enumerate}
\end{proposition}

\begin{proof}
\begin{enumerate}[label=(\roman*)]
  \item Let $d = \gcd(a,b)$ and $m = ab/d$. Write $a = da'$, $b = db'$ with
        $\gcd(a',b') = 1$. Then $m = da'b'$, so $a \mid m$ (since $m = b \cdot a'$)
        and $b \mid m$ (since $m = a \cdot b'$). Hence $m$ is a common multiple.

        Now let $n$ be any positive common multiple. Write $n = ak_1 = bk_2$. Then
        $da'k_1 = db'k_2$, so $a'k_1 = b'k_2$. Since $\gcd(a',b') = 1$, Gauss's
        lemma gives $b' \mid k_1$, say $k_1 = b'\ell$. Then
        $n = a k_1 = da'b'\ell = m\ell$, so $m \mid n$ and thus $m \le n$.
  \item This was shown in the proof of~(i). \qedhere
\end{enumerate}
\end{proof}

% ============================================================
\section{Divisibility lattice}
\label{sec:divisibility-lattice}
\index{divisibility lattice}\index{Hasse diagram}

The divisibility relation $\mid$ defines a partial order on $\N^{*}$.
For a fixed positive integer $n$, the set of divisors of $n$ ordered by
divisibility forms a lattice in which the meet is $\gcd$ and the join is $\lcm$.

\begin{figure}[ht]
\centering
\begin{tikzpicture}[
  every node/.style={draw, circle, minimum size=7mm, inner sep=0pt,
                     fill=blue!10, font=\small},
  edge/.style={thick, blue!60!black},
  scale=1.1
]
  % Divisibility lattice of 36
  \node (1)  at (0,0)   {$1$};
  \node (2)  at (-2,1.2) {$2$};
  \node (3)  at (0,1.2)  {$3$};
  \node (4)  at (-2,2.4) {$4$};
  \node (6)  at (0,2.4)  {$6$};
  \node (9)  at (2,2.4)  {$9$};
  \node (12) at (-1,3.6) {$12$};
  \node (18) at (1,3.6)  {$18$};
  \node (36) at (0,4.8)  {$36$};

  \draw[edge] (1)  -- (2);
  \draw[edge] (1)  -- (3);
  \draw[edge] (2)  -- (4);
  \draw[edge] (2)  -- (6);
  \draw[edge] (3)  -- (6);
  \draw[edge] (3)  -- (9);
  \draw[edge] (4)  -- (12);
  \draw[edge] (6)  -- (12);
  \draw[edge] (6)  -- (18);
  \draw[edge] (9)  -- (18);
  \draw[edge] (12) -- (36);
  \draw[edge] (18) -- (36);
\end{tikzpicture}
\caption{Hasse diagram of the divisibility lattice of $36$.
  For example, $\gcd(12,18) = 6$ (meet) and $\lcm(12,18) = 36$ (join).}
\label{fig:divisibility-lattice}
\end{figure}

% ============================================================
\section{Connections to cryptography}
\label{sec:crypto-gcd}
\index{cryptography}\index{RSA}

The Euclidean algorithm and B\'{e}zout's identity are not merely theoretical
curiosities; they are essential tools in modern cryptography.

In the \textbf{RSA cryptosystem}\index{RSA}, one selects two large primes
$p$ and $q$, computes $n = pq$, and chooses an encryption exponent $e$ with
$\gcd(e, \varphi(n)) = 1$, where $\varphi(n) = (p-1)(q-1)$ is Euler's totient.
The decryption exponent $d$ is then the modular inverse of $e$ modulo
$\varphi(n)$:
\[
  ed \equiv 1 \pmod{\varphi(n)}.
\]
This inverse is computed via the \emph{extended Euclidean algorithm}.
The security of the system rests on the difficulty of factoring $n$, but
its \emph{implementation} depends fundamentally on the efficient computation
of gcd and modular inverses.

\begin{remark}
\label{rem:crypto-efficiency}
In practice, the integers involved have thousands of decimal digits. The
$O((\log n)^2)$ complexity of the Euclidean algorithm (with fast multiplication
improvements) makes this entirely feasible, while the factoring problem for
the same size integers remains computationally intractable.
\end{remark}

% ============================================================
\section{Exercises}
\label{sec:ch1-exercises}

\begin{exercise}[$\star$]
\label{exer:ch1-basic-div}
Prove that if $a \mid b$ and $b \mid c$, then $a \mid (b + c)$.
\end{exercise}

\begin{exercise}[$\star$]
\label{exer:ch1-div-square}
Show that if $n^2 \mid n$ for some $n \in \Z$, then $n \in \{-1, 0, 1\}$.
\end{exercise}

\begin{exercise}[$\star$]
\label{exer:ch1-euclidean-compute}
Use the Euclidean algorithm to compute $\gcd(1001, 385)$. Then find integers
$x, y$ with $1001x + 385y = \gcd(1001, 385)$.
\end{exercise}

\begin{exercise}[$\star$]
\label{exer:ch1-consecutive}
Prove that any two consecutive integers are coprime.
\end{exercise}

\begin{exercise}[$\star$]
\label{exer:ch1-gcd-lcm}
Let $a, b \in \N^{*}$. Prove that $\gcd(a,b) \cdot \lcm(a,b) = ab$.
\end{exercise}

\begin{exercise}[$\star\star$]
\label{exer:ch1-linear-diophantine}
Determine all integer solutions $(x,y)$ of $12x + 8y = 28$.
\end{exercise}

\begin{exercise}[$\star\star$]
\label{exer:ch1-coprime-product}
Let $a, b, c \in \Z$ with $\gcd(a,b) = 1$. Prove that $\gcd(a, bc) = \gcd(a,c)$.
\end{exercise}

\begin{exercise}[$\star\star$]
\label{exer:ch1-fibonacci-gcd}
\index{Fibonacci numbers!gcd identity}
Let $(F_n)$ be the Fibonacci sequence with $F_1 = F_2 = 1$. Prove that
$\gcd(F_m, F_n) = F_{\gcd(m,n)}$ for all $m, n \ge 1$.
\end{exercise}

\begin{exercise}[$\star\star$]
\label{exer:ch1-bezout-unique}
Show that the B\'{e}zout coefficients $x, y$ in $ax + by = \gcd(a,b)$ are not
unique. Describe all solutions $(x,y)$.
\end{exercise}

\begin{exercise}[$\star\star\star$]
\label{exer:ch1-lame}
\index{Lam\'e's theorem}
Prove \textbf{Lam\'{e}'s theorem}: the Euclidean algorithm applied to
$\gcd(a,b)$ with $a > b > 0$ requires at most $\floor{(\log_\phi b)} - 1$
divisions, where $\phi = (1+\sqrt{5})/2$ is the golden ratio.
\textit{Hint:} show that if the algorithm takes $n$ steps, then $b \ge F_{n+1}$,
where $F_k$ is the $k$th Fibonacci number.
\end{exercise}

\begin{exercise}[$\star\star\star$]
\label{exer:ch1-chicken-mcnugget}
\index{Chicken McNugget theorem}\index{Frobenius number}
Let $a, b \in \N^{*}$ with $\gcd(a,b) = 1$. Prove that the largest integer
that \emph{cannot} be expressed as $ax + by$ with $x, y \in \N$ is $ab - a - b$.
This is the \textbf{Frobenius number} $g(a,b)$.
\end{exercise}

% ============================================================
\section*{Chapter summary}
\addcontentsline{toc}{section}{Chapter summary}

\begin{itemize}[nosep]
  \item \textbf{Divisibility} ($a \mid b$) is a partial order on $\N^{*}$;
        it is reflexive, transitive, antisymmetric, and linear in the sense
        that divisibility is preserved under $\Z$-linear combinations.
  \item The \textbf{division algorithm} guarantees unique quotient and remainder
        for any division by a positive integer.
  \item The \textbf{Euclidean algorithm} computes $\gcd(a,b)$ in $O((\log b)^2)$
        time; the \textbf{extended} version also yields B\'{e}zout coefficients.
  \item \textbf{B\'{e}zout's identity}: $\gcd(a,b)$ is the smallest positive
        $\Z$-linear combination of $a$ and $b$.
  \item \textbf{Gauss's lemma}: if $a \mid bc$ and $\gcd(a,b) = 1$ then $a \mid c$.
        As a consequence, \textbf{Euclid's lemma}: if $p$ is prime and $p \mid ab$,
        then $p \mid a$ or $p \mid b$.
  \item The \textbf{lcm} satisfies $\gcd(a,b) \cdot \lcm(a,b) = ab$.
  \item These tools underpin \textbf{RSA cryptography}: modular inverses
        are computed via the extended Euclidean algorithm.
\end{itemize}


% ############################################################
% CHAPTER 2
% ############################################################
\chapter{Prime Numbers: Infinity, Arithmetic, and Distribution}
\label{ch:primes}
\index{prime numbers}

\begin{quote}
\itshape
``Mathematicians have tried in vain to this day to discover some order in the
sequence of prime numbers, and we have reason to believe that it is a mystery
into which the human mind will never penetrate.''\\
--- Leonhard Euler
\end{quote}

\section{Primes and the fundamental dichotomy}
\label{sec:primes-def}
\index{prime number!definition}

\begin{definition}[Prime and composite numbers]
\label{def:prime}
An integer $p > 1$ is \emph{prime}\index{prime} if its only positive divisors
are $1$ and $p$. An integer $n > 1$ that is not prime is called
\emph{composite}\index{composite number}.
\end{definition}

\begin{remark}
\label{rem:unit}
The integer $1$ is neither prime nor composite; it is the \emph{unit} of
$(\Z, \cdot)$. This convention is essential for the uniqueness statement of
the Fundamental Theorem of Arithmetic.
\end{remark}

\begin{example}
\label{ex:small-primes}
The first few primes are $2, 3, 5, 7, 11, 13, 17, 19, 23, 29, 31, 37, 41,
43, 47, \ldots$ The integer $2$ is the only even prime.
\end{example}

\begin{lemma}
\label{lem:smallest-divisor-prime}
Every integer $n > 1$ has a prime divisor. Moreover, the smallest divisor of
$n$ greater than $1$ is prime.
\end{lemma}

\begin{proof}
Let $d$ be the smallest element of $\{k \in \N : k > 1,\; k \mid n\}$; this
set is non-empty since $n$ itself belongs to it. Suppose $d$ is composite,
i.e.\ $d = ab$ with $1 < a, b < d$. Then $a \mid d$ and $d \mid n$, so
$a \mid n$, contradicting the minimality of $d$. Hence $d$ is prime.
\end{proof}

\begin{corollary}
\label{cor:composite-small-factor}
\index{trial division}
If $n > 1$ is composite, then $n$ has a prime factor $p \le \sqrt{n}$.
\end{corollary}

\begin{proof}
If $n = ab$ with $1 < a \le b < n$, then $a^2 \le ab = n$, so $a \le \sqrt{n}$.
By Lemma~\ref{lem:smallest-divisor-prime}, $a$ has a prime divisor $p \le a \le \sqrt{n}$.
\end{proof}

% ============================================================
\section{Infinitude of primes}
\label{sec:infinitude}
\index{prime numbers!infinitude}

The fact that there are infinitely many primes is one of the most celebrated
results in all of mathematics. We present three proofs, each illuminating a
different aspect of the structure.

\subsection{Euclid's proof}
\label{subsec:euclid-proof}
\index{Euclid's theorem}

\begin{theorem}[Euclid, c.~300 BCE]
\label{thm:euclid-infinitude}
There are infinitely many prime numbers.
\end{theorem}

\begin{proof}
Suppose for contradiction that there are only finitely many primes,
$p_1, p_2, \ldots, p_n$. Consider the integer
\[
  N = p_1 p_2 \cdots p_n + 1.
\]
Since $N > 1$, by Lemma~\ref{lem:smallest-divisor-prime}, $N$ has a prime
divisor $p$. Now $p = p_i$ for some $i$, so $p_i \mid N$ and
$p_i \mid p_1 \cdots p_n$. Hence $p_i \mid (N - p_1 \cdots p_n) = 1$,
a contradiction. Therefore the set of primes is infinite.
\end{proof}

\begin{remark}
\label{rem:euclid-constructive}
Euclid's proof is often misquoted as ``$p_1 \cdots p_n + 1$ is always prime.''
This is false in general: $2 \cdot 3 \cdot 5 \cdot 7 \cdot 11 \cdot 13 + 1 = 30031 = 59 \times 509$.
What the proof shows is that $N$ has a prime factor \emph{not in the list},
whether or not $N$ itself is prime.
\end{remark}

\subsection{Euler's analytic proof}
\label{subsec:euler-proof}
\index{Euler product}

\begin{theorem}[Euler, 1737]
\label{thm:euler-infinitude}
The sum $\displaystyle\sum_{p} \frac{1}{p}$ over all primes $p$ diverges.
In particular, there are infinitely many primes.
\end{theorem}

\begin{proof}[Proof sketch]
For $s > 1$ one has the \emph{Euler product}\index{Euler product}:
\[
  \sum_{n=1}^{\infty} \frac{1}{n^s}
  = \prod_{p \text{ prime}} \frac{1}{1 - p^{-s}}.
\]
The left side diverges as $s \to 1^{+}$ (harmonic series). If there were only
finitely many primes, the right side would be a finite product of finite values,
a contradiction. A more careful analysis shows that
\[
  \sum_{p \le x} \frac{1}{p} = \log\log x + M + O\!\left(\frac{1}{\log x}\right),
\]
where $M \approx 0.2615$ is the \emph{Meissel--Mertens constant}\index{Meissel--Mertens constant},
confirming divergence.
\end{proof}

\subsection{A proof via Fermat numbers}
\label{subsec:goldbach-proof}
\index{Fermat numbers}

\begin{theorem}
\label{thm:fermat-coprime}
The Fermat numbers $F_n = 2^{2^n} + 1$ ($n = 0, 1, 2, \ldots$) are pairwise
coprime. Hence there are infinitely many primes.
\end{theorem}

\begin{proof}
We first show that for $m < n$, $F_m \mid F_n - 2$. By induction, one verifies
\[
  F_n - 2 = F_0 F_1 \cdots F_{n-1}, \qquad n \ge 1.
\]
\emph{Base case:} $F_1 - 2 = 3 - 2 = 1$ and $F_0 = 3$; but we actually compute
$F_1 - 2 = 2^2 + 1 - 2 = 3 = F_0$, which works since the product $F_0 = 3$.

\emph{Inductive step:} Suppose $F_n - 2 = F_0 \cdots F_{n-1}$. Then
$F_{n+1} = (F_n - 1)^2 + 1 = F_n^2 - 2F_n + 2$, so
\[
  F_{n+1} - 2 = F_n(F_n - 2) = F_n \cdot F_0 \cdots F_{n-1} = F_0 F_1 \cdots F_n.
\]

Now if $d \mid F_m$ and $d \mid F_n$ with $m < n$, then $d \mid F_n - 2$ and
$d \mid F_n$, so $d \mid 2$. Since each $F_k$ is odd, $d$ is odd, so $d = 1$.
Hence $\gcd(F_m, F_n) = 1$.

Each $F_n > 1$ has at least one prime divisor, and these prime divisors are
distinct for different $n$. Thus there are infinitely many primes.
\end{proof}

% ============================================================
\section{The Fundamental Theorem of Arithmetic}
\label{sec:fta}
\index{Fundamental Theorem of Arithmetic}

\begin{theorem}[Fundamental Theorem of Arithmetic]
\label{thm:fta}
Every integer $n > 1$ can be written as a product of primes:
\[
  n = p_1^{a_1} p_2^{a_2} \cdots p_k^{a_k},
  \qquad p_1 < p_2 < \cdots < p_k, \quad a_i \ge 1.
\]
Moreover, this representation is \emph{unique} (up to the order of factors).
\end{theorem}

\begin{proof}
\textbf{Existence}\index{Fundamental Theorem of Arithmetic!existence}
(by strong induction on $n$).

\emph{Base case:} $n = 2$ is prime, so it is a product of one prime.

\emph{Inductive step:} Assume every integer $m$ with $2 \le m < n$ is a product
of primes. If $n$ is prime, we are done. Otherwise, $n$ is composite:
$n = ab$ with $1 < a, b < n$. By the inductive hypothesis, both $a$ and $b$ are
products of primes, hence so is $n = ab$.

\medskip
\textbf{Uniqueness}\index{Fundamental Theorem of Arithmetic!uniqueness}
(by strong induction on $n$).

Suppose
\[
  n = p_1 p_2 \cdots p_r = q_1 q_2 \cdots q_s,
\]
where the $p_i$ and $q_j$ are primes (written with repetitions, not necessarily
in order). We prove $r = s$ and the $p_i$ are a permutation of the $q_j$.

Since $p_1 \mid q_1 q_2 \cdots q_s$, Corollary~\ref{cor:euclid-lemma-general}
gives $p_1 \mid q_j$ for some $j$. Since $q_j$ is prime and $p_1 > 1$, we have
$p_1 = q_j$. Cancel this factor:
\[
  p_2 \cdots p_r = q_1 \cdots q_{j-1} q_{j+1} \cdots q_s.
\]
If $r = 1$, both sides equal $1$, forcing $s = 1$ and we are done.
If $r > 1$, the common value $n/p_1 < n$, and by the inductive hypothesis
the two factorisations of $n/p_1$ agree up to reordering. Prepending $p_1 = q_j$
gives the result.
\end{proof}

\begin{notation}[Canonical factorisation]
\label{not:canonical}
\index{canonical factorisation}
We write the factorisation of $n$ in \emph{canonical form}:
\[
  n = \prod_{p \text{ prime}} p^{v_p(n)},
\]
where $v_p(n) \ge 0$ is the \emph{$p$-adic valuation}\index{p-adic valuation@$p$-adic valuation} of $n$ and all but
finitely many exponents are zero. Then:
\[
  \gcd(a,b) = \prod_p p^{\min(v_p(a),\, v_p(b))},
  \qquad
  \lcm(a,b) = \prod_p p^{\max(v_p(a),\, v_p(b))}.
\]
\end{notation}

\begin{example}
\label{ex:fta-computation}
$360 = 2^3 \cdot 3^2 \cdot 5$ and $150 = 2 \cdot 3 \cdot 5^2$. Therefore
$\gcd(360, 150) = 2 \cdot 3 \cdot 5 = 30$ and
$\lcm(360, 150) = 2^3 \cdot 3^2 \cdot 5^2 = 1800$.
\end{example}

% ============================================================
\section{The Sieve of Eratosthenes}
\label{sec:sieve}
\index{Sieve of Eratosthenes}

The \emph{Sieve of Eratosthenes} (c.~240 BCE)\index{Eratosthenes} is the oldest
known systematic method for listing all primes up to a given bound $N$. The
idea is simple: write down all integers from $2$ to $N$, then iteratively mark
off the multiples of each prime found.

\medskip
\noindent\textbf{Algorithm.}
\begin{enumerate}[nosep]
  \item List the integers $2, 3, 4, \ldots, N$.
  \item Let $p = 2$ (the smallest prime).
  \item Mark all multiples of $p$ from $p^2$ to $N$ as composite.
        (Multiples $2p, 3p, \ldots, (p-1)p$ have already been handled.)
  \item Find the next unmarked integer $> p$; call it $p'$. Set $p \leftarrow p'$.
  \item If $p^2 \le N$, go to step~3. Otherwise, stop: all unmarked numbers are prime.
\end{enumerate}

\begin{proposition}
\label{prop:sieve-correctness}
The Sieve of Eratosthenes correctly identifies all primes up to $N$. Its
time complexity is $O(N \log\log N)$.
\end{proposition}

\begin{proof}[Proof of correctness]
By Corollary~\ref{cor:composite-small-factor}, every composite $n \le N$ has a
prime factor $p \le \sqrt{N}$. The sieve marks $n$ as composite in the
iteration for $p$. Conversely, a prime $q \le N$ is never marked: it is not a
proper multiple of any smaller prime.
\end{proof}

\begin{figure}[ht]
\centering
\begin{tikzpicture}[
  cell/.style={minimum size=7mm, inner sep=0pt, font=\small\ttfamily},
  prime/.style={cell, fill=green!30, draw=green!60!black, rounded corners=2pt, font=\small\bfseries},
  composite/.style={cell, fill=red!15, draw=red!40, text=gray!60},
  one/.style={cell, fill=gray!15, draw=gray!40, text=gray!50}
]
  \foreach \i in {1,...,50} {
    \pgfmathtruncatemacro{\col}{mod(\i-1,10)}
    \pgfmathtruncatemacro{\row}{div(\i-1,10)}
    \pgfmathtruncatemacro{\ypos}{-\row}
    % Determine if prime
    \pgfmathtruncatemacro{\isprime}{%
      (\i>1) &&
      (\i==2 || mod(\i,2)!=0) &&
      (\i==3 || mod(\i,3)!=0) &&
      (\i<25 || mod(\i,5)!=0) &&
      (\i<49 || mod(\i,7)!=0) ? 1 : 0}
    \ifnum\i=1
      \node[one] at (\col, \ypos) {\i};
    \else
      \ifnum\isprime=1
        \node[prime] at (\col, \ypos) {\i};
      \else
        \node[composite] at (\col, \ypos) {\i};
      \fi
    \fi
  }
  \node[anchor=north, font=\footnotesize] at (4.5, -5.3)
    {Green = prime, Red = composite, Gray = 1 (unit)};
\end{tikzpicture}
\caption{Result of the Sieve of Eratosthenes for integers $1$ to $50$.}
\label{fig:sieve}
\end{figure}

% ============================================================
\section{The prime counting function and the Prime Number Theorem}
\label{sec:pnt}
\index{prime counting function}\index{pi(x)@$\pi(x)$}

\begin{definition}
\label{def:pi-x}
For $x \ge 0$, the \emph{prime counting function} is
\[
  \pi(x) = \#\{p \le x : p \text{ is prime}\}.
\]
\end{definition}

\begin{example}
\label{ex:pi-values}
$\pi(10) = 4$ (the primes $2, 3, 5, 7$), $\pi(100) = 25$, $\pi(1000) = 168$,
$\pi(10^6) = 78{,}498$.
\end{example}

The central question of analytic number theory is: \emph{how does $\pi(x)$
grow as $x \to \infty$?}

\subsection{Chebyshev's bounds}
\label{subsec:chebyshev}
\index{Chebyshev bounds}

In the 1850s, Chebyshev obtained the correct order of growth without the
precise asymptotic constant.

\begin{theorem}[Chebyshev, 1852]
\label{thm:chebyshev}
There exist positive constants $c_1, c_2$ such that for all sufficiently
large $x$:
\[
  c_1 \,\frac{x}{\log x} \le \pi(x) \le c_2 \,\frac{x}{\log x}.
\]
Chebyshev showed one can take $c_1 = \frac{\log 2}{2} \approx 0.347$ and
$c_2 = \frac{6 \log 2}{5} \approx 0.832$.
\end{theorem}

\begin{proof}[Proof idea]
The lower bound follows from an analysis of the central binomial coefficient
$\binom{2n}{n}$. Since $\binom{2n}{n} \le 4^n$ and every prime
$n < p \le 2n$ divides $\binom{2n}{n}$, one obtains $\pi(2n) - \pi(n) \le$
a controlled quantity. The upper bound similarly uses the prime factorisation
of $\binom{2n}{n}$ to bound $\sum_{p \le 2n} \log p$ from above. We refer
to Chapter~5 of Hardy and Wright for complete details.
\end{proof}

\subsection{The Prime Number Theorem}
\label{subsec:pnt}
\index{Prime Number Theorem}

\begin{theorem}[Prime Number Theorem --- Hadamard, de la Vall\'{e}e-Poussin, 1896]
\label{thm:pnt}
\[
  \pi(x) \sim \frac{x}{\log x}
  \qquad\text{as } x \to \infty.
\]
Equivalently, $\displaystyle\lim_{x \to \infty} \frac{\pi(x) \log x}{x} = 1$.
A better approximation is given by the \emph{logarithmic integral}\index{logarithmic integral}:
\[
  \pi(x) \sim \operatorname{Li}(x) = \int_2^{x} \frac{dt}{\log t}.
\]
\end{theorem}

\begin{remark}
\label{rem:pnt-proof}
The original proofs used complex analysis (properties of the Riemann zeta
function $\zeta(s)$ on the line $\operatorname{Re}(s) = 1$). Elementary proofs
were given independently by Erd\H{o}s and Selberg in 1949, though ``elementary''
here means ``avoiding complex analysis,'' not ``simple.''
\end{remark}

\begin{figure}[ht]
\centering
\begin{tikzpicture}
\begin{axis}[
  width=12cm, height=7cm,
  xlabel={$x$},
  ylabel={$\pi(x)$},
  xmin=0, xmax=105,
  ymin=0, ymax=30,
  grid=major,
  grid style={gray!30},
  legend pos=north west,
  legend style={font=\small},
  axis lines=left,
  every axis label/.style={font=\small},
  tick label style={font=\footnotesize}
]
% Exact pi(x) as a step function
\addplot[blue!70!black, thick, const plot]
  coordinates {
    (2,1) (3,2) (5,3) (7,4) (11,5) (13,6) (17,7) (19,8) (23,9) (29,10)
    (31,11) (37,12) (41,13) (43,14) (47,15) (53,16) (59,17) (61,18)
    (67,19) (71,20) (73,21) (79,22) (83,23) (89,24) (97,25) (101,25)
  };
\addlegendentry{$\pi(x)$ (exact)}
% x/ln(x) approximation
\addplot[red!70!black, thick, dashed, domain=2.5:105, samples=100]
  {x/ln(x)};
\addlegendentry{$x/\!\log x$}
\end{axis}
\end{tikzpicture}
\caption{The prime counting function $\pi(x)$ (blue, step function) and
  the approximation $x/\!\log x$ (red, dashed) for $x \le 100$.}
\label{fig:prime-counting}
\end{figure}

% ============================================================
\section{Special primes and open problems}
\label{sec:special-primes}

\subsection{Mersenne primes}
\label{subsec:mersenne}
\index{Mersenne primes}

\begin{definition}[Mersenne number]
\label{def:mersenne}
A \emph{Mersenne number} is an integer of the form $M_n = 2^n - 1$ for
$n \in \N^{*}$. If $M_n$ is prime, it is called a \emph{Mersenne prime}.
\end{definition}

\begin{proposition}
\label{prop:mersenne-necessary}
If $M_n = 2^n - 1$ is prime, then $n$ is prime.
\end{proposition}

\begin{proof}
Suppose $n = ab$ with $1 < a, b < n$. Then
\[
  2^n - 1 = (2^a)^b - 1 = (2^a - 1)\bigl((2^a)^{b-1} + (2^a)^{b-2} + \cdots + 1\bigr),
\]
so $2^a - 1 \mid 2^n - 1$ with $1 < 2^a - 1 < 2^n - 1$. Hence $M_n$ is composite.
\end{proof}

\begin{remark}
\label{rem:mersenne-examples}
The converse is false: $M_{11} = 2047 = 23 \times 89$. As of 2024, $52$
Mersenne primes are known; the largest, $M_{136{,}279{,}841}$, has over $41$ million digits.
Whether infinitely many Mersenne primes exist is an open problem.
\end{remark}

\subsection{Fermat primes}
\label{subsec:fermat-primes}
\index{Fermat primes}

\begin{definition}[Fermat number]
\label{def:fermat-number}
The $n$th \emph{Fermat number} is $F_n = 2^{2^n} + 1$ for $n \ge 0$.
\end{definition}

\begin{remark}
\label{rem:fermat-primes}
$F_0 = 3$, $F_1 = 5$, $F_2 = 17$, $F_3 = 257$, $F_4 = 65537$ are all prime.
Fermat conjectured that every $F_n$ is prime, but Euler showed in 1732 that
$F_5 = 4{,}294{,}967{,}297 = 641 \times 6{,}700{,}417$ is composite.
No Fermat prime beyond $F_4$ has been found.
\end{remark}

\subsection{Twin primes and the Goldbach conjecture}
\label{subsec:twin-goldbach}
\index{twin primes}\index{Goldbach conjecture}

\begin{definition}
\label{def:twin-primes}
\emph{Twin primes} are pairs $(p, p+2)$ where both $p$ and $p+2$ are prime.
\end{definition}

\begin{example}
\label{ex:twin-primes}
The first few twin prime pairs are $(3,5)$, $(5,7)$, $(11,13)$, $(17,19)$,
$(29,31)$, $(41,43)$.
\end{example}

\noindent\textbf{Twin Prime Conjecture.}\index{Twin Prime Conjecture}
There are infinitely many twin primes. This remains unproved, though
Zhang~(2013)\index{Zhang, Yitang} showed there are infinitely many pairs
$(p, p+k)$ of primes with $k \le 70{,}000{,}000$, subsequently improved to
$k \le 246$ by the Polymath project\index{Polymath project}.

\medskip
\noindent\textbf{Goldbach's Conjecture} (1742).\index{Goldbach conjecture}
Every even integer $n \ge 4$ is the sum of two primes. Verified computationally
for $n \le 4 \times 10^{18}$ and known to hold for all sufficiently large odd
integers (Helfgott, 2013: every odd $n \ge 7$ is a sum of three primes).

% ============================================================
\section{Prime spirals: Ulam and Sacks}
\label{sec:prime-spirals}
\index{Ulam spiral}\index{Sacks spiral}

In 1963, Stanislaw Ulam\index{Ulam, Stanislaw} noticed that when integers are
arranged in a spiral pattern, the primes exhibit a striking tendency to fall
along certain diagonal lines. This visual pattern, while not fully explained
theoretically, reflects the fact that certain quadratic polynomials
(e.g., $4n^2 + bn + c$) produce many primes.

\begin{figure}[ht]
\centering
\begin{tikzpicture}[scale=0.38]
  % Draw an Ulam-style spiral grid with primes highlighted
  % We encode primes up to 121 (11x11 grid centred at 1)
  \foreach \n/\x/\y in {%
    2/1/0, 3/1/1, 4/0/1, 5/-1/1, 6/-1/0, 7/-1/-1, 8/0/-1, 9/1/-1,
    10/2/-1, 11/2/0, 12/2/1, 13/2/2, 14/1/2, 15/0/2, 16/-1/2,
    17/-2/2, 18/-2/1, 19/-2/0, 20/-2/-1, 21/-2/-2, 22/-1/-2,
    23/0/-2, 24/1/-2, 25/2/-2, 26/3/-2, 27/3/-1, 28/3/0, 29/3/1,
    30/3/2, 31/3/3, 32/2/3, 33/1/3, 34/0/3, 35/-1/3, 36/-2/3,
    37/-3/3, 38/-3/2, 39/-3/1, 40/-3/0, 41/-3/-1, 42/-3/-2,
    43/-3/-3, 44/-2/-3, 45/-1/-3, 46/0/-3, 47/1/-3, 48/2/-3, 49/3/-3%
  }{%
    % Check if prime (manual list for small values)
    \pgfmathtruncatemacro{\isp}{%
      (\n==2 || \n==3 || \n==5 || \n==7 || \n==11 || \n==13 ||
       \n==17 || \n==19 || \n==23 || \n==29 || \n==31 || \n==37 ||
       \n==41 || \n==43 || \n==47) ? 1 : 0}
    \ifnum\isp=1
      \fill[blue!70!black] (\x,\y) circle (0.35);
    \else
      \fill[gray!25] (\x,\y) circle (0.2);
    \fi
  }
  % Centre: 1
  \fill[orange!80!black] (0,0) circle (0.35);
  \node[font=\tiny, white] at (0,0) {1};
\end{tikzpicture}
\caption{A small Ulam spiral: integers are placed on a spiral path starting
  from~$1$ (orange centre). Blue dots are primes; grey dots are composites.
  Even in this small example, diagonal alignments are visible.}
\label{fig:ulam-spiral}
\end{figure}

% ============================================================
\section{Why large primes matter: RSA revisited}
\label{sec:rsa-primes}
\index{RSA}\index{cryptography!RSA}

The security of RSA encryption depends on the following asymmetry:
\begin{itemize}[nosep]
  \item \textbf{Easy:} Given two large primes $p, q$, compute $n = pq$.
  \item \textbf{Hard:} Given $n$, recover $p$ and $q$ (the \emph{integer
        factorisation problem}\index{integer factorisation problem}).
\end{itemize}

The Fundamental Theorem of Arithmetic guarantees that the factorisation
\emph{exists and is unique}, but says nothing about the \emph{computational
complexity} of finding it. Current RSA implementations use primes with $1024$
or more bits (over $300$ decimal digits each), producing a modulus $n$ of
roughly $2048$ bits. No known classical algorithm factors such integers in
feasible time.

\begin{remark}[Primality testing vs.\ factoring]
\label{rem:primality-vs-factor}
\index{primality testing}\index{AKS algorithm}
It is far easier to \emph{test} whether a number is prime than to \emph{factor}
a composite number. The AKS algorithm (2002) proves primality in deterministic
polynomial time, while no polynomial-time factoring algorithm is known for
classical computers. (Shor's quantum algorithm factors in polynomial time on a
quantum computer, motivating the transition to post-quantum cryptography.)
\end{remark}

% ============================================================
\section{Exercises}
\label{sec:ch2-exercises}

\begin{exercise}[$\star$]
\label{exer:ch2-list-primes}
List all primes up to $100$ using the Sieve of Eratosthenes. How many are
there?
\end{exercise}

\begin{exercise}[$\star$]
\label{exer:ch2-fta-factor}
Find the canonical factorisation of $2520$, $10!$, and $\binom{20}{10}$.
\end{exercise}

\begin{exercise}[$\star$]
\label{exer:ch2-composite-consecutive}
Show that for every $n \ge 2$, the $n-1$ consecutive integers
$n! + 2, n! + 3, \ldots, n! + n$ are all composite.
\end{exercise}

\begin{exercise}[$\star$]
\label{exer:ch2-p-divides-binom}
Let $p$ be prime and $1 \le k \le p-1$. Prove that $p \mid \binom{p}{k}$.
\end{exercise}

\begin{exercise}[$\star\star$]
\label{exer:ch2-sqrt2-irrational}
Use the Fundamental Theorem of Arithmetic to prove that $\sqrt{2}$ is
irrational.
\end{exercise}

\begin{exercise}[$\star\star$]
\label{exer:ch2-primes-4k3}
\index{primes in arithmetic progressions}
Prove that there are infinitely many primes of the form $4k + 3$.
\textit{Hint:} adapt Euclid's argument, noting that a product of integers
$\equiv 1 \pmod{4}$ is again $\equiv 1 \pmod{4}$.
\end{exercise}

\begin{exercise}[$\star\star$]
\label{exer:ch2-legendre-formula}
\index{Legendre's formula}
Prove \textbf{Legendre's formula}: for prime $p$ and $n \ge 1$,
\[
  v_p(n!) = \sum_{i=1}^{\infty} \floor{\frac{n}{p^i}}.
\]
Use this to compute $v_2(100!)$ and $v_5(100!)$, and deduce the number of
trailing zeros of $100!$.
\end{exercise}

\begin{exercise}[$\star\star$]
\label{exer:ch2-bertrand}
\index{Bertrand's postulate}
(\textbf{Bertrand's postulate --- verification for small cases.})
Verify directly that for each $n$ with $1 \le n \le 50$, there is a prime $p$
with $n < p \le 2n$.
\end{exercise}

\begin{exercise}[$\star\star\star$]
\label{exer:ch2-euler-product}
\index{Euler product}
Prove the Euler product formula: for $s > 1$,
\[
  \sum_{n=1}^{\infty} \frac{1}{n^s}
  = \prod_{p \text{ prime}} \frac{1}{1 - p^{-s}}.
\]
\textit{Hint:} expand each factor $(1-p^{-s})^{-1} = \sum_{k=0}^{\infty}p^{-ks}$
and use the Fundamental Theorem of Arithmetic.
\end{exercise}

\begin{exercise}[$\star\star\star$]
\label{exer:ch2-chebyshev-lower}
Prove Chebyshev's lower bound: there exists $c > 0$ such that
$\pi(x) \ge c\, x / \log x$ for all $x \ge 2$.
\textit{Hint:} analyse $v_p(\binom{2n}{n})$ using Legendre's formula and bound
$\log\binom{2n}{n}$ from below.
\end{exercise}

\begin{exercise}[$\star\star\star$]
\label{exer:ch2-prime-gaps}
\index{prime gaps}
Prove that for every $\varepsilon > 0$, there exist consecutive primes
$p_n, p_{n+1}$ with $p_{n+1} - p_n > \varepsilon \log p_n$.
\textit{Hint:} if all gaps were $\le \varepsilon \log p_n$, sum to get a
contradiction with the Prime Number Theorem.
\end{exercise}

% ============================================================
\section*{Chapter summary}
\addcontentsline{toc}{section}{Chapter summary}

\begin{itemize}[nosep]
  \item Every integer $> 1$ has a prime divisor; a composite $n$ has a prime
        factor $\le \sqrt{n}$.
  \item \textbf{Infinitude of primes:} proved by Euclid (contradiction),
        Euler ($\sum 1/p$ diverges), and via Fermat numbers (pairwise coprime).
  \item The \textbf{Fundamental Theorem of Arithmetic} asserts that every
        integer $> 1$ has a \emph{unique} prime factorisation. The proof
        of uniqueness relies on Euclid's lemma.
  \item The \textbf{Sieve of Eratosthenes} finds all primes up to $N$ in
        $O(N \log\log N)$ time.
  \item The \textbf{Prime Number Theorem}: $\pi(x) \sim x/\log x$. Chebyshev
        obtained the correct order of growth; Hadamard and de la Vall\'{e}e-Poussin
        proved the asymptotic.
  \item \textbf{Mersenne primes} ($2^p - 1$) and \textbf{Fermat primes}
        ($2^{2^n}+1$) are special families; their infinitude is unknown.
  \item The \textbf{Twin Prime Conjecture} and \textbf{Goldbach's Conjecture}
        remain major open problems.
  \item \textbf{RSA cryptography} depends on the ease of primality testing
        versus the difficulty of factoring large composites.
\end{itemize}

% === END OF CHUNK preamble+ch1-2 ===
% === START OF CHUNK ch3-4 ===

% ============================================================
%  CHAPTER 3 — Congruences and Modular Arithmetic
% ============================================================

\chapter{Congruences and Modular Arithmetic}
\label{ch:congruences}

\section{Historical Context: Gauss's \textit{Disquisitiones Arithmeticae}}
\label{sec:gauss-disquisitiones}
\index{Gauss, Carl Friedrich}
\index{Disquisitiones Arithmeticae@\textit{Disquisitiones Arithmeticae}}

In 1801, Carl Friedrich Gauss, then only twenty-four years old, published
his masterwork \textit{Disquisitiones Arithmeticae}\index{Disquisitiones Arithmeticae@\textit{Disquisitiones Arithmeticae}}.
This treatise transformed number theory from a collection of scattered
results into a systematic and rigorous discipline.  The very first section
of the \textit{Disquisitiones} introduces the concept of \emph{congruence},
a relation that Gauss recognised as the natural language for discussing
divisibility and residues.

Gauss wrote:
\begin{quote}
\itshape
If a number~$a$ divides the difference of the numbers~$b$ and~$c$,
$b$ and~$c$ are said to be \emph{congruent} with respect to~$a$;
if not, \emph{incongruent}.  The number~$a$ is called the \emph{modulus}.
\end{quote}

\noindent
He introduced the now-universal notation $b \equiv c \pmod{a}$ and
proceeded to develop the arithmetic of congruences with the same
rigour that Euclid had applied to geometry.  The ideas in this chapter
follow the path laid out in the \textit{Disquisitiones}, enriched by
the algebraic perspective that emerged over the following two centuries.

% ------------------------------------------------------------
\section{The Congruence Relation}
\label{sec:congruence-def}

\begin{definition}[Congruence modulo~$n$]
\label{def:congruence}
\index{congruence!definition}
Let $n$ be a positive integer.  Two integers $a$ and $b$ are said to be
\emph{congruent modulo~$n$}, written
\[
  a \equiv b \pmod{n},
\]
if $n \mid (a - b)$, that is, if there exists an integer $k$ such that
$a - b = kn$.  The integer~$n$ is called the \emph{modulus}\index{modulus}
of the congruence.
\end{definition}

\begin{example}[Basic congruences]
\label{ex:basic-cong}
\leavevmode
\begin{enumerate}[label=(\roman*)]
  \item $17 \equiv 2 \pmod{5}$, since $17 - 2 = 15 = 3 \cdot 5$.
  \item $-3 \equiv 4 \pmod{7}$, since $-3 - 4 = -7 = (-1)\cdot 7$.
  \item For any integer~$a$, $a \equiv a \bmod n \pmod{n}$, where $a \bmod n$
        denotes the remainder upon division by~$n$.
\end{enumerate}
\end{example}

\begin{theorem}[Congruence is an equivalence relation]
\label{thm:cong-equiv}
\index{congruence!equivalence relation}
For every positive integer~$n$, the relation $\equiv \pmod{n}$ is an
equivalence relation on~$\Z$.  That is, for all $a,b,c \in \Z$:
\begin{enumerate}[label=(\roman*)]
  \item \textbf{Reflexivity:} $a \equiv a \pmod{n}$.
  \item \textbf{Symmetry:} If $a \equiv b \pmod{n}$, then $b \equiv a \pmod{n}$.
  \item \textbf{Transitivity:} If $a \equiv b \pmod{n}$ and $b \equiv c \pmod{n}$,
        then $a \equiv c \pmod{n}$.
\end{enumerate}
\end{theorem}

\begin{proof}
\leavevmode
\begin{enumerate}[label=(\roman*)]
  \item Since $a - a = 0 = 0 \cdot n$, we have $n \mid 0$, so $a \equiv a \pmod{n}$.
  \item If $n \mid (a-b)$, write $a - b = kn$.  Then $b - a = (-k)n$, so $n \mid (b-a)$.
  \item If $a - b = k_1 n$ and $b - c = k_2 n$, then
        $a - c = (a - b) + (b - c) = (k_1 + k_2)n$, so $n \mid (a - c)$.
        \qedhere
\end{enumerate}
\end{proof}

The equivalence classes are the \emph{residue classes}\index{residue class} modulo~$n$.
The class containing the integer~$a$ is
\[
  \bar{a} = [a]_n = \{ a + kn : k \in \Z \}.
\]
There are exactly $n$ distinct residue classes, represented by
$\{0, 1, 2, \dots, n-1\}$.

\begin{theorem}[Compatibility with arithmetic]
\label{thm:cong-arith}
\index{congruence!arithmetic properties}
Let $a \equiv a' \pmod{n}$ and $b \equiv b' \pmod{n}$.  Then:
\begin{enumerate}[label=(\roman*)]
  \item $a + b \equiv a' + b' \pmod{n}$.
  \item $a - b \equiv a' - b' \pmod{n}$.
  \item $ab \equiv a'b' \pmod{n}$.
  \item For every non-negative integer~$k$, $a^k \equiv (a')^k \pmod{n}$.
\end{enumerate}
\end{theorem}

\begin{proof}
Write $a = a' + sn$ and $b = b' + tn$ for some $s,t \in \Z$.
\begin{enumerate}[label=(\roman*)]
  \item $a + b = a' + b' + (s + t)n$, so $n \mid (a + b - a' - b')$.
  \item Identical reasoning with $s - t$.
  \item $ab = (a' + sn)(b' + tn) = a'b' + (a't + b's + stn)n$, so
        $n \mid (ab - a'b')$.
  \item Induction on~$k$ using~(iii): the base case $k = 0$ gives $1 \equiv 1$,
        and if $a^k \equiv (a')^k$ then
        $a^{k+1} = a \cdot a^k \equiv a' \cdot (a')^k = (a')^{k+1} \pmod{n}$.
        \qedhere
\end{enumerate}
\end{proof}

\begin{remark}[Cancellation requires care]
\label{rem:cancel-care}
\index{congruence!cancellation}
Unlike ordinary equations, one cannot freely cancel a common factor.
For instance, $2 \cdot 3 \equiv 2 \cdot 0 \pmod{6}$, but $3 \not\equiv 0
\pmod{6}$.  Cancellation of~$a$ from $ab \equiv ac \pmod{n}$ is valid
precisely when $\gcd(a, n) = 1$.
\end{remark}

\begin{proposition}[Cancellation law]
\label{prop:cancel}
\index{congruence!cancellation law}
If $ab \equiv ac \pmod{n}$ and $d = \gcd(a,n)$, then
$b \equiv c \pmod{n/d}$.  In particular, if $\gcd(a,n)=1$ then
$b \equiv c \pmod{n}$.
\end{proposition}

\begin{proof}
From $ab \equiv ac \pmod{n}$ we get $n \mid a(b-c)$.  Writing $a = da'$,
$n = dn'$ with $\gcd(a', n') = 1$, we have $dn' \mid da'(b-c)$, so
$n' \mid a'(b-c)$.  Since $\gcd(a', n') = 1$, Euclid's lemma gives
$n' \mid (b - c)$, i.e.\ $b \equiv c \pmod{n/d}$.
\end{proof}

% ------------------------------------------------------------
\section{The Ring $\Z/n\Z$}
\label{sec:ZnZ}
\index{Z/nZ@$\Z/n\Z$}

\begin{definition}[The ring $\Z/n\Z$]
\label{def:ZnZ}
The set of residue classes modulo~$n$,
\[
  \Z/n\Z = \{ \bar{0}, \bar{1}, \dots, \overline{n-1} \},
\]
equipped with the operations
\[
  \bar{a} + \bar{b} = \overline{a+b}, \qquad
  \bar{a} \cdot \bar{b} = \overline{ab},
\]
forms a commutative ring\index{ring!commutative} with identity $\bar{1}$.
\end{definition}

The well-definedness of these operations is precisely the content of
Theorem~\ref{thm:cong-arith}.

\begin{definition}[Units and zero divisors]
\label{def:units-zerodiv}
\index{unit}\index{zero divisor}
An element $\bar{a} \in \Z/n\Z$ is a \emph{unit} (or \emph{invertible element})
if there exists $\bar{b} \in \Z/n\Z$ with $\bar{a}\bar{b} = \bar{1}$.
An element $\bar{a} \neq \bar{0}$ is a \emph{zero divisor} if there exists
$\bar{b} \neq \bar{0}$ with $\bar{a}\bar{b} = \bar{0}$.
\end{definition}

\begin{theorem}[Units and zero divisors in $\Z/n\Z$]
\label{thm:units-ZnZ}
\index{Z/nZ@$\Z/n\Z$!units}
Let $\bar{a} \in \Z/n\Z$ with $a \not\equiv 0 \pmod{n}$.
\begin{enumerate}[label=(\roman*)]
  \item $\bar{a}$ is a unit if and only if $\gcd(a,n)=1$.
  \item $\bar{a}$ is a zero divisor if and only if $\gcd(a,n)>1$.
\end{enumerate}
In particular, every nonzero element of $\Z/n\Z$ is either a unit or
a zero divisor.
\end{theorem}

\begin{proof}
\leavevmode
\begin{enumerate}[label=(\roman*)]
  \item $\bar{a}$ is a unit iff $\exists\, b$ with $ab \equiv 1 \pmod{n}$,
        iff the equation $ax + ny = 1$ has a solution in integers,
        iff $\gcd(a,n) = 1$ (by Bézout's identity).
  \item If $d = \gcd(a,n) > 1$, set $b = n/d$.  Then $1 \le b < n$ and
        $ab = a(n/d) = (a/d)\cdot n \equiv 0 \pmod{n}$, so $\bar{a}$ is a zero
        divisor.  Conversely, if $\gcd(a,n)=1$ then $\bar{a}$ is a unit, hence
        $\bar{a}\bar{b} = \bar{0}$ implies $\bar{b} = \bar{0}$, so $\bar{a}$
        is not a zero divisor.
        \qedhere
\end{enumerate}
\end{proof}

\begin{corollary}[$\Z/n\Z$ is a field iff $n$ is prime]
\label{cor:ZpZ-field}
\index{Z/pZ@$\Z/p\Z$!field}\index{field}
The ring $\Z/n\Z$ is a field (i.e.\ every nonzero element is a unit)
if and only if $n$ is prime.
\end{corollary}

\begin{proof}
If $n = p$ is prime, then for $1 \le a \le p-1$ we have $\gcd(a,p)=1$,
so $\bar{a}$ is a unit.  Conversely, if $n$ is composite, say $n = ab$
with $1 < a,b < n$, then $\bar{a}\bar{b} = \bar{0}$ with
$\bar{a}, \bar{b} \neq \bar{0}$, so $\Z/n\Z$ has zero divisors and
is not a field.
\end{proof}

\begin{notation}[Group of units]
\label{not:units-group}
\index{Z/nZ@$(\Z/n\Z)^\times$}
The group of units of $\Z/n\Z$ is denoted $(\Z/n\Z)^\times$.  Its order
is given by Euler's totient function~$\varphi(n)$, which we study
in Chapter~\ref{ch:fermat-euler-wilson}.
\end{notation}

% --- TikZ: clock arithmetic ---
\begin{figure}[ht]
\centering
\begin{tikzpicture}[scale=1.1]
  % Draw clock for Z/12Z
  \def\n{12}
  \foreach \i in {0,...,11} {
    \pgfmathsetmacro{\angle}{90 - \i*360/\n}
    \node[circle, draw, fill=blue!10, inner sep=3pt, minimum size=22pt,
          font=\small\bfseries] (n\i) at (\angle:2.6) {$\i$};
  }
  % Draw circle
  \draw[gray!50, thick] (0,0) circle (2.6);
  % Highlight addition: 5 + 9 = 2 (mod 12)
  \node[circle, draw, fill=red!30, inner sep=3pt, minimum size=22pt,
        font=\small\bfseries] at (90 - 5*30:2.6) {$5$};
  \node[circle, draw, fill=orange!30, inner sep=3pt, minimum size=22pt,
        font=\small\bfseries] at (90 - 9*30:2.6) {$9$};
  \node[circle, draw, fill=green!30, inner sep=3pt, minimum size=22pt,
        font=\small\bfseries] at (90 - 2*30:2.6) {$2$};
  % Arrow from 5 going 9 steps
  \draw[-{Stealth}, thick, red!70!black, bend left=15]
    (90 - 5*30:3.15) to node[midway, above right, font=\scriptsize] {$+9$}
    (90 - 2*30:3.15);
  \node[below] at (0,-3.2) {Clock arithmetic: $5 + 9 \equiv 2 \pmod{12}$};
\end{tikzpicture}
\caption{Addition in $\Z/12\Z$ visualised as clock arithmetic.}
\label{fig:clock-arith}
\index{clock arithmetic}
\end{figure}

% --- TikZ: multiplication table mod 7 ---
\begin{figure}[ht]
\centering
\begin{tikzpicture}[scale=0.72]
  \def\n{7}
  % Header row
  \node[font=\bfseries] at (0,0.8) {$\times$};
  \foreach \j in {1,...,6} {
    \node[font=\bfseries] at (\j*1.2, 0.8) {$\j$};
  }
  \draw[thick] (-0.5, 0.45) -- (7.7, 0.45);
  % Rows
  \foreach \i in {1,...,6} {
    \node[font=\bfseries] at (0, -\i*0.7 + 0.7) {$\i$};
    \foreach \j in {1,...,6} {
      \pgfmathtruncatemacro{\val}{mod(\i*\j,\n)}
      \ifnum\val=1
        \node[fill=green!25, rounded corners=2pt, inner sep=2pt]
          at (\j*1.2, -\i*0.7 + 0.7) {$\val$};
      \else
        \node at (\j*1.2, -\i*0.7 + 0.7) {$\val$};
      \fi
    }
  }
  \draw[thick] (0.5, 1.15) -- (0.5, -3.9);
  \node[below, font=\small] at (3.6, -4.2)
    {Multiplication table of $(\Z/7\Z)^\times$; entries equal to $1$ are highlighted.};
\end{tikzpicture}
\caption{Multiplication table of $(\Z/7\Z)^\times$.}
\label{fig:mult-table-7}
\index{multiplication table}
\end{figure}

% ------------------------------------------------------------
\section{Linear Congruences}
\label{sec:linear-cong}
\index{linear congruence}

The equation $ax \equiv b \pmod{n}$ is the modular analogue of the
linear equation $ax = b$.  Its solvability depends on the relationship
between $a$, $b$, and~$n$.

\begin{theorem}[Existence and number of solutions]
\label{thm:linear-cong}
\index{linear congruence!solution theory}
Let $a,b \in \Z$ and $n \ge 1$.  Set $d = \gcd(a,n)$.
\begin{enumerate}[label=(\roman*)]
  \item The congruence $ax \equiv b \pmod{n}$ has a solution if and only
        if $d \mid b$.
  \item When $d \mid b$, there are exactly $d$ incongruent solutions modulo~$n$.
        If $x_0$ is one solution, the complete set of solutions is
        \[
          x \equiv x_0 + k \cdot \frac{n}{d} \pmod{n}, \qquad k = 0, 1, \dots, d-1.
        \]
\end{enumerate}
\end{theorem}

\begin{proof}
\leavevmode
\begin{enumerate}[label=(\roman*)]
  \item The congruence $ax \equiv b \pmod{n}$ is equivalent to the
        existence of integers $x, y$ with $ax - ny = b$, i.e.\ $ax + n(-y) = b$.
        By the characterisation of the image of the map
        $(x,y) \mapsto ax + ny$ (which equals $d\Z$), this has a solution
        if and only if $d \mid b$.

  \item Suppose $d \mid b$.  Dividing $ax \equiv b \pmod{n}$ through by~$d$ yields
        \[
          \frac{a}{d}\, x \equiv \frac{b}{d} \pmod{\frac{n}{d}},
        \]
        where $\gcd(a/d,\, n/d) = 1$.  By Proposition~\ref{prop:cancel} (or by
        Bézout), $a/d$ has an inverse modulo~$n/d$, so there is a unique
        solution $x_0$ modulo~$n/d$.

        The solutions modulo~$n$ that reduce to~$x_0$ modulo~$n/d$ are
        precisely
        \[
          x_0,\; x_0 + \frac{n}{d},\; x_0 + 2\cdot\frac{n}{d},\;
          \dots,\; x_0 + (d-1)\cdot\frac{n}{d},
        \]
        and these are $d$ values that are pairwise incongruent modulo~$n$
        (since their differences are non-zero multiples of $n/d$ that are
        less than~$n$).
        \qedhere
\end{enumerate}
\end{proof}

\begin{example}[Solving a linear congruence]
\label{ex:linear-cong}
Solve $12x \equiv 9 \pmod{15}$.

We have $d = \gcd(12, 15) = 3$ and $3 \mid 9$, so solutions exist.
Dividing by~$3$: $4x \equiv 3 \pmod{5}$.  Since $4 \cdot 4 = 16 \equiv 1
\pmod{5}$, the inverse of~$4$ modulo~$5$ is~$4$.  Thus
$x \equiv 4 \cdot 3 = 12 \equiv 2 \pmod{5}$.

The three solutions modulo~$15$ are
\[
  x \equiv 2, \quad x \equiv 7, \quad x \equiv 12 \pmod{15}.
\]
\end{example}

\begin{example}[No solution]
\label{ex:linear-cong-no-sol}
The congruence $6x \equiv 5 \pmod{9}$ has no solution because
$\gcd(6,9) = 3$ and $3 \nmid 5$.
\end{example}

% ------------------------------------------------------------
\section{Divisibility Criteria via Congruences}
\label{sec:divisibility-criteria}
\index{divisibility criteria}

Congruences provide elegant proofs of the familiar divisibility tests
learned in school.  Let $N$ be a positive integer with decimal
representation $N = a_k 10^k + a_{k-1} 10^{k-1} + \cdots + a_1 \cdot 10 + a_0$.

\begin{proposition}[Divisibility by $3$ and $9$]
\label{prop:div-3-9}
\index{divisibility criteria!by 3 and 9}
$N \equiv a_k + a_{k-1} + \cdots + a_0 \pmod{9}$.
In particular, $3 \mid N$ iff $3$ divides the digit sum, and
$9 \mid N$ iff $9$ divides the digit sum.
\end{proposition}

\begin{proof}
Since $10 \equiv 1 \pmod{9}$, we have $10^j \equiv 1 \pmod{9}$ for all $j \ge 0$.
Hence $N \equiv \sum_{j=0}^{k} a_j \cdot 1 = \sum a_j \pmod{9}$.
The criterion for~$3$ follows because $3 \mid 9$.
\end{proof}

\begin{proposition}[Divisibility by $11$]
\label{prop:div-11}
\index{divisibility criteria!by 11}
$N \equiv a_0 - a_1 + a_2 - \cdots + (-1)^k a_k \pmod{11}$.
Hence $11 \mid N$ if and only if $11$ divides the alternating digit sum.
\end{proposition}

\begin{proof}
Since $10 \equiv -1 \pmod{11}$, we have $10^j \equiv (-1)^j \pmod{11}$, so
$N \equiv \sum_{j=0}^{k} (-1)^j a_j \pmod{11}$.
\end{proof}

\begin{proposition}[Divisibility by $7$, $11$, and $13$]
\label{prop:div-7-11-13}
\index{divisibility criteria!by 7, 11, 13}
Since $7 \cdot 11 \cdot 13 = 1001$ and $10^3 \equiv -1 \pmod{1001}$,
one can test divisibility by~$7$, $11$, or~$13$ by forming
alternating sums of three-digit blocks from right to left.
\end{proposition}

\begin{example}
\label{ex:div-criteria}
Consider $N = 2\,358\,271$.  The three-digit blocks from the right are
$271$, $358$, $2$.  The alternating sum is $271 - 358 + 2 = -85$.
Since $-85 = -5 \cdot 17$, we see that $N$ is not divisible by~$7$,
$11$, or~$13$.
\end{example}

% ------------------------------------------------------------
\section{Fast Modular Exponentiation}
\label{sec:fast-exp}
\index{modular exponentiation}\index{exponentiation!fast modular}

Computing $a^k \bmod n$ by repeated multiplication requires $k-1$
multiplications, which is infeasible when $k$ is hundreds of digits long
(as in cryptographic applications).  The method of \emph{repeated squaring}
(also known as \emph{binary exponentiation}\index{binary exponentiation})
reduces the cost to at most $2\floor{\log_2 k}$ multiplications modulo~$n$.

\begin{definition}[Binary exponentiation]
\label{def:binary-exp}
Write $k$ in binary: $k = (b_\ell b_{\ell-1} \cdots b_1 b_0)_2$ with
$b_\ell = 1$.  Then
\[
  a^k = a^{2^\ell b_\ell + \cdots + 2 b_1 + b_0}
      = \prod_{i=0}^{\ell} \bigl(a^{2^i}\bigr)^{b_i}.
\]
Each successive $a^{2^i}$ is obtained by squaring the previous one.
\end{definition}

\begin{example}[Computing $3^{45} \bmod 67$]
\label{ex:fast-exp}
We have $45 = (101101)_2 = 32 + 8 + 4 + 1$.  Build the table:
\[
\begin{array}{r@{\;}l}
  3^1  &\equiv 3 \pmod{67},  \\
  3^2  &\equiv 9,  \\
  3^4  &\equiv 81 \equiv 14,  \\
  3^8  &\equiv 14^2 = 196 \equiv 62 \equiv -5,  \\
  3^{16} &\equiv 25,  \\
  3^{32} &\equiv 625 \equiv 625 - 9\cdot 67 = 625 - 603 = 22.
\end{array}
\]
Therefore $3^{45} \equiv 3^{32} \cdot 3^{8} \cdot 3^{4} \cdot 3^{1}
\equiv 22 \cdot (-5) \cdot 14 \cdot 3 \pmod{67}$.

Computing step by step: $22 \cdot (-5) = -110 \equiv -110 + 2\cdot 67 = 24$,
then $24 \cdot 14 = 336 \equiv 336 - 5\cdot 67 = 336 - 335 = 1$,
then $1 \cdot 3 = 3$.  So $3^{45} \equiv 3 \pmod{67}$.
\end{example}

\begin{remark}[Complexity]
\label{rem:fast-exp-complexity}
The binary exponentiation algorithm computes $a^k \bmod n$ using
$O(\log k)$ modular multiplications, each of which costs $O((\log n)^2)$
bit operations (using schoolbook multiplication).  The total cost is thus
$O((\log k)(\log n)^2)$ bit operations.
\end{remark}

% ------------------------------------------------------------
\section{Connections to Cryptography}
\label{sec:cong-crypto}
\index{cryptography!modular arithmetic in}

Modular arithmetic is the computational backbone of modern public-key
cryptography.  We highlight two points; a full treatment of RSA and
Diffie--Hellman appears in Chapter~\ref{ch:fermat-euler-wilson}.

\begin{enumerate}
  \item \textbf{One-way functions.}\index{one-way function}
        The map $x \mapsto g^x \bmod p$ (for a prime~$p$ and
        generator~$g$ of $(\Z/p\Z)^\times$) is easy to compute via
        repeated squaring, but inverting it---the
        \emph{discrete logarithm problem}\index{discrete logarithm problem}---is
        believed to be computationally hard.

  \item \textbf{Modular inverses in RSA.}
        The RSA cryptosystem relies on computing modular inverses
        via the extended Euclidean algorithm, and on the difficulty
        of factoring a product $n = pq$ of two large primes.
\end{enumerate}

% ------------------------------------------------------------
\section{Exercises}
\label{sec:ch3-exercises}

\begin{exercise}
\label{exer:cong-basic}
Determine the last two digits of $7^{999}$ (i.e.\ compute $7^{999} \bmod 100$).
\end{exercise}

\begin{exercise}
\label{exer:linear-cong-1}
Find all solutions of $35x \equiv 14 \pmod{91}$.
\end{exercise}

\begin{exercise}
\label{exer:linear-cong-2}
Prove that if $p$ is an odd prime and $a \not\equiv 0 \pmod{p}$, then
$ax \equiv b \pmod{p}$ has a unique solution modulo~$p$.
\end{exercise}

\begin{exercise}
\label{exer:ZnZ-units}
List all units and all zero divisors in $\Z/12\Z$.  For each unit, find
its multiplicative inverse.
\end{exercise}

\begin{exercise}
\label{exer:divisibility-37}
\index{divisibility criteria!by 37}
Using $10^3 \equiv 1 \pmod{37}$, derive a test for divisibility by~$37$
based on splitting the decimal representation into three-digit blocks.
\end{exercise}

\begin{exercise}
\label{exer:fast-exp}
Use binary exponentiation to compute $5^{117} \bmod 19$.
\end{exercise}

\begin{exercise}
\label{exer:system-cong}
Solve the system of congruences $x \equiv 3 \pmod{5}$, $x \equiv 4 \pmod{7}$.
(This foreshadows the Chinese Remainder Theorem.)
\end{exercise}

\begin{exercise}
\label{exer:nilpotent}
Let $n = p^k$ for a prime~$p$ and $k \ge 1$.  Prove that $\bar{a} \in \Z/n\Z$
is nilpotent (i.e.\ $\bar{a}^m = \bar{0}$ for some $m \ge 1$) if and only
if $p \mid a$.
\end{exercise}

\begin{exercise}
\label{exer:idempotent}
Find all idempotent elements ($e^2 = e$) in $\Z/n\Z$ for
$n = 6, 10, 15$.  What pattern do you observe?
\end{exercise}

\begin{exercise}[Challenge]
\label{exer:cong-challenge}
Prove that for every $n \ge 2$, the product of all units in $(\Z/n\Z)^\times$
equals $\bar{-1}$ if $n = 1, 2, 4, p^k$, or $2p^k$ (for odd prime~$p$),
and equals $\bar{1}$ otherwise.
\end{exercise}

% ------------------------------------------------------------
\section{Chapter Summary}
\label{sec:ch3-summary}

\begin{enumerate}
  \item \textbf{Congruence} $a \equiv b \pmod{n}$ means $n \mid (a-b)$.
        It is an equivalence relation compatible with addition and multiplication.
  \item The quotient ring $\Z/n\Z$ consists of $n$ residue classes.
        Its units are the classes $\bar{a}$ with $\gcd(a,n)=1$; the remaining
        nonzero classes are zero divisors.  $\Z/n\Z$ is a field iff $n$ is prime.
  \item The linear congruence $ax \equiv b \pmod{n}$ has solutions iff
        $\gcd(a,n) \mid b$, and then exactly $\gcd(a,n)$ incongruent
        solutions.
  \item Classical divisibility tests follow from $10 \equiv 1 \pmod{9}$,
        $10 \equiv -1 \pmod{11}$, etc.
  \item Binary exponentiation computes $a^k \bmod n$ in $O(\log k)$
        multiplications, enabling efficient cryptographic computations.
\end{enumerate}


% ============================================================
%  CHAPTER 4 — Theorems of Fermat, Euler, and Wilson
% ============================================================

\chapter{Theorems of Fermat, Euler, and Wilson}
\label{ch:fermat-euler-wilson}

\section{Historical Context: Fermat's Letter to Frénicle}
\label{sec:fermat-frenicle}
\index{Fermat, Pierre de}\index{Frénicle de Bessy}

On 18~October 1640, Pierre de Fermat wrote a letter to Bernard Frénicle de
Bessy in which he stated, without proof, a remarkable observation:
\begin{quote}
\itshape
If $p$ is prime and $a$ is not divisible by~$p$, then $a^{p-1} - 1$ is
always divisible by~$p$.
\end{quote}
Fermat added, characteristically, ``I would send you the proof, if I did
not fear it being too long.''  The first published proof was given by
Euler\index{Euler, Leonhard} in 1736, who later generalised the result to
composite moduli using what we now call \emph{Euler's totient function}.

Wilson's theorem\index{Wilson, John}, stating that $(p-1)! \equiv -1
\pmod{p}$ for every prime~$p$, was first conjectured by the English
mathematician John Wilson around~1770.  Lagrange\index{Lagrange, Joseph-Louis}
provided the first proof in~1771.

This chapter presents complete proofs of these classical theorems and
explores their far-reaching applications, including the RSA cryptosystem
and the Diffie--Hellman key exchange protocol.

% ------------------------------------------------------------
\section{Euler's Totient Function}
\label{sec:euler-totient}
\index{Euler's totient function}\index{totient function|see{Euler's totient function}}
\index{phi@$\varphi(n)$|see{Euler's totient function}}

\begin{definition}[Euler's totient function]
\label{def:euler-totient}
For $n \ge 1$, \emph{Euler's totient function} $\varphi(n)$\index{phi@$\varphi(n)$}
is the number of integers~$a$ with $1 \le a \le n$ and $\gcd(a,n) = 1$.
Equivalently,
\[
  \varphi(n) = \abs{(\Z/n\Z)^\times}.
\]
\end{definition}

\begin{example}[Small values]
\label{ex:totient-small}
$\varphi(1) = 1$, $\varphi(2) = 1$, $\varphi(3) = 2$, $\varphi(4) = 2$,
$\varphi(5) = 4$, $\varphi(6) = 2$, $\varphi(7) = 6$, $\varphi(8) = 4$,
$\varphi(9) = 6$, $\varphi(10) = 4$, $\varphi(12) = 4$.
\end{example}

\begin{proposition}[Totient of a prime]
\label{prop:totient-prime}
\index{Euler's totient function!of a prime}
If $p$ is prime, then $\varphi(p) = p - 1$.
\end{proposition}

\begin{proof}
Every integer $a$ with $1 \le a \le p-1$ satisfies $\gcd(a,p) = 1$.
\end{proof}

\begin{proposition}[Totient of a prime power]
\label{prop:totient-prime-power}
\index{Euler's totient function!of a prime power}
For a prime~$p$ and $k \ge 1$,
\[
  \varphi(p^k) = p^k - p^{k-1} = p^{k-1}(p-1) = p^k\Bigl(1 - \frac{1}{p}\Bigr).
\]
\end{proposition}

\begin{proof}
Among $\{1, 2, \dots, p^k\}$, the integers \emph{not} coprime to~$p^k$
are exactly the multiples of~$p$: these are $p, 2p, 3p, \dots, p^{k-1}\cdot p$,
and there are $p^{k-1}$ of them.  Hence $\varphi(p^k) = p^k - p^{k-1}$.
\end{proof}

\begin{theorem}[Multiplicativity of $\varphi$]
\label{thm:totient-mult}
\index{Euler's totient function!multiplicativity}
\index{multiplicative function}
If $\gcd(m,n) = 1$, then $\varphi(mn) = \varphi(m)\varphi(n)$.
\end{theorem}

\begin{proof}
By the Chinese Remainder Theorem\index{Chinese Remainder Theorem},
$\Z/mn\Z \cong \Z/m\Z \times \Z/n\Z$ as rings when $\gcd(m,n) = 1$.
This isomorphism restricts to an isomorphism of unit groups:
\[
  (\Z/mn\Z)^\times \cong (\Z/m\Z)^\times \times (\Z/n\Z)^\times.
\]
Taking cardinalities gives $\varphi(mn) = \varphi(m)\varphi(n)$.
\end{proof}

\begin{theorem}[Product formula for $\varphi$]
\label{thm:totient-formula}
\index{Euler's totient function!product formula}
For any $n \ge 2$ with prime factorisation $n = p_1^{k_1} p_2^{k_2} \cdots p_r^{k_r}$,
\[
  \varphi(n) = n \prod_{p \mid n} \Bigl(1 - \frac{1}{p}\Bigr)
             = n \cdot \frac{p_1 - 1}{p_1} \cdot \frac{p_2 - 1}{p_2}
               \cdots \frac{p_r - 1}{p_r}.
\]
\end{theorem}

\begin{proof}
Combine multiplicativity (Theorem~\ref{thm:totient-mult}) with the
prime-power formula (Proposition~\ref{prop:totient-prime-power}):
\[
  \varphi(n)
  = \prod_{i=1}^{r} \varphi(p_i^{k_i})
  = \prod_{i=1}^{r} p_i^{k_i}\Bigl(1 - \frac{1}{p_i}\Bigr)
  = n \prod_{i=1}^{r} \Bigl(1 - \frac{1}{p_i}\Bigr).
  \qedhere
\]
\end{proof}

\begin{example}[Computing $\varphi(360)$]
\label{ex:totient-360}
We have $360 = 2^3 \cdot 3^2 \cdot 5$, so
\[
  \varphi(360) = 360 \cdot \Bigl(1 - \frac{1}{2}\Bigr)
                      \Bigl(1 - \frac{1}{3}\Bigr)
                      \Bigl(1 - \frac{1}{5}\Bigr)
               = 360 \cdot \frac{1}{2} \cdot \frac{2}{3} \cdot \frac{4}{5}
               = 96.
\]
\end{example}

\begin{theorem}[Gauss's divisor sum for $\varphi$]
\label{thm:gauss-divisor-sum}
\index{Euler's totient function!divisor sum}
For every $n \ge 1$,
\[
  \sum_{d \mid n} \varphi(d) = n.
\]
\end{theorem}

\begin{proof}
Partition $\{1, 2, \dots, n\}$ according to $\gcd(a, n)$.  For each
divisor~$d$ of~$n$, the number of integers~$a$ with $1 \le a \le n$ and
$\gcd(a,n) = d$ equals the number of integers~$b$ with $1 \le b \le n/d$
and $\gcd(b, n/d) = 1$, which is $\varphi(n/d)$.  Summing over all
divisors~$d$ of~$n$:
\[
  n = \sum_{d \mid n} \varphi\!\left(\frac{n}{d}\right)
    = \sum_{d \mid n} \varphi(d),
\]
where the last equality holds because as~$d$ ranges over divisors of~$n$,
so does~$n/d$.
\end{proof}

% ------------------------------------------------------------
\section{Euler's Theorem}
\label{sec:euler-thm}
\index{Euler's theorem}

\begin{theorem}[Euler's theorem]
\label{thm:euler}
If $n \ge 1$ and $\gcd(a,n) = 1$, then
\[
  a^{\varphi(n)} \equiv 1 \pmod{n}.
\]
\end{theorem}

\begin{proof}
\label{proof:euler}
Let $(\Z/n\Z)^\times = \{r_1, r_2, \dots, r_{\varphi(n)}\}$ be the group of
units.  Since $\gcd(a,n) = 1$, multiplication by~$\bar{a}$ is a
bijection\index{bijection!on units} from $(\Z/n\Z)^\times$ to itself.
Therefore the set $\{ar_1, ar_2, \dots, ar_{\varphi(n)}\}$ is simply a
rearrangement of $\{r_1, r_2, \dots, r_{\varphi(n)}\}$ modulo~$n$.

Taking the product of all elements in each set:
\[
  \prod_{i=1}^{\varphi(n)} (ar_i) \equiv \prod_{i=1}^{\varphi(n)} r_i \pmod{n}.
\]
The left side equals $a^{\varphi(n)} \prod_{i=1}^{\varphi(n)} r_i$.  Let
$P = \prod_{i=1}^{\varphi(n)} r_i$.  Then
\[
  a^{\varphi(n)} \cdot P \equiv P \pmod{n}.
\]
Since each~$r_i$ is a unit, $P$ is also a unit and may be cancelled,
yielding $a^{\varphi(n)} \equiv 1 \pmod{n}$.
\end{proof}

\begin{remark}
\label{rem:euler-group-theory}
From the viewpoint of group theory, Euler's theorem is an immediate consequence
of Lagrange's theorem\index{Lagrange's theorem}: the order of any element
of a finite group divides the order of the group.  Applied to the element
$\bar{a}$ in $(\Z/n\Z)^\times$, whose order is $\varphi(n)$, we obtain
$\bar{a}^{\varphi(n)} = \bar{1}$.
\end{remark}

% ------------------------------------------------------------
\section{Fermat's Little Theorem}
\label{sec:fermat-little}
\index{Fermat's little theorem}

\begin{corollary}[Fermat's little theorem]
\label{cor:fermat-little}
If $p$ is prime and $p \nmid a$, then
\[
  a^{p-1} \equiv 1 \pmod{p}.
\]
Equivalently, for any integer~$a$,
\[
  a^p \equiv a \pmod{p}.
\]
\end{corollary}

\begin{proof}
The first form follows from Euler's theorem with $n = p$, since
$\varphi(p) = p-1$.  For the second form: if $p \mid a$, both sides are
$\equiv 0$; if $p \nmid a$, multiply the first form by~$a$.
\end{proof}

\begin{example}[Application of Fermat]
\label{ex:fermat-app}
Compute $2^{100} \bmod 13$.  By Fermat, $2^{12} \equiv 1 \pmod{13}$.
Since $100 = 12 \cdot 8 + 4$, we get
$2^{100} = (2^{12})^8 \cdot 2^4 \equiv 1^8 \cdot 16 \equiv 3 \pmod{13}$.
\end{example}

\begin{remark}[Fermat pseudoprimes]
\label{rem:pseudoprimes}
\index{pseudoprime!Fermat}
The converse of Fermat's little theorem is false.  A composite number~$n$
satisfying $2^{n-1} \equiv 1 \pmod{n}$ is called a
\emph{Fermat pseudoprime}\index{Fermat pseudoprime} to base~$2$.
The smallest example is $n = 341 = 11 \cdot 31$.  Composite numbers that
satisfy $a^{n-1} \equiv 1 \pmod{n}$ for \emph{all} $a$ with $\gcd(a,n)=1$
are called \emph{Carmichael numbers}\index{Carmichael number}; the smallest
is~$561 = 3 \cdot 11 \cdot 17$.
\end{remark}

% ------------------------------------------------------------
\section{Wilson's Theorem}
\label{sec:wilson}
\index{Wilson's theorem}

\begin{theorem}[Wilson's theorem]
\label{thm:wilson}
An integer $p \ge 2$ is prime if and only if
\[
  (p-1)! \equiv -1 \pmod{p}.
\]
\end{theorem}

\begin{proof}
We prove both directions.

\medskip\noindent
\textbf{($\Rightarrow$) If $p$ is prime, then $(p-1)! \equiv -1 \pmod{p}$.}

Consider the elements $1, 2, \dots, p-1$ in $(\Z/p\Z)^\times$.  Since
$\Z/p\Z$ is a field, each element~$a$ has a unique multiplicative inverse
$a^{-1}$, and $a = a^{-1}$ if and only if $a^2 \equiv 1 \pmod{p}$, i.e.\
$p \mid (a-1)(a+1)$.  Since $p$ is prime, this forces $a \equiv 1$ or
$a \equiv -1 \pmod{p}$.

Thus, in the product $(p-1)! = 1 \cdot 2 \cdot 3 \cdots (p-1)$, the elements
$2, 3, \dots, p-2$ can be paired off: each~$a$ with its inverse $a^{-1} \neq a$.
Each such pair contributes $a \cdot a^{-1} = 1$ to the product.  The only
unpaired elements are $1$ and $p-1 \equiv -1$.  Therefore
\[
  (p-1)! \equiv 1 \cdot (-1) \cdot \prod_{\text{pairs}} 1 = -1 \pmod{p}.
\]

\medskip\noindent
\textbf{($\Leftarrow$) If $(p-1)! \equiv -1 \pmod{p}$, then $p$ is prime.}

Suppose $p$ is composite.  Then $p$ has a divisor~$d$ with $1 < d < p$,
so $d$ appears as one of the factors in $(p-1)!$, hence $d \mid (p-1)!$.
Since $d \mid p$, we would need $d \mid ((p-1)! + 1)$ and $d \mid (p-1)!$,
forcing $d \mid 1$, a contradiction.  Hence $p$ must be prime.
\end{proof}

\begin{example}
\label{ex:wilson}
For $p = 7$: $6! = 720 = 102 \cdot 7 + 6 = 103 \cdot 7 - 1$, confirming
$6! \equiv -1 \pmod{7}$.

The pairing: $2 \leftrightarrow 4$ (since $2 \cdot 4 = 8 \equiv 1$),
$3 \leftrightarrow 5$ (since $3 \cdot 5 = 15 \equiv 1$).  Unpaired: $1$ and $6$.
So $6! \equiv 1 \cdot (2 \cdot 4) \cdot (3 \cdot 5) \cdot 6 \equiv
1 \cdot 1 \cdot 1 \cdot (-1) = -1 \pmod{7}$.
\end{example}

\begin{corollary}
\label{cor:wilson-square}
\index{Wilson's theorem!quadratic variant}
If $p$ is an odd prime, then
\[
  \bigl(((p-1)/2)!\bigr)^2 \equiv (-1)^{(p+1)/2} \pmod{p}.
\]
\end{corollary}

\begin{proof}
Note that $k \equiv -(p-k) \pmod{p}$ for $1 \le k \le p-1$.  Thus
\begin{align*}
  (p-1)! &= \prod_{k=1}^{(p-1)/2} k \cdot \prod_{k=(p+1)/2}^{p-1} k
           = \prod_{k=1}^{(p-1)/2} k \cdot \prod_{j=1}^{(p-1)/2} (p - j) \\
         &\equiv \prod_{k=1}^{(p-1)/2} k \cdot \prod_{j=1}^{(p-1)/2} (-j)
          = (-1)^{(p-1)/2} \left(\left(\frac{p-1}{2}\right)!\right)^{\!2} \pmod{p}.
\end{align*}
By Wilson, the left side is $\equiv -1 \pmod{p}$, so
$\bigl(((p-1)/2)!\bigr)^2 \equiv (-1) \cdot (-1)^{-(p-1)/2}
= (-1)^{1-(p-1)/2} = (-1)^{(3-p)/2} = (-1)^{(p+1)/2} \pmod{p}$,
where the last equality uses $(-1)^{(3-p)/2} = (-1)^{(p+1)/2}$
since $(3-p)/2 + (p+1)/2 = 2$ is even.
\end{proof}

% ------------------------------------------------------------
\section{Primitive Roots}
\label{sec:primitive-roots}
\index{primitive root}

\begin{definition}[Order of an element]
\label{def:order-element}
\index{order!of an element modulo $n$}
Let $\gcd(a,n) = 1$.  The \emph{order} of~$a$ modulo~$n$, denoted
$\ord_n(a)$\index{ord@$\ord_n(a)$}, is the smallest positive integer~$d$
such that $a^d \equiv 1 \pmod{n}$.
\end{definition}

\begin{proposition}[Basic properties of the order]
\label{prop:order-props}
\index{order!properties}
Let $\gcd(a,n) = 1$ and $d = \ord_n(a)$.  Then:
\begin{enumerate}[label=(\roman*)]
  \item $a^k \equiv 1 \pmod{n}$ if and only if $d \mid k$.
  \item $d \mid \varphi(n)$.
  \item $\ord_n(a^j) = d / \gcd(j, d)$.
\end{enumerate}
\end{proposition}

\begin{proof}
\leavevmode
\begin{enumerate}[label=(\roman*)]
  \item ($\Leftarrow$) If $k = dq$, then $a^k = (a^d)^q \equiv 1$.
        ($\Rightarrow$) Write $k = dq + r$ with $0 \le r < d$.  Then
        $a^r = a^k \cdot (a^d)^{-q} \equiv 1 \pmod{n}$.
        By minimality of~$d$, $r = 0$.
  \item By Euler's theorem, $a^{\varphi(n)} \equiv 1$, so $d \mid \varphi(n)$
        by~(i).
  \item Let $e = \ord_n(a^j)$.  Then $(a^j)^e = a^{je} \equiv 1$, so
        $d \mid je$, i.e.\ $(d/g) \mid e$ where $g = \gcd(j,d)$.
        Conversely, $(a^j)^{d/g} = (a^d)^{j/g} \equiv 1$, so $e \mid (d/g)$.
        Hence $e = d/g$.
        \qedhere
\end{enumerate}
\end{proof}

\begin{definition}[Primitive root]
\label{def:primitive-root}
\index{primitive root!definition}
A \emph{primitive root modulo~$n$} is an integer~$g$ with $\gcd(g,n)=1$
and $\ord_n(g) = \varphi(n)$.  Equivalently, $g$ is a primitive root
iff $\bar{g}$ generates the cyclic group $(\Z/n\Z)^\times$.
\end{definition}

\begin{theorem}[Existence of primitive roots for primes]
\label{thm:primitive-root-prime}
\index{primitive root!existence for primes}
\index{cyclic group}
For every prime~$p$, the group $(\Z/p\Z)^\times$ is cyclic of order~$p-1$.
In other words, there exists a primitive root modulo~$p$.
\end{theorem}

\begin{proof}
We use a counting argument.  For each divisor~$d$ of $p-1$, let
$\psi(d)$ denote the number of elements of order~$d$ in $(\Z/p\Z)^\times$.

\medskip\noindent
\textbf{Step 1: The polynomial $x^d - 1$ has at most $d$ roots in $\Z/p\Z$.}

Since $\Z/p\Z$ is a field, a polynomial of degree~$d$ over $\Z/p\Z$ has
at most $d$ roots.

\medskip\noindent
\textbf{Step 2: For each $d \mid (p-1)$, $x^d \equiv 1 \pmod{p}$ has
exactly $d$ solutions.}

Since $d \mid (p-1)$, we can write $x^{p-1} - 1 = (x^d - 1)\cdot q(x)$
for some polynomial~$q$ of degree $p-1-d$.  The polynomial $x^{p-1} - 1$
has exactly $p-1$ roots (namely $1, 2, \dots, p-1$, by Fermat), and $q(x)$
has at most $p-1-d$ roots.  Therefore $x^d - 1$ must have at least
$(p-1) - (p-1-d) = d$ roots, and by Step~1, at most~$d$.  So it has exactly~$d$
roots.

\medskip\noindent
\textbf{Step 3: If $\psi(d) > 0$ for some $d \mid (p-1)$, then $\psi(d) = \varphi(d)$.}

Suppose $a$ has order~$d$.  Then $a, a^2, \dots, a^d$ are $d$ distinct
elements satisfying $x^d \equiv 1$.  By Step~2 these are \emph{all}
solutions of $x^d \equiv 1$.  Among them, $a^j$ has order~$d$ iff
$\gcd(j,d) = 1$ (by Proposition~\ref{prop:order-props}(iii)), and there
are exactly $\varphi(d)$ such values of~$j$.  Hence $\psi(d) = \varphi(d)$.

\medskip\noindent
\textbf{Step 4: $\psi(d) = \varphi(d)$ for all $d \mid (p-1)$.}

Every element of $(\Z/p\Z)^\times$ has some order~$d$ dividing $p-1$, so
\[
  p - 1 = \sum_{d \mid (p-1)} \psi(d).
\]
From Gauss's identity (Theorem~\ref{thm:gauss-divisor-sum}),
$p - 1 = \sum_{d \mid (p-1)} \varphi(d)$.  Since $\psi(d) \le \varphi(d)$
for all~$d$ (by Step~3 and the fact that $\psi(d) \in \{0, \varphi(d)\}$),
equality of the sums forces $\psi(d) = \varphi(d)$ for every~$d$.

\medskip\noindent
In particular, $\psi(p-1) = \varphi(p-1) \ge 1$, so there exists an
element of order~$p-1$, which is a primitive root.
\end{proof}

\begin{example}[Primitive roots modulo $7$]
\label{ex:prim-root-7}
\index{primitive root!modulo 7}
The elements of $(\Z/7\Z)^\times = \{1,2,3,4,5,6\}$ have orders:
\[
\begin{array}{c|cccccc}
  a & 1 & 2 & 3 & 4 & 5 & 6 \\ \hline
  \ord_7(a) & 1 & 3 & 6 & 3 & 6 & 2
\end{array}
\]
The primitive roots are $3$ and $5$ (each has order $\varphi(7) = 6$).
Note $\varphi(6) = 2$, matching the count.
\end{example}

% --- TikZ: powers of a primitive root ---
\begin{figure}[ht]
\centering
\begin{tikzpicture}[scale=1.2]
  % Cyclic group (Z/7Z)* generated by g=3
  \def\n{6}
  \foreach \i [evaluate=\i as \angle using {90 - (\i-1)*360/\n}] in {1,...,6} {
    % Compute 3^i mod 7
    \pgfmathtruncatemacro{\val}{
      \i == 1 ? 3 :
      \i == 2 ? 2 :
      \i == 3 ? 6 :
      \i == 4 ? 4 :
      \i == 5 ? 5 :
      1
    }
    \node[circle, draw, fill=blue!15, inner sep=3pt, minimum size=26pt,
          font=\small\bfseries] (v\i) at (\angle:2.2) {$\val$};
    \node[font=\tiny, text=red!70!black] at (\angle:1.55) {$3^{\i}$};
  }
  % Arrows showing cyclic structure
  \foreach \i [evaluate=\i as \nexti using {int(mod(\i,\n)+1)}] in {1,...,6} {
    \draw[-{Stealth}, thick, blue!60!black]
      (v\i) -- (v\nexti);
  }
  \node[below, font=\small] at (0,-2.8)
    {The cyclic group $(\Z/7\Z)^\times = \langle 3 \rangle$: successive powers of $g=3$.};
\end{tikzpicture}
\caption{The cyclic structure of $(\Z/7\Z)^\times$ generated by the primitive root $g = 3$.}
\label{fig:cyclic-Z7}
\index{cyclic group!visualization}
\end{figure}

\begin{theorem}[Characterisation of moduli with primitive roots]
\label{thm:prim-root-exist}
\index{primitive root!existence characterisation}
Primitive roots exist modulo~$n$ if and only if
$n \in \{1, 2, 4, p^k, 2p^k\}$, where $p$ is an odd prime and $k \ge 1$.
\end{theorem}

\begin{remark}
\label{rem:prim-root-exist-proof}
The proof of this complete characterisation requires tools beyond the
scope of this chapter (specifically, the structure of the groups
$(\Z/2^k\Z)^\times$ for $k \ge 3$ and the Chinese Remainder Theorem
applied to unit groups).  We have proved the case $n = p$ in
Theorem~\ref{thm:primitive-root-prime}.
\end{remark}

\begin{proposition}[Number of primitive roots]
\label{prop:count-prim-roots}
\index{primitive root!count}
If primitive roots exist modulo~$n$, then there are exactly
$\varphi(\varphi(n))$ of them.
\end{proposition}

\begin{proof}
Let $g$ be a primitive root.  Then every unit modulo~$n$ is $g^j$ for a
unique $j \in \{0, 1, \dots, \varphi(n)-1\}$.  By
Proposition~\ref{prop:order-props}(iii), $g^j$ is a primitive root iff
$\gcd(j, \varphi(n)) = 1$, and there are $\varphi(\varphi(n))$ such values.
\end{proof}

% --- TikZ: power table ---
\begin{figure}[ht]
\centering
\begin{tikzpicture}[scale=0.68]
  % Power table for Z/11Z
  \def\p{11}
  % Header
  \node[font=\bfseries\small] at (0, 0.8) {$a \backslash k$};
  \foreach \k in {1,...,10} {
    \node[font=\bfseries\small] at (\k*1.1, 0.8) {$\k$};
  }
  \draw[thick] (-0.6, 0.45) -- (11.7, 0.45);
  % For each base a, row of powers a^k mod 11
  % a=2 (primitive root, order 10)
  \node[font=\bfseries\small] at (0, 0) {$2$};
  \foreach \k [evaluate=\k as \val using {int(mod(
    \k==1?2:\k==2?4:\k==3?8:\k==4?5:\k==5?10:\k==6?9:\k==7?7:\k==8?3:\k==9?6:1
  ,11))}] in {1,...,10}{
    \ifnum\val=1
      \node[fill=green!25, rounded corners=2pt, inner sep=2pt, font=\small]
        at (\k*1.1, 0) {$\val$};
    \else
      \node[font=\small] at (\k*1.1, 0) {$\val$};
    \fi
  }
  % a=3 (order 5)
  \node[font=\bfseries\small] at (0, -0.7) {$3$};
  \foreach \k [evaluate=\k as \val using {int(
    \k==1?3:\k==2?9:\k==3?5:\k==4?4:\k==5?1:\k==6?3:\k==7?9:\k==8?5:\k==9?4:1
  )}] in {1,...,10}{
    \ifnum\val=1
      \node[fill=green!25, rounded corners=2pt, inner sep=2pt, font=\small]
        at (\k*1.1, -0.7) {$\val$};
    \else
      \node[font=\small] at (\k*1.1, -0.7) {$\val$};
    \fi
  }
  % a=5 (order 5)
  \node[font=\bfseries\small] at (0, -1.4) {$5$};
  \foreach \k [evaluate=\k as \val using {int(
    \k==1?5:\k==2?3:\k==3?4:\k==4?9:\k==5?1:\k==6?5:\k==7?3:\k==8?4:\k==9?9:1
  )}] in {1,...,10}{
    \ifnum\val=1
      \node[fill=green!25, rounded corners=2pt, inner sep=2pt, font=\small]
        at (\k*1.1, -1.4) {$\val$};
    \else
      \node[font=\small] at (\k*1.1, -1.4) {$\val$};
    \fi
  }
  % a=10 (order 2)
  \node[font=\bfseries\small] at (0, -2.1) {$10$};
  \foreach \k [evaluate=\k as \val using {int(
    \k==1?10:\k==2?1:\k==3?10:\k==4?1:\k==5?10:\k==6?1:\k==7?10:\k==8?1:\k==9?10:1
  )}] in {1,...,10}{
    \ifnum\val=1
      \node[fill=green!25, rounded corners=2pt, inner sep=2pt, font=\small]
        at (\k*1.1, -2.1) {$\val$};
    \else
      \node[font=\small] at (\k*1.1, -2.1) {$\val$};
    \fi
  }
  \draw[thick] (0.5, 1.1) -- (0.5, -2.6);
  \node[below, font=\small] at (5.5, -2.9)
    {Powers $a^k \bmod 11$ for selected bases; entries equal to $1$ are highlighted.};
\end{tikzpicture}
\caption{Power table modulo $11$: $a = 2$ is a primitive root (order~$10$),
         while $a = 3,5$ have order~$5$ and $a = 10$ has order~$2$.}
\label{fig:power-table-11}
\index{power table}
\end{figure}

% ------------------------------------------------------------
\section{The RSA Cryptosystem}
\label{sec:rsa}
\index{RSA cryptosystem}
\index{cryptography!RSA}

The RSA cryptosystem, invented by Rivest, Shamir, and Adleman\index{Rivest, Ron}%
\index{Shamir, Adi}\index{Adleman, Leonard} in~1977, is perhaps the most
celebrated application of number theory.  Its security rests on the
difficulty of factoring large integers.

\subsection{Key Generation}
\label{subsec:rsa-keygen}
\index{RSA cryptosystem!key generation}

\begin{enumerate}
  \item Choose two distinct large primes $p$ and $q$.  Set $n = pq$.
  \item Compute $\varphi(n) = (p-1)(q-1)$.
  \item Choose an integer~$e$ with $1 < e < \varphi(n)$ and $\gcd(e, \varphi(n)) = 1$.
        (The common choice is $e = 65537 = 2^{16} + 1$.)
  \item Compute $d \equiv e^{-1} \pmod{\varphi(n)}$, i.e.\ find $d$ with
        $ed \equiv 1 \pmod{\varphi(n)}$.
  \item The \emph{public key}\index{public key} is $(n, e)$; the
        \emph{private key}\index{private key} is $(n, d)$.
\end{enumerate}

\subsection{Encryption and Decryption}
\label{subsec:rsa-enc-dec}
\index{RSA cryptosystem!encryption}\index{RSA cryptosystem!decryption}

To encrypt a message $m$ (an integer with $0 \le m < n$):
\[
  c \equiv m^e \pmod{n} \qquad \text{(ciphertext)}.
\]
To decrypt:
\[
  m \equiv c^d \pmod{n} \qquad \text{(recover plaintext)}.
\]

\subsection{Proof of Correctness}
\label{subsec:rsa-proof}
\index{RSA cryptosystem!correctness proof}

\begin{theorem}[RSA correctness]
\label{thm:rsa-correct}
With notation as above, $c^d \equiv m \pmod{n}$ for every integer $m$
with $0 \le m < n$.
\end{theorem}

\begin{proof}
Since $ed \equiv 1 \pmod{\varphi(n)}$, write $ed = 1 + k\varphi(n)$
for some non-negative integer~$k$.  We must show $m^{ed} \equiv m \pmod{n}$,
i.e.\ $m^{1 + k(p-1)(q-1)} \equiv m \pmod{pq}$.

By the Chinese Remainder Theorem, it suffices to prove the congruence
modulo~$p$ and modulo~$q$ separately.

\medskip\noindent
\textbf{Modulo $p$:}
If $p \mid m$, then $m \equiv 0 \pmod{p}$ and $m^{ed} \equiv 0 \equiv m$.
If $p \nmid m$, then by Fermat's little theorem $m^{p-1} \equiv 1 \pmod{p}$, so
\[
  m^{ed} = m^{1 + k(p-1)(q-1)} = m \cdot \bigl(m^{p-1}\bigr)^{k(q-1)}
  \equiv m \cdot 1^{k(q-1)} = m \pmod{p}.
\]

\medskip\noindent
\textbf{Modulo $q$:} By the identical argument, $m^{ed} \equiv m \pmod{q}$.

\medskip\noindent
Since $\gcd(p,q) = 1$ and $m^{ed} \equiv m$ both mod~$p$ and mod~$q$,
we conclude $m^{ed} \equiv m \pmod{pq} = m \pmod{n}$.
\end{proof}

\begin{example}[Toy RSA]
\label{ex:toy-rsa}
\index{RSA cryptosystem!example}
Let $p = 11$, $q = 13$, so $n = 143$ and $\varphi(n) = 10 \cdot 12 = 120$.
Choose $e = 7$.  Then $d \equiv 7^{-1} \pmod{120}$.  Since
$7 \cdot 103 = 721 = 6 \cdot 120 + 1$, we get $d = 103$.

Encrypt $m = 9$: $c \equiv 9^7 \pmod{143}$.  We compute:
$9^2 = 81$, $9^4 \equiv 81^2 = 6561 \equiv 6561 - 45 \cdot 143 = 6561 - 6435 = 126$,
$9^7 = 9^4 \cdot 9^2 \cdot 9 \equiv 126 \cdot 81 \cdot 9 \pmod{143}$.
Now $126 \cdot 81 = 10206 \equiv 10206 - 71 \cdot 143 = 10206 - 10153 = 53$,
then $53 \cdot 9 = 477 \equiv 477 - 3\cdot 143 = 477 - 429 = 48$.
So $c = 48$.

Decrypt: $m \equiv 48^{103} \pmod{143}$. Using repeated squaring (which
we omit for brevity), one verifies $48^{103} \equiv 9 \pmod{143}$, recovering
the original message.
\end{example}

\begin{remark}[Security of RSA]
\label{rem:rsa-security}
\index{RSA cryptosystem!security}
An adversary who knows $n$ and $e$ can recover $d$ if they can compute
$\varphi(n) = (p-1)(q-1)$.  This is equivalent in difficulty to factoring
$n = pq$, which---for $n$ of several thousand bits---is believed to be
computationally infeasible with current algorithms.  In practice, one
uses primes of at least $1024$ bits each, giving $n$ of at least $2048$ bits.
\end{remark}

% ------------------------------------------------------------
\section{Diffie--Hellman Key Exchange}
\label{sec:diffie-hellman}
\index{Diffie--Hellman key exchange}
\index{cryptography!Diffie--Hellman}

Published in 1976 by Whitfield Diffie\index{Diffie, Whitfield} and Martin
Hellman\index{Hellman, Martin}, the Diffie--Hellman protocol allows two
parties---traditionally called Alice\index{Alice} and Bob\index{Bob}---to
establish a shared secret over an insecure channel.

\subsection{Protocol Description}
\label{subsec:dh-protocol}

\begin{enumerate}
  \item \textbf{Public parameters:} A large prime~$p$ and a primitive root
        $g$ modulo~$p$ are agreed upon publicly.
  \item \textbf{Alice's step:} Alice chooses a secret integer $a$
        ($1 \le a \le p-2$) and sends $A \equiv g^a \pmod{p}$ to Bob.
  \item \textbf{Bob's step:} Bob chooses a secret integer $b$
        ($1 \le b \le p-2$) and sends $B \equiv g^b \pmod{p}$ to Alice.
  \item \textbf{Shared secret:} Alice computes $s \equiv B^a \equiv g^{ab}
        \pmod{p}$.  Bob computes $s \equiv A^b \equiv g^{ab} \pmod{p}$.
        Both arrive at the same shared secret~$s$.
\end{enumerate}

\begin{remark}[Security]
\label{rem:dh-security}
An eavesdropper who observes $p$, $g$, $A = g^a \bmod p$, and $B = g^b \bmod p$
must compute $g^{ab} \bmod p$.  This is the
\emph{Diffie--Hellman problem}\index{Diffie--Hellman problem},
which is believed to be as hard as the discrete logarithm problem---computing
$a$ from $g^a \bmod p$---for which no efficient classical algorithm is known.
\end{remark}

\begin{example}[Toy Diffie--Hellman]
\label{ex:toy-dh}
\index{Diffie--Hellman key exchange!example}
Let $p = 23$ and $g = 5$ (a primitive root modulo~$23$).
\begin{itemize}
  \item Alice picks $a = 6$ and sends $A \equiv 5^6 \equiv 15625 \equiv
        15625 - 679 \cdot 23 = 15625 - 15617 = 8 \pmod{23}$.
  \item Bob picks $b = 15$ and sends $B \equiv 5^{15} \pmod{23}$.
        Using repeated squaring: $5^2 = 25 \equiv 2$, $5^4 \equiv 4$,
        $5^8 \equiv 16$, $5^{15} = 5^8 \cdot 5^4 \cdot 5^2 \cdot 5
        \equiv 16 \cdot 4 \cdot 2 \cdot 5 = 640 \equiv 640 - 27\cdot 23
        = 640 - 621 = 19 \pmod{23}$.
  \item Alice computes $s \equiv B^a = 19^6 \pmod{23}$.
        $19^2 = 361 \equiv 361 - 15\cdot 23 = 16$, $19^4 \equiv 16^2 = 256
        \equiv 256 - 11\cdot 23 = 3$, $19^6 = 19^4 \cdot 19^2 \equiv
        3 \cdot 16 = 48 \equiv 2 \pmod{23}$.
  \item Bob computes $s \equiv A^b = 8^{15} \pmod{23}$.
        $8^2 = 64 \equiv 18$, $8^4 \equiv 324 \equiv 324 - 14\cdot 23 = 2$,
        $8^8 \equiv 4$, $8^{15} = 8^8 \cdot 8^4 \cdot 8^2 \cdot 8
        \equiv 4 \cdot 2 \cdot 18 \cdot 8 = 1152 \equiv 1152 - 50\cdot 23
        = 1152 - 1150 = 2 \pmod{23}$.
  \item Both obtain the shared secret $s = 2$.
\end{itemize}
\end{example}

% --- TikZ: Diffie-Hellman diagram ---
\begin{figure}[ht]
\centering
\begin{tikzpicture}[>=Stealth, node distance=4.5cm]
  % Alice and Bob
  \node[draw, rounded corners, fill=blue!10, minimum width=2.2cm,
        minimum height=1.2cm, font=\bfseries] (alice) {Alice};
  \node[draw, rounded corners, fill=red!10, minimum width=2.2cm,
        minimum height=1.2cm, font=\bfseries, right=of alice] (bob) {Bob};

  % Secret values
  \node[below=0.4cm of alice, font=\small, text=blue!70!black] (sa)
    {secret $a$};
  \node[below=0.4cm of bob, font=\small, text=red!70!black] (sb)
    {secret $b$};

  % Exchanges
  \draw[->, thick, blue!60!black, transform canvas={yshift=6pt}]
    (alice) -- node[above, font=\small] {$A = g^a \bmod p$} (bob);
  \draw[->, thick, red!60!black, transform canvas={yshift=-6pt}]
    (bob) -- node[below, font=\small] {$B = g^b \bmod p$} (alice);

  % Shared secret
  \node[below=1.6cm of $(alice)!0.5!(bob)$, draw, rounded corners,
        fill=green!15, font=\small\bfseries, minimum width=3.5cm] (secret)
    {Shared secret: $s = g^{ab} \bmod p$};
  \draw[->, thick, dashed, blue!60!black] (sa.south) |- (secret.west);
  \draw[->, thick, dashed, red!60!black]  (sb.south) |- (secret.east);
\end{tikzpicture}
\caption{The Diffie--Hellman key exchange protocol.}
\label{fig:diffie-hellman}
\end{figure}

% ------------------------------------------------------------
\section{Exercises}
\label{sec:ch4-exercises}

\begin{exercise}
\label{exer:totient-compute}
Compute $\varphi(n)$ for $n = 100, 252, 1000, 2520$.
\end{exercise}

\begin{exercise}
\label{exer:euler-app}
Use Euler's theorem to compute $3^{1000} \bmod 35$.
\end{exercise}

\begin{exercise}
\label{exer:fermat-app}
Find the remainder when $2^{2025}$ is divided by~$17$.
\end{exercise}

\begin{exercise}
\label{exer:wilson-app}
Using Wilson's theorem, find the least non-negative residue of
$18! \pmod{23}$.
\end{exercise}

\begin{exercise}
\label{exer:prim-root-11}
\index{primitive root!modulo 11}
Find all primitive roots modulo~$11$.  Verify that there are
$\varphi(\varphi(11)) = \varphi(10) = 4$ of them.
\end{exercise}

\begin{exercise}
\label{exer:prim-root-existence}
Prove that $(\Z/8\Z)^\times$ is not cyclic.  (Hence there is no primitive
root modulo~$8$.)
\end{exercise}

\begin{exercise}
\label{exer:order-properties}
Let $p$ be prime and $\gcd(a,p) = 1$.  Prove that if $\ord_p(a) = d$, then
$\{a^0, a^1, \dots, a^{d-1}\}$ are the $d$ distinct roots of $x^d - 1
\equiv 0 \pmod{p}$.
\end{exercise}

\begin{exercise}
\label{exer:rsa-1}
\index{RSA cryptosystem!exercise}
In a toy RSA system with $p = 7$, $q = 11$, $e = 13$:
\begin{enumerate}[label=(\alph*)]
  \item Compute $n$, $\varphi(n)$, and $d$.
  \item Encrypt $m = 5$ and then decrypt to verify correctness.
\end{enumerate}
\end{exercise}

\begin{exercise}
\label{exer:dh-1}
\index{Diffie--Hellman key exchange!exercise}
In a toy Diffie--Hellman exchange with $p = 29$, $g = 2$, Alice picks
$a = 11$ and Bob picks $b = 19$.  Compute the values exchanged and the
shared secret.
\end{exercise}

\begin{exercise}
\label{exer:carmichael}
\index{Carmichael number}
Verify that $561 = 3 \cdot 11 \cdot 17$ is a Carmichael number by
checking that $a^{560} \equiv 1 \pmod{561}$ for all $a$ with
$\gcd(a, 561) = 1$.
\textit{Hint:} Use Korselt's criterion: $n$ is a Carmichael number iff
$n$ is square-free and $(p-1) \mid (n-1)$ for every prime $p \mid n$.
\end{exercise}

\begin{exercise}
\label{exer:sum-totient}
Prove that $\sum_{d \mid n} \varphi(d) = n$ using Möbius
inversion\index{Möbius inversion}, providing an alternative to the
direct counting argument of Theorem~\ref{thm:gauss-divisor-sum}.
\end{exercise}

\begin{exercise}[Challenge]
\label{exer:prim-root-product}
\index{primitive root!product of all}
Let $p$ be an odd prime.  Prove that the product of all primitive roots
modulo~$p$ is congruent to~$1$ modulo~$p$.
\textit{Hint:} If $g$ is a primitive root, the primitive roots are $g^k$
with $\gcd(k, p-1) = 1$.  Sum the exponents and use the fact that
$\sum_{\substack{1 \le k \le n \\ \gcd(k,n)=1}} k = \frac{n \varphi(n)}{2}$
for $n > 2$.
\end{exercise}

\begin{exercise}[Challenge]
\label{exer:order-lcm}
\index{order!of a product}
Let $\gcd(a,n) = \gcd(b,n) = 1$ and suppose $\gcd(\ord_n(a), \ord_n(b)) = 1$.
Prove that $\ord_n(ab) = \ord_n(a) \cdot \ord_n(b)$.  Give a counterexample
when the orders are not coprime.
\end{exercise}

% ------------------------------------------------------------
\section{Chapter Summary}
\label{sec:ch4-summary}

\begin{enumerate}
  \item \textbf{Euler's totient function} $\varphi(n)$ counts integers
        $1 \le a \le n$ with $\gcd(a,n) = 1$.  It is multiplicative and
        satisfies $\varphi(n) = n \prod_{p \mid n}(1 - 1/p)$.
  \item \textbf{Euler's theorem:} $a^{\varphi(n)} \equiv 1 \pmod{n}$
        whenever $\gcd(a,n) = 1$.
  \item \textbf{Fermat's little theorem:} $a^{p-1} \equiv 1 \pmod{p}$ for
        prime~$p$ and $p \nmid a$.  Equivalently, $a^p \equiv a \pmod{p}$
        for all~$a$.
  \item \textbf{Wilson's theorem:} $(p-1)! \equiv -1 \pmod{p}$ if and only
        if $p$ is prime.
  \item \textbf{Primitive roots:} $(\Z/p\Z)^\times$ is cyclic for every
        prime~$p$.  More generally, primitive roots exist for $n = 1, 2, 4,
        p^k, 2p^k$.
  \item \textbf{RSA:} Encryption $c = m^e \bmod n$, decryption $m = c^d
        \bmod n$, where $ed \equiv 1 \pmod{\varphi(n)}$ and $n = pq$.
        Correctness follows from Fermat/Euler.
  \item \textbf{Diffie--Hellman:} Two parties exchange $g^a$ and $g^b$
        modulo a prime, arriving at a shared secret $g^{ab}$.  Security
        relies on the hardness of the discrete logarithm problem.
\end{enumerate}

% === END OF CHUNK ch3-4 ===
% === START OF CHUNK ch5-6 ===

% ============================================================
%  CHAPTER 5 — Chinese Remainder Theorem and Applications
% ============================================================
\chapter{Chinese Remainder Theorem and Applications}
\label{ch:crt}

\section*{Historical Introduction}
\addcontentsline{toc}{section}{Historical Introduction}

\index{Chinese Remainder Theorem!history}
The problem of solving simultaneous congruences appears in the \emph{Sunzi Suanjing}
(\emph{Master Sun's Mathematical Manual}), a Chinese text from the 3rd--5th century CE.
The classic formulation asks:

\begin{quote}
\itshape
There are certain things whose number is unknown. If we count them by threes, the
remainder is~$2$; if we count them by fives, the remainder is~$3$; if we count them by
sevens, the remainder is~$2$. How many things are there?
\end{quote}

\index{Sun Tzu}
Sun Tzu\footnote{Not to be confused with the military strategist of the same name.}
gave the answer~$23$ and a method of solution.
Centuries later, \textbf{Qin Jiushao}\index{Qin Jiushao} (1202--1261) developed
a general algorithm in his \emph{Mathematical Treatise in Nine Sections}~(1247),
which systematically solved systems of simultaneous congruences using what we would now
call the extended Euclidean algorithm. His work preceded the European rediscovery
by several hundred years.

In the West, the theorem was stated by Euler\index{Euler, Leonhard} and
given its modern algebraic formulation by Gauss\index{Gauss, Carl Friedrich}
in the \emph{Disquisitiones Arithmeticae}~(1801).

% ------------------------------------------------------------
\section{Systems of Linear Congruences}
\label{sec:crt-systems}

We begin with the motivating problem: given congruences
\[
  x \equiv a_1 \pmod{m_1}, \quad
  x \equiv a_2 \pmod{m_2}, \quad \ldots, \quad
  x \equiv a_k \pmod{m_k},
\]
when does a simultaneous solution exist, and when is it unique?

\begin{example}[Sun Tzu's problem]
\label{ex:sun-tzu}
\index{Sun Tzu!problem}
We seek $x$ satisfying
\[
  x \equiv 2 \pmod{3}, \quad
  x \equiv 3 \pmod{5}, \quad
  x \equiv 2 \pmod{7}.
\]
By inspection, $x = 23$ works. We shall soon see that this is the unique solution
modulo $3 \cdot 5 \cdot 7 = 105$, so the full solution set is $\{23 + 105k : k \in \Z\}$.
\end{example}

\begin{lemma}[Two-congruence case]
\label{lem:two-congruences}
\index{congruence!system of two}
Let $m_1, m_2 \in \Z_{>0}$. The system
\[
  x \equiv a_1 \pmod{m_1}, \quad x \equiv a_2 \pmod{m_2}
\]
has a solution if and only if $\gcd(m_1, m_2) \mid (a_1 - a_2)$.
When a solution exists, it is unique modulo $\lcm(m_1, m_2)$.
\end{lemma}

\begin{proof}
The system is equivalent to the existence of integers $k_1, k_2$ with
$a_1 + m_1 k_1 = a_2 + m_2 k_2$, i.e., $m_1 k_1 - m_2 k_2 = a_2 - a_1$.
By the characterisation of the image of the linear map
$(k_1, k_2) \mapsto m_1 k_1 - m_2 k_2$, this equation is solvable if and only if
$\gcd(m_1, m_2) \mid (a_2 - a_1)$.

For uniqueness, suppose $x_0$ and $x_0'$ are both solutions. Then
$m_1 \mid (x_0' - x_0)$ and $m_2 \mid (x_0' - x_0)$, whence
$\lcm(m_1, m_2) \mid (x_0' - x_0)$.
\end{proof}

% ------------------------------------------------------------
\section{The Chinese Remainder Theorem}
\label{sec:crt-statement}

\begin{theorem}[Chinese Remainder Theorem --- CRT]
\label{thm:crt}
\index{Chinese Remainder Theorem}
Let $m_1, m_2, \ldots, m_k$ be pairwise coprime positive integers, and set
$M = m_1 m_2 \cdots m_k$. Then for any integers $a_1, a_2, \ldots, a_k$, the system
\begin{equation}
\label{eq:crt-system}
  x \equiv a_i \pmod{m_i}, \qquad i = 1, 2, \ldots, k,
\end{equation}
has a solution, and any two solutions are congruent modulo~$M$.
Explicitly, the unique solution modulo~$M$ is
\begin{equation}
\label{eq:crt-formula}
  x \equiv \sum_{i=1}^{k} a_i \, M_i \, y_i \pmod{M},
\end{equation}
where $M_i = M / m_i$ and $y_i$ is any integer satisfying $M_i \, y_i \equiv 1 \pmod{m_i}$.
\end{theorem}

\begin{proof}
\index{Chinese Remainder Theorem!proof}
We give a constructive proof in two parts.

\medskip\noindent\textbf{Existence.}
For each $i \in \{1, \ldots, k\}$, define $M_i = M / m_i = \prod_{j \neq i} m_j$.
Since the $m_j$ are pairwise coprime, every prime factor of $m_i$ divides $m_i$
but none of the $m_j$ ($j \neq i$), so $\gcd(M_i, m_i) = 1$.
By B\'ezout's identity, there exists $y_i \in \Z$ with
\[
  M_i \, y_i \equiv 1 \pmod{m_i}.
\]
Now set
\[
  x_0 = \sum_{i=1}^{k} a_i \, M_i \, y_i.
\]
We verify that $x_0$ solves each congruence. Fix any index $j$. For $i \neq j$,
the factor $M_i$ contains $m_j$ as one of its factors (since $j \neq i$), so
$m_j \mid M_i$ and thus $a_i M_i y_i \equiv 0 \pmod{m_j}$.
For $i = j$, we have $a_j M_j y_j \equiv a_j \cdot 1 = a_j \pmod{m_j}$.
Therefore
\[
  x_0 \equiv a_j \pmod{m_j}
\]
for every $j$, as required.

\medskip\noindent\textbf{Uniqueness.}
Suppose $x_0$ and $x_0'$ both satisfy \eqref{eq:crt-system}.
Then $m_i \mid (x_0' - x_0)$ for all~$i$.
Since $m_1, \ldots, m_k$ are pairwise coprime, it follows that
$M = m_1 \cdots m_k$ divides $x_0' - x_0$
(this is a straightforward induction using the fact that if $\gcd(a, b) = 1$
and $a \mid n$, $b \mid n$, then $ab \mid n$).
Hence $x_0' \equiv x_0 \pmod{M}$.
\end{proof}

\begin{example}[Solving Sun Tzu's problem via CRT]
\label{ex:sun-tzu-crt}
With $m_1 = 3$, $m_2 = 5$, $m_3 = 7$, we have $M = 105$, and
\[
  M_1 = 35, \quad M_2 = 21, \quad M_3 = 15.
\]
We find the inverses modulo each $m_i$:
\begin{itemize}
  \item $35 y_1 \equiv 1 \pmod{3}$: since $35 \equiv 2 \pmod{3}$, we need $2y_1 \equiv 1 \pmod{3}$, giving $y_1 = 2$.
  \item $21 y_2 \equiv 1 \pmod{5}$: since $21 \equiv 1 \pmod{5}$, we get $y_2 = 1$.
  \item $15 y_3 \equiv 1 \pmod{7}$: since $15 \equiv 1 \pmod{7}$, we get $y_3 = 1$.
\end{itemize}
Therefore
\[
  x_0 = 2 \cdot 35 \cdot 2 + 3 \cdot 21 \cdot 1 + 2 \cdot 15 \cdot 1
       = 140 + 63 + 30 = 233 \equiv 23 \pmod{105}. \qedhere
\]
\end{example}

% ------------------------------------------------------------
\section{Algebraic Formulation}
\label{sec:crt-algebraic}

The CRT admits a clean algebraic reformulation as a ring isomorphism.
\index{Chinese Remainder Theorem!algebraic form}

\begin{theorem}[CRT --- algebraic form]
\label{thm:crt-algebraic}
Let $m, n$ be positive integers with $\gcd(m, n) = 1$. The map
\begin{equation}
\label{eq:crt-iso}
  \varphi \colon \Z / mn\Z \longrightarrow \Z / m\Z \times \Z / n\Z, \qquad
  x \bmod{mn} \longmapsto (x \bmod{m},\; x \bmod{n})
\end{equation}
is a ring isomorphism.
\end{theorem}

\begin{proof}
\index{ring isomorphism!CRT}
\emph{Well-definedness and homomorphism.} If $x \equiv x' \pmod{mn}$, then
$x \equiv x' \pmod{m}$ and $x \equiv x' \pmod{n}$, so $\varphi$ is well-defined.
It clearly preserves addition and multiplication (since reduction modulo $m$ or $n$
is a ring homomorphism) and sends $1$ to $(1, 1)$.

\emph{Injectivity.} If $\varphi(x) = (0, 0)$, then $m \mid x$ and $n \mid x$.
Since $\gcd(m, n) = 1$, we get $mn \mid x$, so $x \equiv 0 \pmod{mn}$.

\emph{Surjectivity.} Both sides are finite sets of the same cardinality
$mn = m \cdot n$, so injectivity implies surjectivity.
Alternatively, surjectivity is exactly the existence statement of
Theorem~\ref{thm:crt}.
\end{proof}

\begin{corollary}[Euler's totient is multiplicative]
\label{cor:euler-multiplicative}
\index{Euler's totient function!multiplicativity}
If $\gcd(m, n) = 1$, then $\varphi(mn) = \varphi(m)\,\varphi(n)$,
where $\varphi$ denotes Euler's totient function.
\end{corollary}

\begin{proof}
The ring isomorphism $\Z/mn\Z \cong \Z/m\Z \times \Z/n\Z$ restricts to a
group isomorphism on the unit groups:
$(\Z/mn\Z)^\times \cong (\Z/m\Z)^\times \times (\Z/n\Z)^\times$.
Taking cardinalities gives $\varphi(mn) = \varphi(m)\,\varphi(n)$.
\end{proof}

\begin{corollary}[General CRT]
\label{cor:crt-general}
If $m_1, \ldots, m_k$ are pairwise coprime, then
\[
  \Z / M\Z \;\cong\; \Z/m_1\Z \times \Z/m_2\Z \times \cdots \times \Z/m_k\Z
\]
as rings, where $M = m_1 m_2 \cdots m_k$.
\end{corollary}

\begin{proof}
Apply Theorem~\ref{thm:crt-algebraic} iteratively. At each step, the remaining
product is coprime to the factor being split off, since pairwise coprimality is
preserved.
\end{proof}

% --- TikZ: CRT isomorphism visualisation ---
\begin{figure}[ht]
\centering
\begin{tikzpicture}[
  box/.style={draw, rounded corners=4pt, minimum width=2.8cm, minimum height=1.2cm,
              font=\large, thick},
  arr/.style={-{Stealth[length=6pt]}, thick, shorten <=4pt, shorten >=4pt}
]
  % Left box
  \node[box, fill=blue!8, draw=blue!60!black] (A) {$\Z/mn\Z$};
  % Right boxes
  \node[box, fill=red!8, draw=red!60!black] (B) at (6,0.9) {$\Z/m\Z$};
  \node[box, fill=green!8, draw=green!60!black] (C) at (6,-0.9) {$\Z/n\Z$};
  % Product symbol
  \node at (6,0) {$\times$};
  % Arrow
  \draw[arr, blue!60!black] (A) -- node[above, font=\small] {$\varphi$}
        node[below, font=\small] {$\cong$} (3.8,0);
  % Projection arrows
  \draw[arr, red!60!black] (3.8,0) -- (B.west);
  \draw[arr, green!60!black] (3.8,0) -- (C.west);
  % Labels
  \node[above=0.4cm of A, font=\small\itshape, text=blue!60!black]
    {$\gcd(m,n)=1$};
  % Concrete example below
  \node[below=1.5cm of A, font=\small, text=gray!70!black] (ex1)
    {e.g.\ $\Z/15\Z$};
  \node[below=0.8cm of B, font=\small, text=gray!70!black] (ex2)
    {$\Z/3\Z$};
  \node[below=0.8cm of C, font=\small, text=gray!70!black] (ex3)
    {$\Z/5\Z$};
  \node at (6,-2.15) [font=\small, text=gray!70!black] {$\times$};
  \draw[-{Stealth}, gray, dashed, shorten <=2pt, shorten >=2pt]
    (ex1) -- (4.2,-2.15);
\end{tikzpicture}
\caption{The CRT isomorphism $\Z/mn\Z \cong \Z/m\Z \times \Z/n\Z$ for coprime $m,n$.}
\label{fig:crt-iso}
\end{figure}

% ------------------------------------------------------------
\section{Applications}
\label{sec:crt-applications}

\subsection{Solving simultaneous congruences}
\label{subsec:crt-solving}

The constructive proof of the CRT yields an explicit algorithm.

\begin{example}[A four-modulus system]
\label{ex:crt-four-moduli}
Solve the system
\[
  x \equiv 1 \pmod{2}, \quad
  x \equiv 2 \pmod{3}, \quad
  x \equiv 3 \pmod{5}, \quad
  x \equiv 4 \pmod{7}.
\]
Here $M = 210$, and
\begin{align*}
  M_1 &= 105,\; y_1: 105 y_1 \equiv 1 \pmod{2} \implies y_1 = 1,\\
  M_2 &= 70,\; y_2: 70 y_2 \equiv 1 \pmod{3} \implies y_2 = 1 \;\;(\text{since } 70 \equiv 1),\\
  M_3 &= 42,\; y_3: 42 y_3 \equiv 1 \pmod{5} \implies y_3 = 3 \;\;(\text{since } 42 \equiv 2,\; 2\cdot 3 = 6 \equiv 1),\\
  M_4 &= 30,\; y_4: 30 y_4 \equiv 1 \pmod{7} \implies y_4 = 4 \;\;(\text{since } 30 \equiv 2,\; 2\cdot 4 = 8 \equiv 1).
\end{align*}
Then
\[
  x_0 = 1 \cdot 105 \cdot 1 + 2 \cdot 70 \cdot 1 + 3 \cdot 42 \cdot 3 + 4 \cdot 30 \cdot 4
      = 105 + 140 + 378 + 480 = 1103 \equiv 53 \pmod{210}.
\]
One checks: $53 = 26 \cdot 2 + 1$, $53 = 17 \cdot 3 + 2$, $53 = 10 \cdot 5 + 3$, $53 = 7 \cdot 7 + 4$.
\end{example}

\subsection{Calendar computations}
\label{subsec:crt-calendar}
\index{Chinese Remainder Theorem!calendar application}

Many calendar systems involve independent cycles of different periods.
For instance, the \emph{Chinese sexagenary cycle}\index{sexagenary cycle}
combines a cycle of $10$ Heavenly Stems with a cycle of $12$ Earthly Branches.
Since $\gcd(10, 12) = 2 \neq 1$, the CRT does not directly apply to the full
pair---but the reduced coprime pair $(5, 12)$ governs the $60$-year cycle,
since $\lcm(10, 12) = 60 = 5 \cdot 12$.

\begin{example}[Day of the week]
\label{ex:crt-weekday}
Suppose an event recurs every $7$ days (weekly) and another every $4$ days.
Since $\gcd(7,4) = 1$, by the CRT the combined pattern repeats every $28$ days.
If the weekly event falls on day~$3$ and the $4$-day event on day~$1$, we solve
\[
  x \equiv 3 \pmod{7}, \qquad x \equiv 1 \pmod{4}.
\]
Using $M_1 = 4$, $y_1 = 2$ (since $4 \cdot 2 = 8 \equiv 1 \pmod{7}$) and
$M_2 = 7$, $y_2 = 3$ (since $7 \cdot 3 = 21 \equiv 1 \pmod{4}$):
\[
  x \equiv 3 \cdot 4 \cdot 2 + 1 \cdot 7 \cdot 3 = 24 + 21 = 45 \equiv 17 \pmod{28}.
\]
\end{example}

\subsection{Error detection and secret sharing}
\label{subsec:crt-codes}
\index{Chinese Remainder Theorem!error detection}
\index{secret sharing}

The CRT underlies several practical schemes:

\begin{enumerate}[label=(\roman*)]
  \item \textbf{Redundant residue number systems.}\index{residue number system}
    Represent an integer by its residues modulo several coprime moduli.
    Adding extra moduli provides error-detection capability, since any single
    residue error yields a tuple outside the image of the CRT isomorphism.

  \item \textbf{Shamir-type secret sharing.}\index{Shamir secret sharing}
    A variant due to Asmuth and Bloom\index{Asmuth--Bloom scheme} uses the CRT:
    choose pairwise coprime moduli $m_1 < m_2 < \cdots < m_n$ and a threshold~$k$.
    The secret $S$ is encoded modulo $m_0$ and shared as $S_i \equiv S \pmod{m_i}$.
    Any $k$ shares suffice to reconstruct $S$ via the CRT, while fewer than $k$
    leave $S$ undetermined.
\end{enumerate}

\subsection{RSA speed-up via CRT}
\label{subsec:crt-rsa}
\index{RSA!CRT optimisation}
\index{Chinese Remainder Theorem!RSA}

In the RSA\index{RSA} cryptosystem, decryption computes $c^d \bmod n$ where $n = pq$.
By the CRT isomorphism $\Z/n\Z \cong \Z/p\Z \times \Z/q\Z$, we may instead compute
\[
  m_p \equiv c^{d \bmod (p-1)} \pmod{p}, \qquad
  m_q \equiv c^{d \bmod (q-1)} \pmod{q},
\]
and then recombine using Garner's formula\index{Garner's formula}:
\[
  m \equiv m_q + q \bigl[q^{-1}(m_p - m_q) \bmod p\bigr] \pmod{n}.
\]
Since modular exponentiation is roughly cubic in the bit-length of the modulus,
working with $p$ and $q$ separately yields an approximate $4\times$ speed-up.

\begin{remark}
\label{rem:crt-garner}
Garner's formula is preferred over the symmetric CRT formula in practice because
it avoids reduction modulo the full product~$n$.
\end{remark}

% ------------------------------------------------------------
\section{Exercises}
\label{sec:crt-exercises}

\begin{exercise}
\label{exer:crt-basic}
Solve the system $x \equiv 3 \pmod{7}$, $x \equiv 5 \pmod{11}$, $x \equiv 8 \pmod{13}$.
\end{exercise}

\begin{exercise}
\label{exer:crt-non-coprime}
\index{congruence!non-coprime moduli}
Show that $x \equiv 3 \pmod{6}$, $x \equiv 5 \pmod{10}$ has a solution,
and find it. What is the modulus of the general solution?
\textup{(}Hint: $\gcd(6,10) = 2$ and $2 \mid (5-3)$.\textup{)}
\end{exercise}

\begin{exercise}
\label{exer:crt-units}
Use the CRT to show that $(\Z/n\Z)^\times$ is not cyclic when
$n$ has at least two distinct odd prime factors.
\textup{(}Hint: if $n = p_1^{a_1} \cdots p_k^{a_k}$, the group splits as a
direct product, each factor of even order.\textup{)}
\end{exercise}

\begin{exercise}
\label{exer:crt-polynomial}
Let $p, q$ be distinct primes. Show that the number of solutions of $x^2 \equiv 1 \pmod{pq}$
is exactly $4$.
\end{exercise}

\begin{exercise}
\label{exer:crt-simultaneous-prime}
Find the smallest positive integer satisfying
\[
  x \equiv 1 \pmod{2}, \quad x \equiv 1 \pmod{3}, \quad x \equiv 1 \pmod{5},
  \quad x \equiv 0 \pmod{7}.
\]
\end{exercise}

\begin{exercise}
\label{exer:crt-totient}
Using the CRT and the formula $\varphi(p^a) = p^{a-1}(p-1)$, derive the general
formula
\[
  \varphi(n) = n \prod_{p \mid n} \Bigl(1 - \frac{1}{p}\Bigr).
\]
\end{exercise}

\begin{exercise}
\label{exer:crt-idempotents}
\index{idempotent}
An element $e$ of a ring $R$ is called \emph{idempotent} if $e^2 = e$.
Show that $\Z/n\Z$ contains exactly $2^k$ idempotents, where $k$ is the
number of distinct prime divisors of $n$.
\textup{(}Hint: use the CRT and note that the only idempotents in $\Z/p^a\Z$
are $0$ and $1$.\textup{)}
\end{exercise}

\begin{exercise}[RSA speed-up]
\label{exer:crt-rsa}
Let $p = 61$, $q = 53$, $n = 3233$, $e = 17$, $d = 2753$.
Compute $c^d \bmod n$ for $c = 2790$ using the CRT method described in
Section~\ref{subsec:crt-rsa}. Verify that you recover $m = 65$.
\end{exercise}

% ------------------------------------------------------------
\section*{Chapter Summary}
\addcontentsline{toc}{section}{Chapter Summary}

\begin{itemize}
  \item The \textbf{Chinese Remainder Theorem} states that a system of congruences
    with pairwise coprime moduli $m_1, \ldots, m_k$ always has a unique solution
    modulo $M = m_1 \cdots m_k$.

  \item The explicit solution is
    $x \equiv \sum_{i=1}^k a_i M_i y_i \pmod{M}$, where $M_i = M/m_i$ and
    $M_i y_i \equiv 1 \pmod{m_i}$.

  \item Algebraically, $\Z/mn\Z \cong \Z/m\Z \times \Z/n\Z$ when $\gcd(m,n)=1$,
    which implies the multiplicativity of Euler's totient function.

  \item Applications range from ancient calendar problems to modern cryptography (RSA)
    and secret-sharing schemes.
\end{itemize}

\bigskip
\hrule
\bigskip

% ============================================================
%  CHAPTER 6 — Quadratic Residues and Quadratic Reciprocity
% ============================================================
\chapter{Quadratic Residues and Quadratic Reciprocity}
\label{ch:qr}

\section*{Historical Introduction}
\addcontentsline{toc}{section}{Historical Introduction}

\index{quadratic reciprocity!history}
The study of quadratic residues is one of the crown jewels of elementary number theory.
\textbf{Euler}\index{Euler, Leonhard} investigated which primes could be represented by
binary quadratic forms such as $x^2 + ny^2$, leading him to discover (around 1783)
the law of quadratic reciprocity empirically, though he could not prove it.

\textbf{Legendre}\index{Legendre, Adrien-Marie} (1785) stated the law precisely and
introduced the symbol $\legendre{a}{p}$ that bears his name, but his proof contained
a gap (it relied on Dirichlet's theorem on primes in arithmetic progressions, which
would not be proven for another fifty years).

It was \textbf{Gauss}\index{Gauss, Carl Friedrich} who, at the age of eighteen, gave
the first complete proof in 1796.
He called it the \emph{theorema aureum}\index{theorema aureum} --- the golden theorem ---
and over the course of his life produced no fewer than \textbf{eight} distinct proofs,
a testament to the theorem's depth and centrality.
Today, over 240 proofs are known.

% ------------------------------------------------------------
\section{Quadratic Residues}
\label{sec:qr-definition}

\begin{definition}[Quadratic residue]
\label{def:qr}
\index{quadratic residue}
Let $p$ be an odd prime and $a$ an integer with $p \nmid a$.
We say that $a$ is a \emph{quadratic residue modulo~$p$} (abbreviated QR) if the
congruence $x^2 \equiv a \pmod{p}$ has a solution. Otherwise, $a$ is a
\emph{quadratic non-residue} (abbreviated QNR).
\index{quadratic non-residue}
\end{definition}

\begin{proposition}[Counting quadratic residues]
\label{prop:count-qr}
\index{quadratic residue!counting}
Let $p$ be an odd prime. Among the residues $1, 2, \ldots, p-1$, exactly
$\frac{p-1}{2}$ are quadratic residues and $\frac{p-1}{2}$ are quadratic non-residues.
\end{proposition}

\begin{proof}
Consider the squaring map $\sigma \colon (\Z/p\Z)^\times \to (\Z/p\Z)^\times$
defined by $\sigma(x) = x^2$. We claim that for each QR~$a$, the equation
$x^2 \equiv a \pmod{p}$ has exactly two solutions $\pm x_0$.

Indeed, if $x_0^2 \equiv a$ then $(-x_0)^2 \equiv a$, and these are distinct since
$p$ is odd ($x_0 \not\equiv -x_0 \pmod{p}$ because $p \nmid 2x_0$).
Conversely, if $x^2 \equiv x_0^2$ then $p \mid (x - x_0)(x + x_0)$, so $x \equiv \pm x_0$.

Thus $\sigma$ is a $2$-to-$1$ map from $(\Z/p\Z)^\times$ (of order $p-1$) onto
the set of quadratic residues, which therefore has cardinality $\frac{p-1}{2}$.
\end{proof}

\begin{example}[Quadratic residues modulo $13$]
\label{ex:qr-mod-13}
We compute $x^2 \bmod 13$ for $x = 1, \ldots, 6$
(the other squares are the same by $(-x)^2 = x^2$):
\[
  1^2 = 1,\quad 2^2 = 4,\quad 3^2 = 9,\quad 4^2 = 3,\quad 5^2 = 12,\quad 6^2 = 10.
\]
So the QRs mod $13$ are $\{1, 3, 4, 9, 10, 12\}$ and the QNRs are $\{2, 5, 6, 7, 8, 11\}$.
\end{example}

% ------------------------------------------------------------
\section{Euler's Criterion}
\label{sec:euler-criterion}

\begin{theorem}[Euler's criterion]
\label{thm:euler-criterion}
\index{Euler's criterion}
Let $p$ be an odd prime and $a$ an integer with $p \nmid a$. Then
\begin{equation}
\label{eq:euler-criterion}
  a^{(p-1)/2} \equiv
  \begin{cases}
    \phantom{-}1 \pmod{p} & \text{if } a \text{ is a QR mod } p, \\
    -1 \pmod{p} & \text{if } a \text{ is a QNR mod } p.
  \end{cases}
\end{equation}
\end{theorem}

\begin{proof}
\index{Euler's criterion!proof}
By Fermat's little theorem, $a^{p-1} \equiv 1 \pmod{p}$, so
\[
  \bigl(a^{(p-1)/2}\bigr)^2 \equiv 1 \pmod{p}.
\]
Since $\Z/p\Z$ is a field, the polynomial $t^2 - 1$ has at most two roots, namely
$t \equiv \pm 1$. Thus $a^{(p-1)/2} \equiv \pm 1 \pmod{p}$.

\emph{Case 1: $a$ is a QR.} Write $a \equiv x_0^2$. Then
\[
  a^{(p-1)/2} \equiv x_0^{p-1} \equiv 1 \pmod{p}
\]
by Fermat's little theorem.

\emph{Case 2: $a$ is a QNR.} Since $a^{(p-1)/2} \equiv \pm 1$ and Case~1 accounts
for all $(p-1)/2$ quadratic residues as roots of $t^{(p-1)/2} - 1 \equiv 0$,
and this polynomial has at most $(p-1)/2$ roots, the remaining $(p-1)/2$ elements
(the QNRs) must all satisfy $a^{(p-1)/2} \equiv -1$.
\end{proof}

% ------------------------------------------------------------
\section{The Legendre Symbol}
\label{sec:legendre-symbol}

\begin{definition}[Legendre symbol]
\label{def:legendre}
\index{Legendre symbol}
For an odd prime $p$ and an integer $a$, the \emph{Legendre symbol} is
\[
  \legendre{a}{p} =
  \begin{cases}
    0  & \text{if } p \mid a, \\
    1  & \text{if } a \text{ is a QR mod } p, \\
    -1 & \text{if } a \text{ is a QNR mod } p.
  \end{cases}
\]
\end{definition}

By Euler's criterion, we have the useful identity
\begin{equation}
\label{eq:legendre-euler}
  \legendre{a}{p} \equiv a^{(p-1)/2} \pmod{p}.
\end{equation}

\begin{proposition}[Properties of the Legendre symbol]
\label{prop:legendre-properties}
\index{Legendre symbol!properties}
Let $p$ be an odd prime and $a, b \in \Z$. Then:
\begin{enumerate}[label=\textup{(\roman*)}]
  \item \label{it:leg-cong}
    If $a \equiv b \pmod{p}$, then $\legendre{a}{p} = \legendre{b}{p}$.
  \item \label{it:leg-mult}
    $\legendre{ab}{p} = \legendre{a}{p} \legendre{b}{p}$.
    \textup{(}Complete multiplicativity.\textup{)}
  \item \label{it:leg-square}
    $\legendre{a^2}{p} = 1$ whenever $p \nmid a$.
  \item \label{it:leg-minus-one}
    $\legendre{-1}{p} = (-1)^{(p-1)/2}$.
    In particular, $-1$ is a QR mod $p$ if and only if $p \equiv 1 \pmod{4}$.
\end{enumerate}
\end{proposition}

\begin{proof}
\ref{it:leg-cong} is immediate from the definition. For \ref{it:leg-mult},
use Euler's criterion:
\[
  \legendre{ab}{p} \equiv (ab)^{(p-1)/2} = a^{(p-1)/2} b^{(p-1)/2}
  \equiv \legendre{a}{p} \legendre{b}{p} \pmod{p}.
\]
Since both sides are $\pm 1$ or $0$ and $p \geq 3$, congruence modulo $p$
implies equality.
Part \ref{it:leg-square} follows from \ref{it:leg-mult}, and
\ref{it:leg-minus-one} is Euler's criterion applied to $a = -1$.
\end{proof}

\begin{example}
\label{ex:legendre-computation}
Is $-3$ a quadratic residue modulo $23$?
We compute
\[
  \legendre{-3}{23} = \legendre{-1}{23}\legendre{3}{23}.
\]
Since $23 \equiv 3 \pmod{4}$, we have $\legendre{-1}{23} = (-1)^{11} = -1$.
For $\legendre{3}{23}$, we use Euler's criterion: $3^{11} \bmod 23$.
Compute: $3^2 = 9$, $3^4 = 81 \equiv 12$, $3^8 \equiv 144 \equiv 6$,
so $3^{11} = 3^8 \cdot 3^2 \cdot 3 \equiv 6 \cdot 9 \cdot 3 = 162 \equiv 162 - 7 \cdot 23 = 1$.
Thus $\legendre{3}{23} = 1$ and $\legendre{-3}{23} = (-1)(1) = -1$; hence $-3$ is a QNR mod~$23$.
\end{example}

% ------------------------------------------------------------
\section{Gauss's Lemma}
\label{sec:gauss-lemma}

\begin{theorem}[Gauss's lemma]
\label{thm:gauss-lemma}
\index{Gauss's lemma}
Let $p$ be an odd prime and $a$ an integer with $p \nmid a$. Consider the set
\[
  S = \Bigl\{ a \cdot 1, \; a \cdot 2, \; \ldots, \; a \cdot \tfrac{p-1}{2} \Bigr\}
\]
reduced modulo $p$ into the range $\{1, 2, \ldots, p-1\}$.
Let $\nu$ be the number of elements of this reduced set that are greater than $\frac{p}{2}$
\textup{(}i.e., lie in $\bigl\{\frac{p+1}{2}, \ldots, p-1\bigr\}$\textup{)}.
Then
\begin{equation}
\label{eq:gauss-lemma}
  \legendre{a}{p} = (-1)^\nu.
\end{equation}
\end{theorem}

\begin{proof}
\index{Gauss's lemma!proof}
For $j = 1, 2, \ldots, \frac{p-1}{2}$, let $r_j$ be the least positive residue of $aj$
modulo $p$, so $r_j \in \{1, \ldots, p-1\}$.
Define
\[
  \varepsilon_j =
  \begin{cases}
    +1 & \text{if } r_j \leq \frac{p-1}{2},\\
    -1 & \text{if } r_j > \frac{p-1}{2},
  \end{cases}
  \qquad \text{and} \qquad
  s_j =
  \begin{cases}
    r_j     & \text{if } r_j \leq \frac{p-1}{2},\\
    p - r_j & \text{if } r_j > \frac{p-1}{2}.
  \end{cases}
\]
Then $r_j \equiv \varepsilon_j \, s_j \pmod{p}$ (since if $r_j > \frac{p-1}{2}$,
then $p - r_j \equiv -r_j \pmod{p}$), and $s_j \in \bigl\{1, \ldots, \frac{p-1}{2}\bigr\}$.

\medskip\noindent\textbf{Claim:} The values $s_1, s_2, \ldots, s_{(p-1)/2}$ are a permutation
of $\bigl\{1, 2, \ldots, \frac{p-1}{2}\bigr\}$.

\emph{Proof of claim.}
It suffices to show they are all distinct. Suppose $s_j = s_k$ for $j \neq k$.
Then $r_j \equiv \pm r_k \pmod{p}$. If $r_j \equiv r_k$, then $aj \equiv ak$,
so $j \equiv k \pmod{p}$, impossible since $1 \leq j, k \leq \frac{p-1}{2} < p$.
If $r_j \equiv -r_k$, then $a(j + k) \equiv 0 \pmod{p}$, so $p \mid (j + k)$,
but $2 \leq j + k \leq p - 1$, contradiction. This proves the claim.

\medskip
Now multiply all the congruences $aj \equiv \varepsilon_j \, s_j \pmod{p}$
for $j = 1, \ldots, \frac{p-1}{2}$:
\[
  a^{(p-1)/2} \cdot \prod_{j=1}^{(p-1)/2} j \;\equiv\;
  \Bigl(\prod_{j=1}^{(p-1)/2} \varepsilon_j\Bigr)
  \cdot \prod_{j=1}^{(p-1)/2} s_j
  \pmod{p}.
\]
Since $\{s_j\}$ is a permutation of $\bigl\{1, \ldots, \frac{p-1}{2}\bigr\}$,
both products $\prod j$ and $\prod s_j$ equal $\bigl(\frac{p-1}{2}\bigr)!$,
which is coprime to $p$. Cancelling gives
\[
  a^{(p-1)/2} \equiv \prod_{j=1}^{(p-1)/2} \varepsilon_j = (-1)^\nu \pmod{p}.
\]
By Euler's criterion, $a^{(p-1)/2} \equiv \legendre{a}{p} \pmod{p}$, so
$\legendre{a}{p} = (-1)^\nu$.
\end{proof}

\begin{example}
\label{ex:gauss-lemma-2mod7}
Let $p = 7$, $a = 2$. The set $S = \{2, 4, 6\}$ reduced modulo $7$ is $\{2, 4, 6\}$.
The elements greater than $\frac{7}{2} = 3.5$ are $4$ and $6$, so $\nu = 2$.
Hence $\legendre{2}{7} = (-1)^2 = 1$. Indeed, $3^2 = 9 \equiv 2 \pmod{7}$.
\end{example}

% ------------------------------------------------------------
\section{The First and Second Supplements}
\label{sec:supplements}

Using Gauss's lemma, we can determine $\legendre{-1}{p}$ and $\legendre{2}{p}$.

\begin{theorem}[First supplement]
\label{thm:first-supplement}
\index{quadratic reciprocity!first supplement}
For any odd prime $p$,
\begin{equation}
\label{eq:first-supplement}
  \legendre{-1}{p} = (-1)^{(p-1)/2} =
  \begin{cases}
    \phantom{-}1 & \text{if } p \equiv 1 \pmod{4},\\
    -1           & \text{if } p \equiv 3 \pmod{4}.
  \end{cases}
\end{equation}
\end{theorem}

\begin{proof}
This follows immediately from Euler's criterion: $\legendre{-1}{p} \equiv (-1)^{(p-1)/2} \pmod{p}$.
Since both sides are $\pm 1$ and $p \geq 3$, congruence implies equality.
\end{proof}

\begin{theorem}[Second supplement]
\label{thm:second-supplement}
\index{quadratic reciprocity!second supplement}
For any odd prime $p$,
\begin{equation}
\label{eq:second-supplement}
  \legendre{2}{p} = (-1)^{(p^2-1)/8} =
  \begin{cases}
    \phantom{-}1 & \text{if } p \equiv \pm 1 \pmod{8},\\
    -1           & \text{if } p \equiv \pm 3 \pmod{8}.
  \end{cases}
\end{equation}
\end{theorem}

\begin{proof}
\index{quadratic reciprocity!second supplement!proof}
We apply Gauss's lemma with $a = 2$. The set is $S = \{2, 4, 6, \ldots, p-1\}$,
and we need to count how many of the elements $2j$ (for $j = 1, \ldots, \frac{p-1}{2}$)
satisfy $2j > \frac{p}{2}$, i.e., $j > \frac{p}{4}$.

The number of such $j$ is
\[
  \nu = \frac{p-1}{2} - \floor{\frac{p}{4}} = \frac{p-1}{2} - \floor{\frac{p}{4}}.
\]
A case-by-case check on $p \bmod 8$ shows:
\begin{center}
\begin{tabular}{c|cccc}
  \toprule
  $p \bmod 8$ & $1$ & $3$ & $5$ & $7$ \\
  \midrule
  $\frac{p-1}{2}$ & $\frac{p-1}{2}$ & $\frac{p-1}{2}$ & $\frac{p-1}{2}$ & $\frac{p-1}{2}$ \\[3pt]
  $\floor{p/4}$ & $\frac{p-1}{4}$ & $\frac{p-3}{4}$ & $\frac{p-1}{4}$ & $\frac{p-3}{4}$ \\[3pt]
  $\nu$ & $\frac{p-1}{4}$ & $\frac{p+1}{4}$ & $\frac{p-1}{4}$ & $\frac{p+1}{4}$ \\[3pt]
  $\nu \bmod 2$ & $0$ & $1$ & $1$ & $0$ \\
  \bottomrule
\end{tabular}
\end{center}
Thus $(-1)^\nu = 1$ when $p \equiv \pm 1 \pmod{8}$ and $(-1)^\nu = -1$ when
$p \equiv \pm 3 \pmod{8}$.

It remains to verify that $(-1)^\nu = (-1)^{(p^2-1)/8}$.
Observe that $p^2 - 1 = (p-1)(p+1)$. Writing $p = 8k + r$ for $r \in \{1,3,5,7\}$:
\begin{align*}
  p \equiv 1 \pmod{8}&\colon \tfrac{p^2-1}{8} = \tfrac{(8k)(8k+2)}{8} = k(4k+1), \text{ even;}\\
  p \equiv 3 \pmod{8}&\colon \tfrac{p^2-1}{8} = \tfrac{(8k+2)(8k+4)}{8} = (4k+1)(2k+1), \text{ odd;}\\
  p \equiv 5 \pmod{8}&\colon \tfrac{p^2-1}{8} = \tfrac{(8k+4)(8k+6)}{8} = (2k+1)(4k+3), \text{ odd;}\\
  p \equiv 7 \pmod{8}&\colon \tfrac{p^2-1}{8} = \tfrac{(8k+6)(8k+8)}{8} = (4k+3)(k+1), \text{ even.}
\end{align*}
This matches the table above, completing the proof.
\end{proof}

% ------------------------------------------------------------
\section{Quadratic Reciprocity}
\label{sec:qr-law}

\begin{theorem}[Law of quadratic reciprocity]
\label{thm:qr-law}
\index{quadratic reciprocity!law of}
\index{theorema aureum|see{quadratic reciprocity}}
Let $p$ and $q$ be distinct odd primes. Then
\begin{equation}
\label{eq:qr-law}
  \legendre{p}{q}\legendre{q}{p} = (-1)^{\frac{p-1}{2}\cdot\frac{q-1}{2}}.
\end{equation}
Equivalently:
\begin{enumerate}[label=\textup{(\alph*)}]
  \item If $p \equiv 1 \pmod{4}$ or $q \equiv 1 \pmod{4}$
    \textup{(}i.e., not both $\equiv 3 \pmod 4$\textup{)},
    then $\legendre{p}{q} = \legendre{q}{p}$.
  \item If $p \equiv q \equiv 3 \pmod{4}$, then $\legendre{p}{q} = -\legendre{q}{p}$.
\end{enumerate}
\end{theorem}

We present Eisenstein's elegant proof\index{Eisenstein, Gotthold}\index{quadratic reciprocity!Eisenstein's proof}
via lattice-point counting.

\begin{proof}
Consider the rectangle $R = (0, p) \times (0, q)$ in $\R^2$ and the line
$\ell \colon qx = py$ passing through the origin and the corner $(p, q)$.

\medskip\noindent\textbf{Step 1: Relate the Legendre symbols to lattice-point counts.}

By Gauss's lemma (Theorem~\ref{thm:gauss-lemma}), $\legendre{q}{p} = (-1)^{\mu}$
where $\mu$ is the number of $j \in \{1, \ldots, \frac{p-1}{2}\}$ such that
the least positive residue of $jq$ modulo $p$ exceeds $\frac{p}{2}$.
One can reformulate this count as follows.
For each integer $j$ with $1 \leq j \leq \frac{p-1}{2}$, the residue of $jq$
modulo $p$ exceeds $\frac{p}{2}$ if and only if $\floor{\frac{jq}{p}}$ and $jq$
have different parities.
A careful parity analysis (see Lemma~\ref{lem:eisenstein-key} below) shows that
\begin{equation}
\label{eq:gauss-lemma-floor}
  \legendre{q}{p} = (-1)^{\sum_{j=1}^{(p-1)/2} \floor{jq/p}}.
\end{equation}

\medskip\noindent\textbf{Step 2: Count lattice points.}

The quantity $\floor{\frac{jq}{p}}$ equals the number of integers $k$ with
$1 \leq k \leq \frac{jq}{p}$, i.e., the number of lattice points $(j, k)$ with
$1 \leq k < \frac{jq}{p}$ (strict inequality since $p \nmid jq$ for
$1 \leq j \leq \frac{p-1}{2}$). Thus
\[
  \sum_{j=1}^{(p-1)/2} \floor{\frac{jq}{p}}
  = \#\Bigl\{(j,k) \in \Z^2 : 1 \leq j \leq \tfrac{p-1}{2},\;
    1 \leq k,\; k < \tfrac{jq}{p}\Bigr\}.
\]
This is the number of lattice points in the interior of the triangle $T_1$
below the line $\ell$ with $1 \leq j \leq \frac{p-1}{2}$ and $k \geq 1$.

Similarly,
\[
  \sum_{k=1}^{(q-1)/2} \floor{\frac{kp}{q}}
  = \#\Bigl\{(j,k) \in \Z^2 : 1 \leq k \leq \tfrac{q-1}{2},\;
    1 \leq j,\; j < \tfrac{kp}{q}\Bigr\},
\]
which counts lattice points in the triangle $T_2$ above the line $\ell$
with $1 \leq k \leq \frac{q-1}{2}$ and $j \geq 1$.

\medskip\noindent\textbf{Step 3: No lattice points on the line.}

Since $p$ and $q$ are distinct primes, the equation $qj = pk$ with
$1 \leq j \leq \frac{p-1}{2}$ and $1 \leq k \leq \frac{q-1}{2}$ would
require $p \mid j$ and $q \mid k$, which is impossible in these ranges.
So no lattice point in the rectangle $[1, \frac{p-1}{2}] \times [1, \frac{q-1}{2}]$
lies on $\ell$.

\medskip\noindent\textbf{Step 4: Combine.}

Every lattice point $(j,k)$ with $1 \leq j \leq \frac{p-1}{2}$,
$1 \leq k \leq \frac{q-1}{2}$ lies either in $T_1$ (below $\ell$) or
$T_2$ (above $\ell$). Therefore
\[
  \sum_{j=1}^{(p-1)/2} \floor{\frac{jq}{p}}
  + \sum_{k=1}^{(q-1)/2} \floor{\frac{kp}{q}}
  = \frac{p-1}{2} \cdot \frac{q-1}{2}.
\]
By \eqref{eq:gauss-lemma-floor} and its analogue for $\legendre{p}{q}$:
\[
  \legendre{q}{p}\legendre{p}{q}
  = (-1)^{\sum \floor{jq/p} + \sum \floor{kp/q}}
  = (-1)^{\frac{p-1}{2}\cdot\frac{q-1}{2}}. \qedhere
\]
\end{proof}

We now state and prove the key lemma used in Step~1.

\begin{lemma}[Eisenstein's reformulation of Gauss's lemma]
\label{lem:eisenstein-key}
\index{Gauss's lemma!Eisenstein's reformulation}
Let $p$ be an odd prime and $a$ an integer with $p \nmid a$ and $a$ odd. Then
\[
  \legendre{a}{p} = (-1)^{\sum_{j=1}^{(p-1)/2} \floor{ja/p}}.
\]
\end{lemma}

\begin{proof}
With notation as in the proof of Gauss's lemma (Theorem~\ref{thm:gauss-lemma}),
we have $aj = \floor{\frac{aj}{p}} \cdot p + r_j$, where $r_j$ is the least
positive residue. Recall that $s_j = r_j$ when $r_j \leq \frac{p-1}{2}$ and
$s_j = p - r_j$ otherwise, and $\{s_j\}$ is a permutation of
$\{1, \ldots, \frac{p-1}{2}\}$.

Summing $aj = \floor{\frac{aj}{p}} p + r_j$ over $j = 1, \ldots, \frac{p-1}{2}$:
\[
  a \sum_{j=1}^{(p-1)/2} j = p \sum_{j=1}^{(p-1)/2} \floor{\frac{aj}{p}}
    + \sum_{j=1}^{(p-1)/2} r_j.
\]
Also, $r_j = s_j$ when $\varepsilon_j = +1$ and $r_j = p - s_j$ when $\varepsilon_j = -1$,
so $\sum r_j = \sum s_j + \nu p - 2 \sum_{j:\, \varepsilon_j = -1} s_j$. But actually,
more simply: $r_j + s_j = 0$ or $r_j + s_j = p$ depending on whether $\varepsilon_j = +1$ or $-1$.
So $r_j = s_j + \frac{1 - \varepsilon_j}{2} \cdot (p - 2s_j)$.

It is cleaner to note:
\[
  \sum r_j = \sum_{\varepsilon_j = 1} s_j + \sum_{\varepsilon_j = -1} (p - s_j)
  = \sum s_j + \nu p - 2\!\!\sum_{\varepsilon_j = -1}\!\! s_j.
\]
Since $\{s_j\}$ is a permutation of $\{1, \ldots, \frac{p-1}{2}\}$,
we have $\sum s_j = \sum_{j=1}^{(p-1)/2} j$. Substituting:
\[
  a \sum j = p \sum \floor{\frac{aj}{p}} + \sum j + \nu p
  - 2\!\!\sum_{\varepsilon_j = -1}\!\! s_j.
\]
Reducing modulo $2$ (note $p$ is odd, so $p \equiv 1 \pmod{2}$):
\[
  (a-1)\sum j \equiv \sum \floor{\frac{aj}{p}} + \nu \pmod{2}.
\]
Since $a$ is odd, $a - 1$ is even, so the left side is $\equiv 0 \pmod{2}$.
Therefore $\nu \equiv \sum \floor{\frac{aj}{p}} \pmod{2}$, and
\[
  \legendre{a}{p} = (-1)^\nu = (-1)^{\sum \floor{aj/p}}. \qedhere
\]
\end{proof}

% --- TikZ: Lattice-point counting for Eisenstein's proof ---
\begin{figure}[ht]
\centering
\begin{tikzpicture}[scale=0.62]
  % Parameters: p=7, q=5
  \def\p{7}
  \def\q{5}
  \pgfmathsetmacro{\halfp}{(\p-1)/2}
  \pgfmathsetmacro{\halfq}{(\q-1)/2}

  % Axes
  \draw[->] (-0.3,0) -- (\halfp+0.8,0) node[right] {$j$};
  \draw[->] (0,-0.3) -- (0,\halfq+0.8) node[above] {$k$};

  % Rectangle boundary (dashed)
  \draw[dashed, gray] (0,0) rectangle (\halfp,\halfq);

  % The line qx = py, i.e., y = (q/p)x
  \draw[thick, blue!70!black] (0,0) -- (\halfp+0.3, {(\halfp+0.3)*\q/\p})
    node[right, font=\small] {$\ell\colon qj = pk$};

  % Lattice points below the line (triangle T1)
  % Manual placement for p=7, q=5
  % j=1: floor(5/7)=0, no points
  % j=2: floor(10/7)=1, k=1: 1 < 10/7=1.43, yes
  \fill[red!70!black] (2,1) circle (3pt);
  % j=3: floor(15/7)=2, k=1: yes, k=2: 2 < 15/7=2.14, yes
  \fill[red!70!black] (3,1) circle (3pt);
  \fill[red!70!black] (3,2) circle (3pt);

  % Lattice points above the line (triangle T2)
  % k=1: floor(7/5)=1, j=1: 1 < 7/5=1.4, yes
  \fill[green!50!black] (1,1) circle (3pt);
  % k=2: floor(14/5)=2, j=1: yes, j=2: 2 < 14/5=2.8, yes
  \fill[green!50!black] (1,2) circle (3pt);
  \fill[green!50!black] (2,2) circle (3pt);

  % Shade triangles lightly
  \fill[red!10] (0,0) -- (\halfp,0) -- (\halfp, {\halfp*\q/\p}) -- cycle;
  \fill[green!10] (0,0) -- (0,\halfq) -- ({\halfq*\p/\q}, \halfq) -- cycle;

  % Tick labels
  \foreach \j in {1,...,3} {
    \node[below, font=\small] at (\j,0) {$\j$};
  }
  \foreach \k in {1,...,2} {
    \node[left, font=\small] at (0,\k) {$\k$};
  }

  % Labels for half-values
  \node[below right, font=\small] at (\halfp, -0.15) {$\frac{p-1}{2}$};
  \node[left, font=\small] at (-0.15, \halfq) {$\frac{q-1}{2}$};

  % Legend
  \node[red!70!black, font=\small] at (5.5, 1.8)
    {$\bullet$ below $\ell$: $\sum\floor{\frac{jq}{p}}$};
  \node[green!50!black, font=\small] at (5.5, 1.2)
    {$\bullet$ above $\ell$: $\sum\floor{\frac{kp}{q}}$};
  \node[font=\small, text=gray] at (5.5, 0.5)
    {Total: $3 \times 2 = 6$};

  % Title
  \node[above, font=\bfseries\small] at (1.5, \halfq+0.8) {$p=7,\; q=5$};
\end{tikzpicture}
\caption{Eisenstein's lattice-point proof for $p=7$, $q=5$.
  The $3 + 3 = 6 = \frac{6}{2}\cdot\frac{4}{2}$ lattice points split
  between the two triangles below and above the line $\ell$.}
\label{fig:eisenstein-lattice}
\end{figure}

\begin{example}[Applying quadratic reciprocity]
\label{ex:qr-application}
\index{quadratic reciprocity!example}
Is $5$ a quadratic residue modulo $41$? Since $41 \equiv 1 \pmod{4}$,
quadratic reciprocity gives
\[
  \legendre{5}{41} = \legendre{41}{5} = \legendre{1}{5} = 1.
\]
So $5$ is a QR mod~$41$.
Indeed, $28^2 = 784 = 19 \cdot 41 + 5$, confirming $28^2 \equiv 5 \pmod{41}$.
\end{example}

\begin{example}[A more involved computation]
\label{ex:qr-involved}
Evaluate $\legendre{70}{571}$. We factor and use multiplicativity:
\[
  \legendre{70}{571} = \legendre{2}{571}\legendre{5}{571}\legendre{7}{571}.
\]
Since $571 \equiv 3 \pmod{8}$, the second supplement gives $\legendre{2}{571} = -1$.

For $\legendre{5}{571}$: since $571 \equiv 3 \pmod{4}$ and $5 \equiv 1 \pmod{4}$
(not both $\equiv 3$), reciprocity gives
$\legendre{5}{571} = \legendre{571}{5} = \legendre{1}{5} = 1$.

For $\legendre{7}{571}$: both $7 \equiv 3$ and $571 \equiv 3 \pmod{4}$, so
$\legendre{7}{571} = -\legendre{571}{7} = -\legendre{4}{7} = -(1) = -1$,
since $571 = 81 \cdot 7 + 4$ and $\legendre{4}{7} = \legendre{2^2}{7} = 1$.

Therefore $\legendre{70}{571} = (-1)(1)(-1) = 1$.
\end{example}

% --- TikZ: QR pattern visualisation ---
\begin{figure}[ht]
\centering
\begin{tikzpicture}[scale=0.38]
  % Visualize QR pattern for small primes
  % Rows = primes p, columns = a in {1,...,p-1}
  % Color: blue = QR, red = QNR
  \def\primes{{3,5,7,11,13}}
  \def\qrdata{%
    {1,0,0,0,0,0,0,0,0,0,0,0},%  p=3:  QR={1}
    {1,0,0,1,0,0,0,0,0,0,0,0},%  p=5:  QR={1,4}
    {1,1,0,1,0,1,0,0,0,0,0,0},%  p=7:  QR={1,2,4}
    {1,0,1,1,1,0,0,0,1,0,0,0},%  p=11: QR={1,3,4,5,9}
    {1,0,1,1,0,0,0,0,1,1,0,1}%   p=13: QR={1,3,4,9,10,12}
  }
  \def\primelist{3,5,7,11,13}
  \def\pcount{5}

  % Header
  \node[font=\bfseries\small] at (6.5, 6) {Quadratic Residue Patterns};

  % Column headers
  \foreach \a in {1,...,12} {
    \node[font=\tiny] at (\a, 5.2) {$\a$};
  }

  % Rows
  \node[font=\small, left] at (0.3, 4) {$p=3$};
  \fill[blue!50] (1,3.6) rectangle (1.8,4.4); % 1 is QR
  \fill[red!30] (2,3.6) rectangle (2.8,4.4);  % 2 is QNR

  \node[font=\small, left] at (0.3, 3) {$p=5$};
  \fill[blue!50] (1,2.6) rectangle (1.8,3.4);
  \fill[red!30] (2,2.6) rectangle (2.8,3.4);
  \fill[red!30] (3,2.6) rectangle (3.8,3.4);
  \fill[blue!50] (4,2.6) rectangle (4.8,3.4);

  \node[font=\small, left] at (0.3, 2) {$p=7$};
  \fill[blue!50] (1,1.6) rectangle (1.8,2.4);
  \fill[blue!50] (2,1.6) rectangle (2.8,2.4);
  \fill[red!30] (3,1.6) rectangle (3.8,2.4);
  \fill[blue!50] (4,1.6) rectangle (4.8,2.4);
  \fill[red!30] (5,1.6) rectangle (5.8,2.4);
  \fill[red!30] (6,1.6) rectangle (6.8,2.4);

  \node[font=\small, left] at (0.3, 1) {$p=11$};
  \fill[blue!50] (1,0.6) rectangle (1.8,1.4);
  \fill[red!30] (2,0.6) rectangle (2.8,1.4);
  \fill[blue!50] (3,0.6) rectangle (3.8,1.4);
  \fill[blue!50] (4,0.6) rectangle (4.8,1.4);
  \fill[blue!50] (5,0.6) rectangle (5.8,1.4);
  \fill[red!30] (6,0.6) rectangle (6.8,1.4);
  \fill[red!30] (7,0.6) rectangle (7.8,1.4);
  \fill[red!30] (8,0.6) rectangle (8.8,1.4);
  \fill[blue!50] (9,0.6) rectangle (9.8,1.4);
  \fill[red!30] (10,0.6) rectangle (10.8,1.4);

  \node[font=\small, left] at (0.3, 0) {$p=13$};
  \fill[blue!50] (1,-0.4) rectangle (1.8,0.4);
  \fill[red!30] (2,-0.4) rectangle (2.8,0.4);
  \fill[blue!50] (3,-0.4) rectangle (3.8,0.4);
  \fill[blue!50] (4,-0.4) rectangle (4.8,0.4);
  \fill[red!30] (5,-0.4) rectangle (5.8,0.4);
  \fill[red!30] (6,-0.4) rectangle (6.8,0.4);
  \fill[red!30] (7,-0.4) rectangle (7.8,0.4);
  \fill[red!30] (8,-0.4) rectangle (8.8,0.4);
  \fill[blue!50] (9,-0.4) rectangle (9.8,0.4);
  \fill[blue!50] (10,-0.4) rectangle (10.8,0.4);
  \fill[red!30] (11,-0.4) rectangle (11.8,0.4);
  \fill[blue!50] (12,-0.4) rectangle (12.8,0.4);

  % Legend
  \fill[blue!50] (2,-1.5) rectangle (2.8,-0.9);
  \node[font=\small, right] at (3,-1.2) {QR};
  \fill[red!30] (5,-1.5) rectangle (5.8,-0.9);
  \node[font=\small, right] at (6,-1.2) {QNR};
\end{tikzpicture}
\caption{Quadratic residue patterns for small primes: blue cells are QRs,
  red cells are QNRs. Note the symmetric/antisymmetric patterns governed by
  $\legendre{a}{p} = \legendre{p-a}{p}$ when $p \equiv 1 \pmod{4}$.}
\label{fig:qr-pattern}
\end{figure}

% ------------------------------------------------------------
\section{The Jacobi Symbol}
\label{sec:jacobi-symbol}

The Legendre symbol is defined only for prime moduli. The Jacobi symbol extends
it to arbitrary odd moduli, which greatly facilitates computations.

\begin{definition}[Jacobi symbol]
\label{def:jacobi}
\index{Jacobi symbol}
Let $n > 1$ be an odd integer with prime factorisation $n = p_1^{e_1} p_2^{e_2} \cdots p_r^{e_r}$.
For any integer $a$, the \emph{Jacobi symbol} is
\[
  \jacobi{a}{n} = \legendre{a}{p_1}^{e_1} \legendre{a}{p_2}^{e_2} \cdots \legendre{a}{p_r}^{e_r}.
\]
\end{definition}

\begin{remark}
\label{rem:jacobi-warning}
\index{Jacobi symbol!warning}
The Jacobi symbol $\jacobi{a}{n} = 1$ does \emph{not} imply that $a$ is a quadratic
residue modulo $n$. For instance, $\jacobi{2}{9} = \legendre{2}{3}^2 = (-1)^2 = 1$,
but $x^2 \equiv 2 \pmod{9}$ has no solution (checking $x = 0, 1, \ldots, 8$).
However, $\jacobi{a}{n} = -1$ \emph{does} imply that $a$ is a QNR modulo $n$.
\end{remark}

\begin{proposition}[Properties of the Jacobi symbol]
\label{prop:jacobi-properties}
\index{Jacobi symbol!properties}
Let $m, n$ be odd positive integers and $a, b \in \Z$. Then:
\begin{enumerate}[label=\textup{(\roman*)}]
  \item $\jacobi{ab}{n} = \jacobi{a}{n}\jacobi{b}{n}$.
  \item $\jacobi{a}{mn} = \jacobi{a}{m}\jacobi{a}{n}$.
  \item If $a \equiv b \pmod{n}$, then $\jacobi{a}{n} = \jacobi{b}{n}$.
  \item $\jacobi{1}{n} = 1$ and $\jacobi{-1}{n} = (-1)^{(n-1)/2}$.
  \item $\jacobi{2}{n} = (-1)^{(n^2-1)/8}$.
\end{enumerate}
\end{proposition}

\begin{proof}
Properties (i)--(iii) follow directly from the definition and the corresponding
properties of the Legendre symbol.

For (iv), $\jacobi{-1}{n} = \prod \legendre{-1}{p_i}^{e_i} = \prod (-1)^{e_i(p_i-1)/2}
= (-1)^{\sum e_i(p_i-1)/2}$.
One verifies that $\sum e_i(p_i-1)/2 \equiv (n-1)/2 \pmod{2}$ by induction on the
number of prime factors, using the identity: if $n = ab$ with $a, b$ odd, then
$\frac{ab-1}{2} \equiv \frac{a-1}{2} + \frac{b-1}{2} \pmod{2}$.

Property (v) follows similarly from the identity
$\frac{(ab)^2 - 1}{8} \equiv \frac{a^2-1}{8} + \frac{b^2-1}{8} \pmod{2}$
for odd $a, b$.
\end{proof}

\begin{theorem}[Reciprocity for the Jacobi symbol]
\label{thm:jacobi-reciprocity}
\index{Jacobi symbol!reciprocity}
\index{quadratic reciprocity!Jacobi symbol}
Let $m, n$ be odd positive integers with $\gcd(m, n) = 1$. Then
\begin{equation}
\label{eq:jacobi-reciprocity}
  \jacobi{m}{n}\jacobi{n}{m} = (-1)^{\frac{m-1}{2}\cdot\frac{n-1}{2}}.
\end{equation}
\end{theorem}

\begin{proof}
Write $m = p_1 \cdots p_s$ and $n = q_1 \cdots q_t$ (with repetition according
to multiplicity). Then
\[
  \jacobi{m}{n}\jacobi{n}{m}
  = \prod_{i,j} \legendre{p_i}{q_j}\legendre{q_j}{p_i}
  = \prod_{i,j} (-1)^{\frac{p_i-1}{2}\cdot\frac{q_j-1}{2}}
  = (-1)^{\sum_{i,j}\frac{p_i-1}{2}\cdot\frac{q_j-1}{2}}.
\]
Now $\sum_{i,j} \frac{p_i-1}{2}\cdot\frac{q_j-1}{2}
= \bigl(\sum_i \frac{p_i-1}{2}\bigr)\bigl(\sum_j \frac{q_j-1}{2}\bigr)$.
By the same parity identity used in Proposition~\ref{prop:jacobi-properties}(iv),
$\sum_i \frac{p_i-1}{2} \equiv \frac{m-1}{2} \pmod{2}$ and likewise for $n$.
\end{proof}

\begin{example}[Efficient computation with the Jacobi symbol]
\label{ex:jacobi-computation}
\index{Jacobi symbol!computation}
Evaluate $\jacobi{1001}{9907}$, where $9907$ is prime (so the Jacobi symbol
equals the Legendre symbol and determines quadratic residuosity).
\begin{align*}
  \jacobi{1001}{9907}
  &= \jacobi{7}{9907}\jacobi{11}{9907}\jacobi{13}{9907}. \\
\intertext{Since $9907 \equiv 3 \pmod{4}$, $7 \equiv 3$, $11 \equiv 3$, $13 \equiv 1$:}
  \jacobi{7}{9907} &= -\jacobi{9907}{7} = -\jacobi{3}{7},\\
  \jacobi{3}{7} &= \jacobi{7}{3} \cdot (-1)^{\frac{2}{2}\cdot\frac{6}{2}}
    = \jacobi{7}{3}\cdot(-1)^3 = -\jacobi{1}{3} = -1,\\
  \text{so }\jacobi{7}{9907} &= -(-1) = 1.\\[4pt]
  \jacobi{11}{9907} &= -\jacobi{9907}{11} = -\jacobi{9907 \bmod 11}{11} = -\jacobi{7}{11},\\
  \jacobi{7}{11} &= \jacobi{11}{7}\cdot(-1)^{\frac{6}{2}\cdot\frac{10}{2}}
    = \jacobi{11}{7}\cdot(-1)^{15} = -\jacobi{4}{7} = -(1) = -1,\\
  \text{so }\jacobi{11}{9907} &= -(-1) = 1.\\[4pt]
  \jacobi{13}{9907} &= \jacobi{9907}{13} = \jacobi{9907 \bmod 13}{13} = \jacobi{3}{13},\\
  \jacobi{3}{13} &= \jacobi{13}{3}\cdot(-1)^{\frac{2}{2}\cdot\frac{12}{2}}
    = \jacobi{1}{3}\cdot(-1)^6 = 1.
\end{align*}
Therefore $\jacobi{1001}{9907} = 1 \cdot 1 \cdot 1 = 1$,
so $1001$ is a quadratic residue modulo $9907$.
\end{example}

\begin{remark}[Computational complexity]
\label{rem:jacobi-complexity}
\index{Jacobi symbol!algorithm}
The Jacobi symbol $\jacobi{a}{n}$ can be computed in $O(\log^2 n)$ bit operations
using a procedure analogous to the Euclidean algorithm: repeatedly apply reciprocity
and reduce modulo the smaller argument. This does \emph{not} require factoring $n$,
which is why the Jacobi symbol is useful in primality testing
(e.g., the Solovay--Strassen test\index{Solovay--Strassen test}).
\end{remark}

% ------------------------------------------------------------
\section{Exercises}
\label{sec:qr-exercises}

\begin{exercise}
\label{exer:qr-list}
List all quadratic residues modulo $17$ and modulo $19$.
\end{exercise}

\begin{exercise}
\label{exer:euler-criterion-compute}
Use Euler's criterion to determine whether $3$ is a quadratic residue modulo $31$.
\end{exercise}

\begin{exercise}
\label{exer:qr-product}
\index{quadratic residue!product}
Let $p$ be an odd prime. Show that the product of all quadratic residues
modulo $p$ is congruent to $(-1)^{(p+1)/2} \pmod{p}$.
\textup{(}Hint: pair each QR $a$ with its ``complement'' $a^{-1}$.\textup{)}
\end{exercise}

\begin{exercise}
\label{exer:minus-one-sum}
Prove that for an odd prime $p$,
\[
  \sum_{a=1}^{p-1} \legendre{a}{p} = 0.
\]
\end{exercise}

\begin{exercise}
\label{exer:qr-minus3}
Using quadratic reciprocity and the supplements, determine exactly which primes $p$
satisfy $\legendre{-3}{p} = 1$. Show that $-3$ is a QR mod $p$ if and only if
$p = 3$ or $p \equiv 1 \pmod{3}$.
\end{exercise}

\begin{exercise}
\label{exer:gauss-lemma-apply}
Apply Gauss's lemma directly (without Euler's criterion) to compute
$\legendre{3}{11}$ and $\legendre{5}{13}$.
\end{exercise}

\begin{exercise}
\label{exer:second-supplement-derive}
Give an alternative derivation of the second supplement $\legendre{2}{p} = (-1)^{(p^2-1)/8}$
using Euler's criterion and the identity
\[
  2^{(p-1)/2} \equiv (-1)^{(p^2-1)/8} \pmod{p}
\]
by working in a suitable extension (e.g., $\Z[i]$ or $\Z[\sqrt{2}]$).
\end{exercise}

\begin{exercise}
\label{exer:jacobi-compute}
Compute $\jacobi{295}{1009}$ using only the properties of the Jacobi symbol
(reciprocity, supplements, multiplicativity). Is $295$ a QR mod $1009$?
\textup{(}Note: $1009$ is prime.\textup{)}
\end{exercise}

\begin{exercise}
\label{exer:qr-consecutive}
\index{quadratic residue!consecutive}
Show that for any prime $p > 5$, there exist consecutive integers $a, a+1$
that are both quadratic residues modulo $p$.
\textup{(}Hint: consider $x^2$ and $(x+1)^2$ and count.\textup{)}
\end{exercise}

\begin{exercise}
\label{exer:jacobi-nonresidue}
Give an example showing that $\jacobi{a}{n} = 1$ does not imply $a$ is a QR mod $n$
when $n$ is not prime. Prove that if $n = pq$ with $p \neq q$ odd primes,
then $a$ is a QR mod $n$ if and only if $\legendre{a}{p} = \legendre{a}{q} = 1$.
\end{exercise}

\begin{exercise}[Primes of the form $x^2 + 5y^2$]
\label{exer:qr-representation}
\index{representation by quadratic forms}
Show that if an odd prime $p$ divides $x^2 + 5y^2$ with $\gcd(p, y) = 1$,
then $\legendre{-5}{p} = 1$. Use reciprocity to show this holds if and only if
$p \equiv 1, 3, 7, 9 \pmod{20}$.
\end{exercise}

% ------------------------------------------------------------
\section*{Chapter Summary}
\addcontentsline{toc}{section}{Chapter Summary}

\begin{itemize}
  \item An integer $a$ with $p \nmid a$ is a \textbf{quadratic residue} mod $p$ if
    $x^2 \equiv a \pmod{p}$ is solvable. Exactly half of $\{1, \ldots, p-1\}$ are QRs.

  \item \textbf{Euler's criterion:}
    $\legendre{a}{p} \equiv a^{(p-1)/2} \pmod{p}$.

  \item The \textbf{Legendre symbol} $\legendre{a}{p}$ encodes quadratic residuosity
    and is completely multiplicative.

  \item \textbf{Gauss's lemma} expresses $\legendre{a}{p}$ as $(-1)^\nu$, where
    $\nu$ counts how many of $a, 2a, \ldots, \frac{p-1}{2}a$ have
    residues exceeding $\frac{p}{2}$.

  \item The \textbf{law of quadratic reciprocity}:
    $\legendre{p}{q}\legendre{q}{p} = (-1)^{\frac{p-1}{2}\cdot\frac{q-1}{2}}$
    for distinct odd primes $p, q$.

  \item \textbf{Supplements:} $\legendre{-1}{p} = (-1)^{(p-1)/2}$ and
    $\legendre{2}{p} = (-1)^{(p^2-1)/8}$.

  \item The \textbf{Jacobi symbol} extends the Legendre symbol to composite odd moduli,
    preserving multiplicativity and reciprocity, enabling efficient computation
    without factoring.
\end{itemize}

% === END OF CHUNK ch5-6 ===
% === START OF CHUNK ch7-8 ===

% ============================================================
%  CHAPTER 7 — Representations by Quadratic Forms
% ============================================================
\chapter{Representations by Quadratic Forms}\label{ch:quadratic-forms}
\index{quadratic form}

\begin{quote}
\textit{``Numerorum primorum, qui sunt aggregata duorum quadratorum, insignes
proprietates.''} --- Pierre de Fermat, 1640
\end{quote}

\noindent
One of the oldest and most beautiful questions in number theory asks:
which integers can be expressed as sums of squares?  Fermat claimed in 1640
that every prime $p \equiv 1 \pmod{4}$ is a sum of two squares, but his proof
was never published.  The first complete proofs were given by Euler (1749) after
years of effort.  In this chapter we develop the theory systematically, beginning
with binary quadratic forms, proving Fermat's two-square theorem completely,
characterising all integers representable as sums of two squares, and stating
Lagrange's four-square theorem.

% ------------------------------------------------------------
\section{Binary Quadratic Forms}\label{sec:binary-quad-forms}
\index{binary quadratic form}

\begin{definition}[Binary quadratic form]\label{def:binary-quad-form}
A \textbf{binary quadratic form}\index{binary quadratic form} is a polynomial
\[
  f(x,y) = ax^2 + bxy + cy^2, \qquad a,b,c \in \Z.
\]
The form is called \textbf{positive definite}\index{positive definite form}
if $f(x,y) > 0$ for all $(x,y) \neq (0,0)$.
The \textbf{discriminant}\index{discriminant!of a quadratic form} of $f$ is
\[
  D = b^2 - 4ac.
\]
\end{definition}

\begin{remark}\label{rem:disc-sign}
A binary quadratic form $ax^2 + bxy + cy^2$ with $a > 0$ is positive definite
if and only if $D < 0$.  Indeed, we may complete the square:
\[
  f(x,y) = a\!\left(x + \frac{b}{2a}\,y\right)^{\!2}
          + \frac{4ac - b^2}{4a}\,y^2,
\]
which is positive for all $(x,y)\neq(0,0)$ precisely when $a>0$ and $4ac - b^2 > 0$.
\end{remark}

\begin{definition}[Representation by a form]\label{def:representation}
An integer $n$ is \textbf{represented}\index{representation!by a quadratic form}
by the form $f$ if there exist $x_0, y_0 \in \Z$ with $f(x_0,y_0)=n$.
The representation is \textbf{proper}\index{proper representation} if
$\gcd(x_0,y_0) = 1$.
\end{definition}

\begin{definition}[Equivalence of forms]\label{def:equiv-forms}
\index{equivalence!of quadratic forms}
Two forms $f$ and $g$ are \textbf{(properly) equivalent}, written $f \sim g$,
if there exists a matrix
$\begin{psmallmatrix} \alpha & \beta \\ \gamma & \delta \end{psmallmatrix}
\in \mathrm{SL}_2(\Z)$ such that
\[
  g(x,y) = f(\alpha x + \beta y,\, \gamma x + \delta y).
\]
Equivalent forms have the same discriminant and represent exactly the same
set of integers.
\end{definition}

\begin{example}\label{ex:forms-disc-neg4}
The forms $x^2 + y^2$ and $x^2 + xy + y^2$ have discriminants $D = -4$ and
$D = -3$ respectively.  Both are positive definite.  For $D = -4$, the only
reduced form is $x^2 + y^2$, so the class number is $h(-4) = 1$.
\end{example}

% ------------------------------------------------------------
\section{Sums of Two Squares: Preliminary Results}\label{sec:two-sq-prelim}
\index{sums of two squares}

We aim to determine which integers are representable as $x^2 + y^2$.
The key algebraic tool is the following identity.

\begin{lemma}[Brahmagupta--Fibonacci identity]\label{lem:brahmagupta}
\index{Brahmagupta--Fibonacci identity}
For all integers $a,b,c,d$,
\[
  (a^2 + b^2)(c^2 + d^2) = (ac - bd)^2 + (ad + bc)^2
                          = (ac + bd)^2 + (ad - bc)^2.
\]
\end{lemma}

\begin{proof}
Direct expansion of both sides.  Alternatively, this is the multiplicativity
of the norm in the Gaussian integers\index{Gaussian integers}:
$|z|^2 \cdot |w|^2 = |zw|^2$ where $z = a + bi$ and $w = c + di$.
\end{proof}

\begin{corollary}\label{cor:product-two-sq}
If $m$ and $n$ are each sums of two squares, then so is $mn$.
\end{corollary}

\begin{lemma}\label{lem:minus-one-qr}
If $p$ is an odd prime with $p \equiv 1 \pmod{4}$, then $-1$ is a quadratic
residue modulo~$p$.  That is, there exists $x \in \Z$ with $x^2 \equiv -1 \pmod{p}$.
\end{lemma}

\begin{proof}
By Wilson's theorem\index{Wilson's theorem},
$(p-1)! \equiv -1 \pmod{p}$.  Write
\[
  (p-1)! = 1 \cdot 2 \cdots \frac{p-1}{2} \cdot \frac{p+1}{2} \cdots (p-1).
\]
Since $p-k \equiv -k \pmod{p}$, we get
$(p-1)! \equiv (-1)^{(p-1)/2}\bigl(\tfrac{p-1}{2}!\bigr)^2 \pmod{p}$.
When $p \equiv 1 \pmod{4}$, the exponent $(p-1)/2$ is even, so
$(p-1)! \equiv \bigl(\tfrac{p-1}{2}!\bigr)^2 \pmod{p}$.
Hence $\bigl(\tfrac{p-1}{2}!\bigr)^2 \equiv -1 \pmod{p}$, and
$x = \bigl(\tfrac{p-1}{2}\bigr)!$ works.
\end{proof}

% ------------------------------------------------------------
\section{Fermat's Two-Square Theorem}\label{sec:fermat-two-sq}
\index{Fermat's two-square theorem}

\begin{theorem}[Fermat's two-square theorem]\label{thm:fermat-two-sq}
An odd prime $p$ is a sum of two squares if and only if $p \equiv 1 \pmod{4}$.
\end{theorem}

\begin{proof}
\textbf{Necessity.}\index{Fermat's two-square theorem!necessity}
Suppose $p = x^2 + y^2$.  Since $p$ is odd, $x$ and $y$ have different parities.
Modulo~4, squares are $\equiv 0$ or $1$, so $x^2 + y^2 \equiv 0 + 1 = 1 \pmod{4}$.
Thus $p \equiv 1 \pmod{4}$.

\medskip
\textbf{Sufficiency.}\index{Fermat's two-square theorem!sufficiency}
We give Zagier's celebrated ``one-sentence proof'' restructured for clarity,
then a second self-contained descent proof.

\smallskip
\textit{Step~1: There exists $m$ with $1 \le m < p$ and $mp = x^2 + y^2$.}

By Lemma~\ref{lem:minus-one-qr}, there exists $r$ with $r^2 \equiv -1 \pmod{p}$
and $0 < r < p$.  Then $r^2 + 1 \equiv 0 \pmod{p}$, so $p \mid (r^2 + 1^2)$
and $mp = r^2 + 1$ for some $m$ with $1 \le m < p$ (since $r^2 + 1 < p^2$).

\smallskip
\textit{Step~2 (Descent): Reduce $m$ to $1$.}

Let $m$ be the smallest positive integer such that $mp = x^2 + y^2$ for some
$x,y \in \Z$.  We show $m = 1$.

Suppose for contradiction that $m > 1$.  If $m$ is even, then $x$ and $y$ have
the same parity, and
\[
  \frac{m}{2}\,p = \left(\frac{x+y}{2}\right)^{\!2} + \left(\frac{x-y}{2}\right)^{\!2},
\]
contradicting the minimality of $m$.  So $m$ is odd and $m \ge 3$.

Choose $u,v$ with $u \equiv x \pmod{m}$, $v \equiv y \pmod{m}$, and
$\abs{u}, \abs{v} \le (m-1)/2 < m/2$.  Then
\[
  u^2 + v^2 \equiv x^2 + y^2 \equiv 0 \pmod{m},
\]
so $u^2 + v^2 = m\,m'$ for some positive integer $m'$.  Note that
$u^2 + v^2 \le 2(m/2)^2 = m^2/2 < m^2$, so $m' < m$.

Moreover $m'$ cannot be zero (that would force $u = v = 0$, hence $m \mid x$ and
$m \mid y$, giving $m^2 \mid (x^2+y^2) = mp$, so $m \mid p$; since $1 < m < p$,
this is impossible).

By the Brahmagupta--Fibonacci identity (Lemma~\ref{lem:brahmagupta}),
\[
  m^2\,m'\,p = (u^2+v^2)(x^2+y^2)
            = (ux+vy)^2 + (uy-vx)^2.
\]
Since $ux + vy \equiv x^2 + y^2 \equiv 0 \pmod{m}$ and similarly
$uy - vx \equiv xy - yx = 0 \pmod{m}$, we can divide:
\[
  m'\,p = \left(\frac{ux+vy}{m}\right)^{\!2}
        + \left(\frac{uy-vx}{m}\right)^{\!2},
\]
where both quotients are integers.  Since $0 < m' < m$, this contradicts
the minimality of $m$.

Therefore $m = 1$ and $p = x^2 + y^2$.
\end{proof}

\begin{remark}[Uniqueness]\label{rem:unique-two-sq}
\index{uniqueness of two-square representation}
Fermat also claimed (and Euler proved) that the representation $p = x^2 + y^2$
with $x,y > 0$ and $x \ge y$ is \emph{unique}.  This follows from the fact that
$\Z[i]$ is a unique factorisation domain, or equivalently that the class
number $h(-4) = 1$.
\end{remark}

% ------------------------------------------------------------
\section{The Gaussian Integers Approach}\label{sec:gaussian-integers}
\index{Gaussian integers}

The theory of sums of two squares is intimately connected to the ring
$\Z[i] = \{a + bi : a,b \in \Z\}$ of Gaussian integers.

\begin{definition}[Gaussian integers and norm]\label{def:gaussian-int}
The ring $\Z[i]$\index{$\Z[i]$} has a multiplicative norm
$N(a+bi) = a^2 + b^2$.\index{norm!Gaussian}
An element $\pi \in \Z[i]$ is a \textbf{Gaussian prime}\index{Gaussian prime}
if it is not a unit and its only divisors are units and associates.
\end{definition}

\begin{theorem}[{Primes in $\Z[i]$}]\label{thm:gaussian-primes}
\index{Gaussian integers!primes in}
The Gaussian primes are, up to associates:
\begin{enumerate}[label=(\roman*)]
  \item $1+i$ (with $N(1+i)=2$),
  \item $a + bi$ where $a^2+b^2 = p$ is a rational prime $\equiv 1\pmod{4}$,
  \item rational primes $p \equiv 3 \pmod{4}$ (which remain prime in $\Z[i]$).
\end{enumerate}
\end{theorem}

\begin{center}
\begin{tikzpicture}[scale=0.42]
  \draw[gray!30,very thin] (-6.5,-6.5) grid (6.5,6.5);
  \draw[->,gray] (-6.8,0) -- (6.8,0) node[right]{\small $\operatorname{Re}$};
  \draw[->,gray] (0,-6.8) -- (0,6.8) node[above]{\small $\operatorname{Im}$};
  % Lattice points on circles: r^2 = 1,2,4,5,8,9,10,13,16,17,18,20,25
  % r^2 = 5: (1,2),(2,1),(-1,2),(-2,1),(1,-2),(2,-1),(-1,-2),(-2,-1)
  \foreach \x/\y in {1/2,2/1,-1/2,-2/1,1/-2,2/-1,-1/-2,-2/-1}{
    \fill[blue!70!black] (\x,\y) circle (4pt);
  }
  % r^2 = 10: (1,3),(3,1),(-1,3),(-3,1),(1,-3),(3,-1),(-1,-3),(-3,-1)
  \foreach \x/\y in {1/3,3/1,-1/3,-3/1,1/-3,3/-1,-1/-3,-3/-1}{
    \fill[red!70!black] (\x,\y) circle (4pt);
  }
  % r^2 = 13: (2,3),(3,2), etc.
  \foreach \x/\y in {2/3,3/2,-2/3,-3/2,2/-3,3/-2,-2/-3,-3/-2}{
    \fill[green!50!black] (\x,\y) circle (4pt);
  }
  % r^2 = 25: (0,5),(5,0),(3,4),(4,3), etc.
  \foreach \x/\y in {0/5,5/0,0/-5,-5/0,3/4,4/3,-3/4,-4/3,3/-4,4/-3,-3/-4,-4/-3}{
    \fill[orange!80!black] (\x,\y) circle (3.5pt);
  }
  % Circles
  \draw[blue!50,thin] (0,0) circle ({sqrt(5)});
  \draw[red!50,thin] (0,0) circle ({sqrt(10)});
  \draw[green!40!black,thin] (0,0) circle ({sqrt(13)});
  \draw[orange!60!black,thin] (0,0) circle (5);
  % Labels
  \node[blue!70!black,right] at (2.5,1.5) {\scriptsize $r^2=5$};
  \node[red!70!black,right] at (3.5,1) {\scriptsize $r^2=10$};
  \node[green!50!black,left] at (-3.5,2.5) {\scriptsize $r^2=13$};
  \node[orange!80!black,right] at (4.5,3) {\scriptsize $r^2=25$};
\end{tikzpicture}
\captionof{figure}{Lattice points $(x,y)\in\Z^2$ on circles $x^2+y^2 = n$
for $n = 5, 10, 13, 25$.}\label{fig:lattice-circles}
\index{lattice points!on circles}
\end{center}

% ------------------------------------------------------------
\section{Which Integers Are Sums of Two Squares?}\label{sec:characterise-two-sq}

\begin{theorem}[Characterisation of sums of two squares]\label{thm:two-sq-char}
\index{sums of two squares!characterisation}
A positive integer $n$ is a sum of two squares if and only if every prime
factor $p \equiv 3 \pmod{4}$ appears to an even power in the prime
factorisation of~$n$.
\end{theorem}

\begin{proof}
\textbf{Sufficiency.}
Suppose $n = 2^{a_0} \prod p_i^{a_i} \prod q_j^{2b_j}$ where each $p_i \equiv 1 \pmod 4$
and each $q_j \equiv 3 \pmod 4$.  Now $2 = 1^2 + 1^2$, each $p_i = x_i^2 + y_i^2$
by Theorem~\ref{thm:fermat-two-sq}, and $q_j^{2b_j} = (q_j^{b_j})^2 + 0^2$.
By repeated application of the Brahmagupta--Fibonacci identity
(Corollary~\ref{cor:product-two-sq}), $n$ is a sum of two squares.

\smallskip
\textbf{Necessity.}
Suppose $n = x^2 + y^2$ and let $q \equiv 3 \pmod{4}$ be a prime with
$q^k \| n$ (meaning $q^k \mid n$ but $q^{k+1} \nmid n$).  We must show
$k$ is even.

Let $d = \gcd(x,y)$ and write $x = dx'$, $y = dy'$ with $\gcd(x',y')=1$.
Then $n = d^2(x'^2 + y'^2)$.  Let $q^\ell \| d$, so $q^{2\ell} \| d^2$.
Set $n' = x'^2 + y'^2$; then $q^{k - 2\ell} \| n'$.

If $q \mid n'$, then $q \mid (x'^2 + y'^2)$, so $x'^2 \equiv -y'^2 \pmod{q}$.
Since $\gcd(x',y')=1$, we have $q \nmid y'$, so $(x' y'^{-1})^2 \equiv -1 \pmod{q}$,
meaning $-1$ is a quadratic residue mod~$q$.  But $q \equiv 3 \pmod{4}$ implies
$\legendre{-1}{q} = (-1)^{(q-1)/2} = -1$, a contradiction.

Therefore $q \nmid n'$, so $k - 2\ell = 0$, i.e.\ $k = 2\ell$ is even.
\end{proof}

\begin{example}\label{ex:two-sq-examples}
\leavevmode
\begin{enumerate}[label=(\alph*)]
  \item $45 = 3^2 \cdot 5$.  The prime $3 \equiv 3 \pmod{4}$ appears to an
        even power, and $5 \equiv 1 \pmod{4}$.  Indeed $45 = 6^2 + 3^2$.
  \item $21 = 3 \cdot 7$.  Both $3$ and $7$ are $\equiv 3 \pmod{4}$ and
        appear to odd powers, so $21$ is \emph{not} a sum of two squares.
  \item $50 = 2 \cdot 5^2 = 1^2 + 7^2 = 5^2 + 5^2$.
\end{enumerate}
\end{example}

\begin{proposition}[Counting representations]\label{prop:r2-count}
\index{$r_2(n)$}
Let $r_2(n)$ denote the number of representations of $n$ as $x^2 + y^2$
(with order and signs).  Then
\[
  r_2(n) = 4\!\sum_{d \mid n} \chi(d),
  \qquad \text{where } \chi(d) = \begin{cases}
    +1 & \text{if } d \equiv 1 \pmod{4},\\
    -1 & \text{if } d \equiv 3 \pmod{4},\\
     0  & \text{if } 2 \mid d.
  \end{cases}
\]
Here $\chi$ is the non-principal character modulo~$4$.
\end{proposition}

% ------------------------------------------------------------
\section{Lagrange's Four-Square Theorem}\label{sec:four-sq}
\index{Lagrange's four-square theorem}

Fermat's theorem handles primes $\equiv 1\pmod{4}$.  What about all positive
integers?  Lagrange proved the following remarkable result in 1770.

\begin{theorem}[Lagrange's four-square theorem]\label{thm:lagrange-four-sq}
Every positive integer is a sum of four squares of non-negative integers.
That is, for every $n \ge 1$ there exist $a,b,c,d \in \Z_{\ge 0}$ with
\[
  n = a^2 + b^2 + c^2 + d^2.
\]
\end{theorem}

\begin{proof}[Proof sketch]
The proof follows the same pattern as for two squares.

\textit{Step~1 (Euler's four-square identity).}\index{Euler's four-square identity}
The product of two sums of four squares is itself a sum of four squares.
This follows from the multiplicativity of the norm in the quaternions
$\mathbb{H}$\index{quaternions}: $N(q_1 q_2) = N(q_1)N(q_2)$ where
$N(a + bi + cj + dk) = a^2+b^2+c^2+d^2$.  It therefore suffices to prove the
theorem for primes.

\textit{Step~2.}
For any prime $p$, there exist $a,b$ with $0 \le a,b \le (p-1)/2$
and $a^2 + b^2 + 1 \equiv 0 \pmod{p}$.  This is proved by a pigeonhole
argument on the sets $\{a^2 : 0 \le a \le (p-1)/2\}$ and
$\{-1-b^2 : 0 \le b \le (p-1)/2\}$, each of size $(p+1)/2$.

\textit{Step~3 (Descent).}
Starting from $mp = a^2+b^2+c^2+d^2$ with $1 \le m < p$, a descent
argument analogous to that in Theorem~\ref{thm:fermat-two-sq} reduces $m$
to~$1$.
\end{proof}

\begin{example}\label{ex:four-sq}
\leavevmode
\begin{itemize}
  \item $7 = 4 + 1 + 1 + 1 = 2^2 + 1^2 + 1^2 + 1^2$.
  \item $15 = 9 + 4 + 1 + 1 = 3^2 + 2^2 + 1^2 + 1^2$.
  \item $23 = 9 + 9 + 4 + 1 = 3^2 + 3^2 + 2^2 + 1^2$.
\end{itemize}
\end{example}

\begin{remark}[Three squares]\label{rem:three-sq}
\index{Legendre's three-square theorem}
Legendre (1798) and Gauss showed that a positive integer $n$ is a sum of three
squares if and only if $n$ is \emph{not} of the form $4^a(8b+7)$ for
non-negative integers $a,b$.  Thus $7, 15, 23, 28, 60, \ldots$ require four squares.
\end{remark}

% ------------------------------------------------------------
\section{Waring's Problem and Further Directions}\label{sec:waring}
\index{Waring's problem}

\begin{definition}[Waring's problem]\label{def:waring}
For each positive integer $k$, let $g(k)$ denote the smallest $s$ such that
every positive integer is a sum of at most $s$ perfect $k$-th powers.
Waring's problem asks for the value of $g(k)$ for each~$k$.
\end{definition}

\begin{remark}\label{rem:waring-values}
Known values include $g(1) = 1$, $g(2) = 4$ (Lagrange),
$g(3) = 9$ (Wieferich--Kempner), and $g(4) = 19$ (Balasubramanian--Deshouillers--Dress).
Hilbert proved in 1909 that $g(k)$ is finite for all~$k$.
\index{Hilbert--Waring theorem}
\end{remark}

\begin{remark}[Ternary quadratic forms]\label{rem:ternary}
\index{ternary quadratic form}
Ramanujan (1916) determined all \emph{universal}\index{universal quadratic form}
positive definite diagonal quaternary forms $ax^2+by^2+cz^2+dw^2$ that represent
all positive integers.  The ``$15$-theorem'' of Conway and Schneeberger (1993)
states that a positive definite integer-matrix quadratic form represents
all positive integers if and only if it represents every integer from $1$ to $15$.
\index{15-theorem}
\end{remark}

% ------------------------------------------------------------
\section{Exercises}\label{sec:ch7-exercises}

\begin{exercise}\label{exer:two-sq-100}
Find all integers $n$ with $1 \le n \le 50$ that are sums of two squares.
For each such $n$, give an explicit representation.
\end{exercise}

\begin{exercise}\label{exer:two-sq-product}
Use the Brahmagupta--Fibonacci identity to write $5 \cdot 13 = 65$ as a
sum of two squares in two different ways.
\end{exercise}

\begin{exercise}\label{exer:p-mod-4-necessity}
Show directly (without using the Legendre symbol) that if $p \equiv 3 \pmod{4}$
is prime, then $p$ cannot be written as $x^2 + y^2$.
\end{exercise}

\begin{exercise}\label{exer:gaussian-units}
Determine all units of $\Z[i]$ and prove that $\Z[i]$ is a Euclidean domain
with respect to the norm $N(a+bi) = a^2+b^2$.
\end{exercise}

\begin{exercise}\label{exer:sum-2-sq-n}
\index{sums of two squares!of $n^2$}
Prove that if $n$ is a sum of two squares, then so is $n^2$.  Give two proofs:
one using the Brahmagupta--Fibonacci identity and one using the characterisation
theorem (Theorem~\ref{thm:two-sq-char}).
\end{exercise}

\begin{exercise}\label{exer:four-sq-7}
Show that the integers of the form $4^a(8b+7)$ are not sums of three squares.
\textit{Hint:} Show that if $n \equiv 7 \pmod{8}$ then $n$ is not a sum of
three squares by considering squares modulo~$8$.
\end{exercise}

\begin{exercise}\label{exer:x2+2y2}
\index{$x^2+2y^2$}
Prove that an odd prime $p$ is represented by the form $x^2 + 2y^2$ if and
only if $p \equiv 1$ or $3 \pmod{8}$.
\textit{Hint:} Show that $-2$ is a quadratic residue modulo $p$ if and only if
$p \equiv 1$ or $3\pmod{8}$, then apply a descent argument.
\end{exercise}

\begin{exercise}\label{exer:three-sq-count}
Find all representations of $30$ as a sum of three squares
(where order and signs matter).  How many are there?
\end{exercise}

\begin{exercise}\label{exer:large-form}
Show that for every positive integer $k$, there exist infinitely many integers
that require exactly $k$ terms when written as sums of squares (for $k \le 4$).
\end{exercise}

\begin{exercise}\label{exer:r2-formula}
\index{$r_2(n)$!computation}
Compute $r_2(n)$ for $n = 1, 2, 5, 10, 25$ using the formula in
Proposition~\ref{prop:r2-count}, and verify by direct enumeration.
\end{exercise}

% ------------------------------------------------------------
\section*{Chapter Summary}

\begin{itemize}
  \item A \textbf{binary quadratic form} $ax^2+bxy+cy^2$ has discriminant
        $D = b^2 - 4ac$.  Equivalent forms represent the same integers.
  \item \textbf{Fermat's two-square theorem}: an odd prime $p = x^2+y^2$ if
        and only if $p \equiv 1 \pmod{4}$.  The proof uses the existence of a
        square root of $-1$ modulo $p$ and a descent argument.
  \item A positive integer $n$ is a sum of two squares $\iff$ every prime
        $q \equiv 3\pmod{4}$ dividing $n$ appears to an even power.
  \item \textbf{Lagrange}: every positive integer is a sum of four squares.
  \item The Gaussian integers $\Z[i]$ provide a structural explanation:
        $n = x^2+y^2$ iff $n = N(\alpha)$ for some $\alpha\in\Z[i]$.
  \item \textbf{Waring's problem}: $g(k)$ is finite for every $k$ (Hilbert, 1909).
\end{itemize}

\newpage

% ============================================================
%  CHAPTER 8 — Arithmetic Functions
% ============================================================
\chapter{Arithmetic Functions}\label{ch:arith-functions}
\index{arithmetic function}

\begin{quote}
\textit{``The theory of numbers is the queen of mathematics, and the theory
of forms and functions is the queen of number theory.''}
--- adapted from Gauss
\end{quote}

\noindent
Arithmetic functions assign a value to each positive integer and encode deep
information about the multiplicative structure of $\Z$.  In this chapter we
study the principal arithmetic functions---Euler's totient, divisor functions,
the M\"obius function, and the von Mangoldt function---prove their key
properties, and develop the powerful tool of Dirichlet convolution and
M\"obius inversion.

% ------------------------------------------------------------
\section{The Main Arithmetic Functions}\label{sec:arith-func-def}

\begin{definition}[Arithmetic function]\label{def:arith-function}
An \textbf{arithmetic function}\index{arithmetic function!definition} is any
function $f \colon \Z_{>0} \to \C$.  It is \textbf{multiplicative}
\index{multiplicative function} if $f(1)=1$ and
\[
  f(mn) = f(m)f(n) \quad \text{whenever } \gcd(m,n) = 1.
\]
It is \textbf{completely multiplicative}\index{completely multiplicative function}
if this holds for \emph{all} $m,n$ (without the coprimality condition).
\end{definition}

\subsection{Euler's Totient Function}\label{subsec:euler-phi}
\index{Euler's totient function}

\begin{definition}[Euler's totient]\label{def:euler-phi}
For $n \ge 1$, $\varphi(n)$\index{$\varphi(n)$} is the number of integers
$k$ with $1 \le k \le n$ and $\gcd(k,n)=1$:
\[
  \varphi(n) = \#\{k \in \Z : 1 \le k \le n,\; \gcd(k,n)=1\}
             = n\prod_{p \mid n}\!\left(1 - \frac{1}{p}\right).
\]
\end{definition}

\begin{proposition}[Multiplicativity of $\varphi$]\label{prop:phi-mult}
$\varphi$ is multiplicative.
\end{proposition}

\begin{proof}
Let $\gcd(m,n) = 1$.  By the Chinese Remainder Theorem\index{Chinese Remainder Theorem},
$\Z/mn\Z \cong \Z/m\Z \times \Z/n\Z$, and an element $(a,b)$ is a unit
in the product ring if and only if $a$ is a unit mod~$m$ and $b$ is a unit
mod~$n$.  Hence
\[
  \varphi(mn) = (\Z/mn\Z)^\times = |\!(\Z/m\Z)^\times\!|
               \cdot |\!(\Z/n\Z)^\times\!|
             = \varphi(m)\,\varphi(n). \qedhere
\]
\end{proof}

\begin{proposition}[Formula for prime powers]\label{prop:phi-primepower}
For a prime $p$ and $a \ge 1$,
$\varphi(p^a) = p^a - p^{a-1} = p^{a-1}(p-1)$.
\end{proposition}

\begin{proof}
Among $1, 2, \ldots, p^a$, those \emph{not} coprime to $p^a$ are exactly the
multiples of $p$: there are $p^{a-1}$ of them.  So
$\varphi(p^a) = p^a - p^{a-1}$.
\end{proof}

\subsection{Divisor Functions}\label{subsec:divisor-functions}
\index{divisor function}

\begin{definition}[Divisor functions $\tau$ and $\sigma$]\label{def:tau-sigma}
For $n \ge 1$, define:
\begin{align*}
  \tau(n) &= \sum_{d \mid n} 1 = \text{number of positive divisors of } n,
    \index{$\tau(n)$} \\
  \sigma(n) &= \sum_{d \mid n} d = \text{sum of positive divisors of } n.
    \index{$\sigma(n)$}
\end{align*}
More generally, $\sigma_s(n) = \sum_{d\mid n} d^s$ for $s \in \C$,
\index{$\sigma_s(n)$}
so $\tau(n) = \sigma_0(n)$ and $\sigma(n) = \sigma_1(n)$.
\end{definition}

\begin{proposition}\label{prop:tau-sigma-mult}
The functions $\tau$ and $\sigma$ (and more generally $\sigma_s$ for each~$s$)
are multiplicative.
\end{proposition}

\begin{proof}
Let $\gcd(m,n)=1$.  Every divisor $d$ of $mn$ can be uniquely written as
$d = d_1 d_2$ with $d_1 \mid m$ and $d_2 \mid n$ (and $\gcd(d_1,d_2)=1$).
Therefore
\[
  \sigma_s(mn) = \sum_{d\mid mn} d^s
               = \sum_{d_1\mid m}\sum_{d_2\mid n} (d_1 d_2)^s
               = \biggl(\sum_{d_1\mid m} d_1^s\biggr)
                 \biggl(\sum_{d_2\mid n} d_2^s\biggr)
               = \sigma_s(m)\,\sigma_s(n). \qedhere
\]
\end{proof}

\begin{proposition}[Formulas on prime powers]\label{prop:tau-sigma-pp}
For a prime $p$ and $a \ge 1$:
\[
  \tau(p^a) = a+1, \qquad
  \sigma(p^a) = \frac{p^{a+1}-1}{p-1}.
\]
\end{proposition}

\begin{proof}
The divisors of $p^a$ are $1, p, p^2, \ldots, p^a$, giving $a+1$ divisors
and sum $1 + p + \cdots + p^a = (p^{a+1}-1)/(p-1)$.
\end{proof}

\subsection{The M\"obius Function}\label{subsec:mobius}
\index{M\"obius function}

\begin{definition}[M\"obius function]\label{def:mobius}
The M\"obius function $\mu \colon \Z_{>0} \to \{-1,0,1\}$
\index{$\mu(n)$} is defined by
\[
  \mu(n) = \begin{cases}
    1       & \text{if } n = 1, \\
    (-1)^k  & \text{if } n = p_1 p_2 \cdots p_k \text{ for distinct primes } p_i, \\
    0       & \text{if } p^2 \mid n \text{ for some prime } p.
  \end{cases}
\]
\end{definition}

\begin{proposition}\label{prop:mu-mult}
$\mu$ is multiplicative.
\end{proposition}

\begin{proof}
Let $\gcd(m,n)=1$.  If $p^2 \mid m$ or $p^2 \mid n$, then $p^2 \mid mn$ and
both sides of $\mu(mn) = \mu(m)\mu(n)$ are zero.  Otherwise
$m = p_1\cdots p_j$ and $n = q_1 \cdots q_k$ with all primes distinct
(since $\gcd(m,n)=1$), so
$\mu(mn) = (-1)^{j+k} = (-1)^j(-1)^k = \mu(m)\mu(n)$.
\end{proof}

\begin{proposition}\label{prop:mu-sum}
\index{M\"obius function!sum over divisors}
For all $n \ge 1$,
\[
  \sum_{d \mid n} \mu(d) = \begin{cases} 1 & \text{if } n = 1,\\ 0 & \text{if } n > 1.\end{cases}
\]
In other words, $\mu \ast \mathbf{1} = \varepsilon$, where
$\mathbf{1}(n)=1$ for all $n$ and $\varepsilon(n) = [n=1]$.
\end{proposition}

\begin{proof}
The case $n=1$ is immediate: $\mu(1) = 1$.

For $n > 1$, write $n = p_1^{a_1}\cdots p_r^{a_r}$.  Since $\mu(d)=0$ whenever
$d$ has a squared prime factor, only the squarefree divisors contribute:
\[
  \sum_{d\mid n} \mu(d)
  = \sum_{S \subseteq \{p_1,\ldots,p_r\}} \mu\!\Bigl(\prod_{p\in S} p\Bigr)
  = \sum_{k=0}^{r} \binom{r}{k}(-1)^k
  = (1-1)^r = 0. \qedhere
\]
\end{proof}

\subsection{The von Mangoldt Function}\label{subsec:von-mangoldt}
\index{von Mangoldt function}

\begin{definition}[Von Mangoldt function]\label{def:von-mangoldt}
The \textbf{von Mangoldt function}\index{$\Lambda(n)$} is
\[
  \Lambda(n) = \begin{cases}
    \log p & \text{if } n = p^k \text{ for some prime } p \text{ and } k \ge 1,\\
    0      & \text{otherwise.}
  \end{cases}
\]
\end{definition}

\begin{proposition}\label{prop:mangoldt-sum}
For all $n \ge 1$, $\displaystyle\sum_{d \mid n} \Lambda(d) = \log n$.
\end{proposition}

\begin{proof}
Write $n = p_1^{a_1}\cdots p_r^{a_r}$.  The divisors $d$ of $n$ with
$\Lambda(d)\neq 0$ are $p_i^j$ for $1 \le j \le a_i$ and $1 \le i \le r$.
Therefore
\[
  \sum_{d\mid n} \Lambda(d) = \sum_{i=1}^{r} \sum_{j=1}^{a_i} \log p_i
  = \sum_{i=1}^{r} a_i \log p_i = \log\!\Bigl(\prod p_i^{a_i}\Bigr)
  = \log n. \qedhere
\]
\end{proof}

\begin{remark}\label{rem:mangoldt-not-mult}
The von Mangoldt function is \emph{not} multiplicative: $\Lambda(1) = 0 \neq 1$.
It plays a central role in analytic number theory, particularly in the
prime number theorem\index{prime number theorem} through the Chebyshev functions
$\psi(x) = \sum_{n\le x}\Lambda(n)$.
\end{remark}

% ------------------------------------------------------------
\section{The Identity $\sum_{d \mid n} \varphi(d) = n$}\label{sec:phi-sum}
\index{Euler's totient function!divisor sum identity}

\begin{theorem}\label{thm:phi-divisor-sum}
For every positive integer $n$,
\[
  \sum_{d \mid n} \varphi(d) = n.
\]
\end{theorem}

\begin{proof}
We give two proofs.

\textbf{Proof~1 (Counting).}
Partition the set $\{1, 2, \ldots, n\}$ according to the value of
$\gcd(k, n)$.  For each divisor $d$ of $n$, the integers $k$ with
$1 \le k \le n$ and $\gcd(k,n) = d$ are exactly the integers of the form
$k = d\ell$ where $1 \le \ell \le n/d$ and $\gcd(\ell, n/d) = 1$.
There are $\varphi(n/d)$ such integers.  Since every $k \in \{1,\ldots,n\}$
has $\gcd(k,n) = d$ for exactly one divisor $d$,
\[
  n = \sum_{d \mid n} \varphi\!\left(\frac{n}{d}\right)
    = \sum_{d \mid n} \varphi(d),
\]
where the last equality holds because as $d$ runs over divisors of $n$,
so does $n/d$.

\medskip
\textbf{Proof~2 (Multiplicative functions).}
Define $F(n) = \sum_{d\mid n}\varphi(d)$.  Since $\varphi$ is multiplicative,
$F$ is also multiplicative (a general fact about Dirichlet convolution,
Proposition~\ref{prop:mult-conv}).
For a prime power $p^a$:
\[
  F(p^a) = \sum_{j=0}^{a}\varphi(p^j)
         = 1 + \sum_{j=1}^{a} p^{j-1}(p-1)
         = 1 + (p-1)\cdot\frac{p^a - 1}{p-1}
         = p^a.
\]
The identity function $\mathrm{id}(n) = n$ is also multiplicative and
satisfies $\mathrm{id}(p^a)=p^a$.  Since two multiplicative functions that
agree on prime powers agree everywhere, $F = \mathrm{id}$.
\end{proof}

% ------------------------------------------------------------
\section{Dirichlet Convolution}\label{sec:dirichlet-conv}
\index{Dirichlet convolution}

\begin{definition}[Dirichlet convolution]\label{def:dirichlet-conv}
Given arithmetic functions $f$ and $g$, their \textbf{Dirichlet convolution}
$f \ast g$ is defined by
\[
  (f \ast g)(n) = \sum_{d \mid n} f(d)\,g\!\left(\frac{n}{d}\right)
               = \sum_{\substack{ab = n \\ a,b \ge 1}} f(a)\,g(b).
\]
\end{definition}

\begin{proposition}[Ring structure]\label{prop:conv-ring}
\index{Dirichlet convolution!ring structure}
The set of arithmetic functions with pointwise addition and Dirichlet
convolution forms a commutative ring with identity.
\begin{enumerate}[label=(\roman*)]
  \item \textbf{Commutativity:} $f \ast g = g \ast f$.
  \item \textbf{Associativity:} $(f \ast g) \ast h = f \ast (g \ast h)$.
  \item \textbf{Identity:} $f \ast \varepsilon = f$, where
        $\varepsilon(n) = [n=1]$ is the Dirichlet identity.
        \index{Dirichlet identity}
  \item \textbf{Distributivity:} $f \ast (g + h) = f\ast g + f\ast h$.
\end{enumerate}
\end{proposition}

\begin{proof}
(i) follows from the bijection $(a,b) \leftrightarrow (b,a)$ on pairs with $ab=n$.

(ii) Both sides equal $\sum_{abc=n} f(a)\,g(b)\,h(c)$.

(iii) $(f\ast\varepsilon)(n) = \sum_{d\mid n} f(d)\,\varepsilon(n/d) = f(n)\cdot 1 = f(n)$
since $\varepsilon(n/d) = 0$ unless $d=n$.

(iv) Direct computation.
\end{proof}

\begin{notation}\label{not:standard-functions}
We use the following standard arithmetic functions:
\begin{align*}
  \varepsilon(n) &= [n=1], & \mathbf{1}(n) &= 1, &
  \mathrm{id}(n) &= n, & \mathrm{id}_s(n) &= n^s.
\end{align*}
\end{notation}

With this notation, the identities we have proved become elegant convolution
equations:
\begin{align}
  \tau &= \mathbf{1} \ast \mathbf{1},    \label{eq:tau-conv} \\
  \sigma &= \mathrm{id} \ast \mathbf{1}, \label{eq:sigma-conv} \\
  \mathrm{id} &= \varphi \ast \mathbf{1}, \label{eq:phi-conv} \\
  \varepsilon &= \mu \ast \mathbf{1},     \label{eq:mu-conv} \\
  \Lambda &= \mu \ast (-\log).            \label{eq:mangoldt-conv}
\end{align}

\begin{proposition}[Multiplicativity preserved by convolution]\label{prop:mult-conv}
\index{multiplicative function!convolution}
If $f$ and $g$ are multiplicative, then so is $f \ast g$.
\end{proposition}

\begin{proof}
Let $\gcd(m,n)=1$.  Every divisor $d$ of $mn$ factors uniquely as $d=d_1 d_2$
with $d_1 \mid m$, $d_2 \mid n$, and $\gcd(d_1,d_2)=1$.
Then
\begin{align*}
  (f\ast g)(mn) &= \sum_{d\mid mn} f(d)\,g\!\left(\frac{mn}{d}\right)
  = \sum_{d_1\mid m}\sum_{d_2\mid n}
    f(d_1 d_2)\,g\!\left(\frac{m}{d_1}\cdot\frac{n}{d_2}\right) \\
  &= \sum_{d_1\mid m}\sum_{d_2\mid n}
    f(d_1)f(d_2)\,g\!\left(\frac{m}{d_1}\right)g\!\left(\frac{n}{d_2}\right) \\
  &= \biggl(\sum_{d_1\mid m}f(d_1)\,g\!\left(\frac{m}{d_1}\right)\!\biggr)
     \biggl(\sum_{d_2\mid n}f(d_2)\,g\!\left(\frac{n}{d_2}\right)\!\biggr)
  = (f\ast g)(m)\,(f\ast g)(n).
\end{align*}
Also $(f\ast g)(1) = f(1)g(1) = 1$.
\end{proof}

\begin{theorem}[Dirichlet inverse]\label{thm:dirichlet-inverse}
\index{Dirichlet inverse}
An arithmetic function $f$ has a Dirichlet inverse (i.e., there exists $g$
with $f \ast g = \varepsilon$) if and only if $f(1) \neq 0$.  In this case,
$g$ is unique and is given recursively by $g(1) = 1/f(1)$ and, for $n > 1$,
\[
  g(n) = -\frac{1}{f(1)} \sum_{\substack{d \mid n \\ d < n}}
         f\!\left(\frac{n}{d}\right) g(d).
\]
Moreover, if $f$ is multiplicative, then so is $f^{-1}$.
\end{theorem}

\begin{proof}
From $(f\ast g)(1) = f(1)g(1) = \varepsilon(1) = 1$, we need $g(1) = 1/f(1)$,
which requires $f(1)\neq 0$.  For $n > 1$,
$(f\ast g)(n) = 0$ gives
$f(1)g(n) + \sum_{d\mid n,\, d<n} f(n/d)\,g(d) = 0$,
yielding the recursion.  Uniqueness follows since $g$ is determined
recursively.

For multiplicativity of $f^{-1}$ when $f$ is multiplicative: since
$f \ast f^{-1} = \varepsilon$ is multiplicative, one shows by induction
on $mn$ (for $\gcd(m,n)=1$) that $f^{-1}(mn) = f^{-1}(m)f^{-1}(n)$.
\end{proof}

% Dirichlet convolution visualisation
\begin{center}
\begin{tikzpicture}[
    node distance=1.8cm,
    funcbox/.style={draw, rounded corners, fill=blue!8,
      minimum width=1.5cm, minimum height=0.8cm, font=\small},
    convbox/.style={draw, rounded corners, fill=red!8,
      minimum width=2cm, minimum height=0.8cm, font=\small},
    arr/.style={->,>=Stealth,thick}
  ]
  \node[funcbox] (mu)    {$\mu$};
  \node[funcbox, right=of mu] (one)   {$\mathbf{1}$};
  \node[convbox, below=1cm of $(mu)!0.5!(one)$] (eps) {$\varepsilon$};
  \draw[arr] (mu) -- node[above]{\scriptsize $\ast$} (one);
  \draw[arr,dashed] ($(mu)!0.5!(one)$) -- node[right]{\scriptsize $=$} (eps);

  \node[funcbox, right=3cm of one] (phi) {$\varphi$};
  \node[funcbox, right=of phi] (one2)  {$\mathbf{1}$};
  \node[convbox, below=1cm of $(phi)!0.5!(one2)$] (id) {$\mathrm{id}$};
  \draw[arr] (phi) -- node[above]{\scriptsize $\ast$} (one2);
  \draw[arr,dashed] ($(phi)!0.5!(one2)$) -- node[right]{\scriptsize $=$} (id);
\end{tikzpicture}
\captionof{figure}{Dirichlet convolution identities:
$\mu \ast \mathbf{1} = \varepsilon$ and
$\varphi \ast \mathbf{1} = \mathrm{id}$.}\label{fig:dirichlet-conv}
\index{Dirichlet convolution!visualisation}
\end{center}

% ------------------------------------------------------------
\section{M\"obius Inversion}\label{sec:mobius-inversion}
\index{M\"obius inversion}

\begin{theorem}[M\"obius inversion formula]\label{thm:mobius-inversion}
Let $f$ and $g$ be arithmetic functions.  Then
\[
  g(n) = \sum_{d \mid n} f(d) \quad \text{for all } n \ge 1
\]
if and only if
\[
  f(n) = \sum_{d \mid n} \mu(d)\,g\!\left(\frac{n}{d}\right)
       = \sum_{d \mid n} \mu\!\left(\frac{n}{d}\right) g(d)
  \quad \text{for all } n \ge 1.
\]
In convolution notation: $g = f \ast \mathbf{1}$ if and only if
$f = g \ast \mu = \mu \ast g$.
\end{theorem}

\begin{proof}
\textbf{Forward direction.}
Suppose $g = f \ast \mathbf{1}$.  Then
\[
  (g \ast \mu)(n) = ((f \ast \mathbf{1})\ast \mu)(n)
                  = (f \ast (\mathbf{1} \ast \mu))(n)
                  = (f \ast \varepsilon)(n) = f(n),
\]
using associativity (Proposition~\ref{prop:conv-ring}(ii)) and the
identity $\mathbf{1}\ast\mu = \varepsilon$ (Proposition~\ref{prop:mu-sum}).

\textbf{Backward direction.}
Conversely, if $f = g \ast \mu$, then
\[
  (f \ast \mathbf{1})(n) = ((g \ast \mu)\ast \mathbf{1})(n)
                         = (g \ast (\mu \ast \mathbf{1}))(n)
                         = (g \ast \varepsilon)(n) = g(n).
\]
Both directions are instances of the same algebraic fact:
$\mu = \mathbf{1}^{-1}$ in the Dirichlet ring.
\end{proof}

\begin{corollary}[M\"obius inversion for $\varphi$]\label{cor:phi-mobius}
\index{Euler's totient function!M\"obius inversion}
From $\mathrm{id} = \varphi \ast \mathbf{1}$, M\"obius inversion gives
\[
  \varphi(n) = \sum_{d \mid n} \mu\!\left(\frac{n}{d}\right) d
             = n \sum_{d \mid n} \frac{\mu(d)}{d}
             = n \prod_{p \mid n}\!\left(1 - \frac{1}{p}\right).
\]
\end{corollary}

\begin{proof}
By M\"obius inversion, $\varphi = \mathrm{id} \ast \mu$.  Thus
\[
  \varphi(n) = \sum_{d\mid n} \mu(d)\,\frac{n}{d}
             = n \sum_{d\mid n}\frac{\mu(d)}{d}.
\]
For $n = p_1^{a_1}\cdots p_r^{a_r}$, only squarefree divisors $d$
(products of subsets of $\{p_1,\ldots,p_r\}$) give $\mu(d)\neq 0$:
\[
  \sum_{d\mid n}\frac{\mu(d)}{d}
  = \prod_{i=1}^{r}\!\left(1 - \frac{1}{p_i}\right). \qedhere
\]
\end{proof}

\begin{example}[Application: counting primitive $n$-th roots of unity]
\label{ex:primitive-roots-unity}
\index{roots of unity!primitive}
The cyclotomic polynomial $\Phi_n(x) = \prod_{\substack{1\le k \le n \\
\gcd(k,n)=1}}(x - e^{2\pi i k/n})$ has degree $\varphi(n)$.  Since
$x^n - 1 = \prod_{d\mid n}\Phi_d(x)$, taking degrees gives
$n = \sum_{d\mid n}\varphi(d)$, recovering Theorem~\ref{thm:phi-divisor-sum}.
\end{example}

\begin{example}[Counting square-free integers]\label{ex:squarefree}
\index{squarefree integer}
Let $Q(x) = \#\{n \le x : n \text{ is squarefree}\}$.  Using the identity
$[\text{$n$ squarefree}] = \sum_{d^2\mid n}\mu(d)$, one shows
$Q(x) = \frac{6}{\pi^2}\,x + O(\sqrt{x})$.
\end{example}

% ------------------------------------------------------------
\section{Dirichlet Series and Euler Products}\label{sec:dirichlet-series}
\index{Dirichlet series}

\begin{definition}[Dirichlet series]\label{def:dirichlet-series}
The \textbf{Dirichlet series} associated to an arithmetic function $f$ is the
formal series
\[
  F(s) = \sum_{n=1}^{\infty} \frac{f(n)}{n^s}.
\]
\end{definition}

\begin{proposition}\label{prop:dir-series-conv}
If $F(s) = \sum f(n)/n^s$ and $G(s) = \sum g(n)/n^s$, then (in the region
of absolute convergence)
\[
  F(s)\,G(s) = \sum_{n=1}^{\infty} \frac{(f\ast g)(n)}{n^s}.
\]
That is, Dirichlet convolution corresponds to multiplication of Dirichlet series.
\end{proposition}

\begin{proof}
\[
  F(s)G(s) = \sum_{m=1}^{\infty}\sum_{n=1}^{\infty}
  \frac{f(m)\,g(n)}{(mn)^s}
  = \sum_{N=1}^{\infty}\frac{1}{N^s}\sum_{mn=N}f(m)\,g(n)
  = \sum_{N=1}^{\infty}\frac{(f\ast g)(N)}{N^s}. \qedhere
\]
\end{proof}

\begin{theorem}[Euler product]\label{thm:euler-product}
\index{Euler product}
If $f$ is multiplicative, then (in the region of absolute convergence)
\[
  \sum_{n=1}^{\infty} \frac{f(n)}{n^s}
  = \prod_{p\text{ prime}} \left(1 + \frac{f(p)}{p^s}
    + \frac{f(p^2)}{p^{2s}} + \cdots\right).
\]
\end{theorem}

\begin{proof}
By the fundamental theorem of arithmetic, every $n$ factors uniquely as
$n = p_1^{a_1}\cdots p_r^{a_r}$.  The multiplicativity of $f$ gives
$f(n) = f(p_1^{a_1})\cdots f(p_r^{a_r})$, and $n^{-s} = p_1^{-a_1 s}\cdots
p_r^{-a_r s}$.  Expanding the product over all primes and collecting terms
recovers the Dirichlet series.
\end{proof}

\begin{example}[Standard Euler products]\label{ex:euler-products}
\index{Riemann zeta function}
\leavevmode
\begin{enumerate}[label=(\alph*)]
  \item $\displaystyle\zeta(s) = \sum_{n=1}^\infty \frac{1}{n^s}
         = \prod_p \frac{1}{1-p^{-s}}$ \quad (Riemann zeta function,
         $\operatorname{Re}(s) > 1$).
  \item $\displaystyle\frac{1}{\zeta(s)} = \sum_{n=1}^\infty \frac{\mu(n)}{n^s}
         = \prod_p (1 - p^{-s})$.
  \item $\displaystyle\frac{\zeta(s-1)}{\zeta(s)}
         = \sum_{n=1}^\infty \frac{\varphi(n)}{n^s}$
         \quad (from $\varphi = \mathrm{id}\ast\mu$).
  \item $\displaystyle -\frac{\zeta'(s)}{\zeta(s)}
         = \sum_{n=1}^\infty \frac{\Lambda(n)}{n^s}$.
\end{enumerate}
\end{example}

% ------------------------------------------------------------
\section{Perfect Numbers and $\sigma(n)$}\label{sec:perfect-numbers}
\index{perfect number}

\begin{definition}[Perfect number]\label{def:perfect-number}
A positive integer $n$ is \textbf{perfect}\index{perfect number!definition}
if $\sigma(n) = 2n$, i.e., $n$ equals the sum of its proper divisors.
\end{definition}

\begin{theorem}[Euler--Euclid theorem]\label{thm:euler-euclid-perfect}
\index{Euler--Euclid theorem}
An even integer $n$ is perfect if and only if $n = 2^{p-1}(2^p - 1)$ where
$2^p - 1$ is a Mersenne prime\index{Mersenne prime}.
\end{theorem}

\begin{proof}
\textbf{Sufficiency} (Euclid).
Let $M = 2^p - 1$ be prime and $n = 2^{p-1}M$.  Since $\gcd(2^{p-1},M) = 1$,
\[
  \sigma(n) = \sigma(2^{p-1})\,\sigma(M)
            = (2^p - 1)(M + 1)
            = M \cdot 2^p = 2n.
\]

\textbf{Necessity} (Euler).
Let $n = 2^{a} m$ with $a \ge 1$ and $m$ odd.  Then
$\sigma(n) = \sigma(2^a)\sigma(m) = (2^{a+1}-1)\sigma(m)$.
Setting $\sigma(n) = 2n = 2^{a+1}m$ gives
\[
  (2^{a+1}-1)\,\sigma(m) = 2^{a+1}\,m.
\]
Since $\gcd(2^{a+1}-1, 2^{a+1}) = 1$, we need $(2^{a+1}-1) \mid m$.
Write $m = (2^{a+1}-1)\,q$.  Then
$\sigma(m) = 2^{a+1}\,q$.  But $\sigma(m) \ge m + 1 = (2^{a+1}-1)q + 1$
(if $m > 1$) and also $\sigma(m) \ge m + q + 1$ (if $q > 1$ and $q \neq m$).

If $q = 1$, then $m = 2^{a+1}-1$ and $\sigma(m) = 2^{a+1}$, so
$\sigma(m) = m + 1$, forcing $m$ to be prime.
Thus $n = 2^a(2^{a+1}-1)$ with $2^{a+1}-1$ prime, as desired.

If $q > 1$, then $q$ is a proper divisor of $m$ distinct from $1$ and $m$,
so $\sigma(m) \ge 1 + q + m = 1 + q + (2^{a+1}-1)q > 2^{a+1}q = \sigma(m)$,
a contradiction.
\end{proof}

\begin{remark}\label{rem:odd-perfect}
\index{perfect number!odd}
It is unknown whether any odd perfect numbers exist.  It has been verified
that none exist below $10^{2200}$ (Ochem--Rao, 2012).  Any odd perfect number
must have at least $10$ distinct prime factors and exceed $10^{1500}$.
\end{remark}

% Divisor lattice
\begin{center}
\begin{tikzpicture}[scale=1, every node/.style={font=\small}]
  % Divisor lattice of 12
  \node (1) at (0,0) {$1$};
  \node (2) at (-2,1.3) {$2$};
  \node (3) at (0,1.3) {$3$};
  \node (4) at (-2,2.6) {$4$};
  \node (6) at (0,2.6) {$6$};
  \node (12) at (-1,3.9) {$12$};
  \draw (1)--(2); \draw (1)--(3);
  \draw (2)--(4); \draw (2)--(6); \draw (3)--(6);
  \draw (4)--(12); \draw (6)--(12);
  \node at (-1,-0.7) {Divisor lattice of $12$};

  % Divisor lattice of 30
  \node (a1) at (5,0) {$1$};
  \node (a2) at (3.5,1.1) {$2$};
  \node (a3) at (5,1.1) {$3$};
  \node (a5) at (6.5,1.1) {$5$};
  \node (a6) at (3.5,2.3) {$6$};
  \node (a10) at (5,2.3) {$10$};
  \node (a15) at (6.5,2.3) {$15$};
  \node (a30) at (5,3.5) {$30$};
  \draw (a1)--(a2); \draw (a1)--(a3); \draw (a1)--(a5);
  \draw (a2)--(a6); \draw (a2)--(a10);
  \draw (a3)--(a6); \draw (a3)--(a15);
  \draw (a5)--(a10); \draw (a5)--(a15);
  \draw (a6)--(a30); \draw (a10)--(a30); \draw (a15)--(a30);
  \node at (5,-0.7) {Divisor lattice of $30$};
\end{tikzpicture}
\captionof{figure}{Divisor lattices (Hasse diagrams) of $12 = 2^2\cdot 3$
and $30 = 2\cdot 3\cdot 5$.}\label{fig:divisor-lattice}
\index{divisor lattice}
\end{center}

% ------------------------------------------------------------
\section{Summary Table of Arithmetic Functions}\label{sec:arith-table}

\begin{center}
\begin{tabular}{@{}lllll@{}}
\toprule
\textbf{Function} & \textbf{Definition} & \textbf{Mult.?}
  & \textbf{At $p^a$} & \textbf{Dirichlet series} \\
\midrule
$\varphi(n)$ & $\#\{k \le n : \gcd(k,n)=1\}$
  & Yes & $p^{a-1}(p-1)$ & $\zeta(s-1)/\zeta(s)$ \\[3pt]
$\tau(n)$    & $\#\{d : d\mid n\}$
  & Yes & $a+1$ & $\zeta(s)^2$ \\[3pt]
$\sigma(n)$  & $\sum_{d\mid n}d$
  & Yes & $\frac{p^{a+1}-1}{p-1}$ & $\zeta(s)\zeta(s-1)$ \\[3pt]
$\mu(n)$     & $(-1)^k$ if $n=p_1\cdots p_k$
  & Yes & $\begin{cases}-1&a=1\\0&a\ge 2\end{cases}$ & $1/\zeta(s)$ \\[3pt]
$\Lambda(n)$ & $\log p$ if $n=p^k$
  & No  & $\log p$ & $-\zeta'(s)/\zeta(s)$ \\
\bottomrule
\end{tabular}
\captionof{table}{Key arithmetic functions and their properties.}\label{tab:arith-funcs}
\end{center}

% ------------------------------------------------------------
\section{Exercises}\label{sec:ch8-exercises}

\begin{exercise}\label{exer:phi-values}
Compute $\varphi(n)$ for $n = 1, 2, \ldots, 20$.  Verify that $\varphi(n)$
is even for all $n \ge 3$.
\end{exercise}

\begin{exercise}\label{exer:tau-sigma-mult}
\index{$\tau(n)$!computation}
Compute $\tau(n)$ and $\sigma(n)$ for $n = 360$.  Is $360$ a
\emph{highly composite number}\index{highly composite number}?
\end{exercise}

\begin{exercise}\label{exer:phi-divisor-sum-verify}
Verify the identity $\sum_{d\mid n}\varphi(d) = n$ for $n = 12$ by computing
each term explicitly.
\end{exercise}

\begin{exercise}\label{exer:mobius-inversion-apply}
\index{M\"obius inversion!application}
Let $f(n) = n^2$ and $g(n) = \sum_{d\mid n}f(d) = \sum_{d\mid n}d^2$.
Use M\"obius inversion to express $f$ in terms of $g$ and $\mu$, then
verify for $n = 6$.
\end{exercise}

\begin{exercise}\label{exer:mu-squared-sum}
\index{M\"obius function!squared}
Show that $\sum_{d\mid n}\abs{\mu(d)} = 2^{\omega(n)}$, where $\omega(n)$
is the number of distinct prime factors of $n$.
\end{exercise}

\begin{exercise}\label{exer:conv-associativity}
Prove directly (without using Dirichlet series) that Dirichlet convolution
is associative.
\end{exercise}

\begin{exercise}\label{exer:mangoldt-inversion}
\index{von Mangoldt function!M\"obius inversion}
Use M\"obius inversion on $\sum_{d\mid n}\Lambda(d) = \log n$ to derive the
formula $\Lambda(n) = -\sum_{d\mid n}\mu(d)\log d$.
\textit{Hint:} Apply M\"obius inversion to $f = \Lambda$ and $g = \log$.
\end{exercise}

\begin{exercise}\label{exer:sigma-bound}
\index{$\sigma(n)$!bound}
Prove that $\sigma(n) < 2n\log n$ for all $n \ge 2$.
\textit{Hint:} Show $\sigma(n)/n = \sum_{d\mid n}1/d \le \sum_{k=1}^{n}1/k$.
\end{exercise}

\begin{exercise}\label{exer:perfect-number-form}
Show that $2^{p-1}(2^p-1)$ is even and that if $2^p - 1$ is composite then
$2^{p-1}(2^p-1)$ is not perfect.
\end{exercise}

\begin{exercise}\label{exer:dirichlet-inverse-mu}
Compute the Dirichlet inverse of $\mathrm{id}(n) = n$.  That is, find the
arithmetic function $h$ such that $\mathrm{id} \ast h = \varepsilon$.
Express $h$ in terms of $\mu$ and verify for small values of $n$.
\end{exercise}

\begin{exercise}\label{exer:euler-product-derive}
Starting from $\zeta(s) = \prod_p(1-p^{-s})^{-1}$, derive the Euler product
for $\zeta(s)^2$ and identify the corresponding arithmetic function.
\end{exercise}

\begin{exercise}\label{exer:multiplicative-determined}
\index{multiplicative function!determination}
Prove that a multiplicative function is completely determined by its values
on prime powers.  That is, if $f$ and $g$ are multiplicative and
$f(p^a) = g(p^a)$ for all primes $p$ and $a \ge 1$, then $f = g$.
\end{exercise}

% ------------------------------------------------------------
\section*{Chapter Summary}

\begin{itemize}
  \item The main arithmetic functions are $\varphi(n)$, $\tau(n)$, $\sigma(n)$,
        $\mu(n)$, and $\Lambda(n)$; all except $\Lambda$ are multiplicative.
  \item Key identity: $\sum_{d\mid n}\varphi(d) = n$, proved by counting or
        by multiplicative function arguments.
  \item \textbf{Dirichlet convolution} $(f\ast g)(n) = \sum_{d\mid n}f(d)\,g(n/d)$
        makes the arithmetic functions into a commutative ring with identity
        $\varepsilon$.  Multiplicativity is preserved under convolution.
  \item \textbf{M\"obius inversion}: $g = f\ast\mathbf{1}$ iff
        $f = g\ast\mu$.  This powerful tool converts summatory relations
        into explicit formulas.
  \item \textbf{Dirichlet series} $\sum f(n)/n^s$ convert convolution into
        multiplication.  When $f$ is multiplicative, the series has an
        \textbf{Euler product} over primes.
  \item An even number is \textbf{perfect} iff it has the form
        $2^{p-1}(2^p-1)$ with $2^p-1$ a Mersenne prime (Euler--Euclid).
\end{itemize}

% === END OF CHUNK ch7-8 ===
% === START OF CHUNK ch9-10+end ===

%%%%%%%%%%%%%%%%%%%%%%%%%%%%%%%%%%%%%%%%%%%%%%%%%%%%%%%%%%%%%
%% CHAPTER 9: Introduction to p-adic Numbers
%%%%%%%%%%%%%%%%%%%%%%%%%%%%%%%%%%%%%%%%%%%%%%%%%%%%%%%%%%%%%
\chapter{Introduction to \texorpdfstring{$p$}{p}-adic Numbers}
\label{ch:p-adic}
\index{p-adic numbers@$p$-adic numbers|(}

\begin{quote}
\emph{``The $p$-adic numbers are just as `real' as the real numbers.''}
--- attributed to various number theorists
\end{quote}

\noindent
Throughout this course, we have worked within the familiar world of the integers
and rational numbers, occasionally invoking properties of the real or complex
numbers. In this chapter we introduce a fundamentally different way to complete
the rationals: the \emph{$p$-adic numbers}. Whereas the real numbers arise from
completing $\Q$ with respect to the usual absolute value, the $p$-adic numbers
$\Q_p$ arise from a completion governed entirely by divisibility by a prime~$p$.
Far from being a curiosity, $p$-adic methods have become indispensable in modern
number theory, from solving Diophantine equations to the proof of Fermat's Last
Theorem.

%------------------------------------------------------------------
\section{Historical Background: Hensel and \texorpdfstring{$p$}{p}-adic Analysis}
\label{sec:p-adic-history}
\index{Hensel, Kurt}
%------------------------------------------------------------------

Kurt Hensel\index{Hensel, Kurt} introduced the $p$-adic numbers in 1897, motivated
by an analogy between the ring of integers $\Z$ and the ring of polynomials
$k[x]$ over a field~$k$. Just as a polynomial can be expanded as a power series
in $(x - a)$ around a point~$a$, Hensel proposed that every integer could be
expanded as a series in powers of a prime~$p$. This bold analogy proved
extraordinarily fruitful: it opened the door to applying techniques of analysis
--- convergence, completeness, and continuity --- to purely arithmetic questions.

The idea gained further momentum through the work of Hasse\index{Hasse, Helmut},
who in the 1920s formulated the celebrated \emph{local-global principle}, stating
that certain equations have rational solutions if and only if they have solutions
in every $p$-adic field $\Q_p$ and in $\R$. Today, $p$-adic analysis underpins
vast swathes of algebraic number theory, arithmetic geometry, and the Langlands
programme.

%------------------------------------------------------------------
\section{The \texorpdfstring{$p$}{p}-adic Valuation}
\label{sec:p-adic-valuation}
\index{p-adic valuation@$p$-adic valuation}
%------------------------------------------------------------------

Fix a prime number~$p$.

\begin{definition}[$p$-adic valuation]
\label{def:p-adic-valuation}
\index{valuation!p-adic@$p$-adic}
For a nonzero integer $n \in \Z$, the \emph{$p$-adic valuation} of~$n$,
denoted $v_p(n)$, is the largest exponent $k \geq 0$ such that $p^k \mid n$.
We extend this to $\Q^*$ by setting
\[
  v_p\!\left(\frac{a}{b}\right) = v_p(a) - v_p(b),
\]
and by convention $v_p(0) = +\infty$.
\end{definition}

\begin{example}[Computing $v_p$]
\label{ex:vp-computation}
We have $v_2(12) = 2$ since $12 = 2^2 \cdot 3$, and $v_3(12) = 1$, while
$v_5(12)=0$. For a rational number, $v_3(45/98) = v_3(45) - v_3(98) = 2 - 0 = 2$.
\end{example}

\begin{proposition}[Properties of the $p$-adic valuation]
\label{prop:vp-properties}
\index{p-adic valuation@$p$-adic valuation!properties}
For all $x,y \in \Q^*$:
\begin{enumerate}
  \item $v_p(xy) = v_p(x) + v_p(y)$;
  \item $v_p(x+y) \geq \min\bigl(v_p(x),\, v_p(y)\bigr)$, with equality when
    $v_p(x) \neq v_p(y)$.
\end{enumerate}
\end{proposition}

\begin{proof}
Property~(1) follows directly from unique factorisation. For~(2), write
$x = p^a \cdot (r/s)$ and $y = p^b \cdot (u/v)$ where $p \nmid ruv$ and
$a = v_p(x)$, $b = v_p(y)$. Without loss of generality assume $a \leq b$.
Then $x + y = p^a\bigl(r/s + p^{b-a}\cdot u/v\bigr)$. The term in
parentheses has $p$-adic valuation $\geq 0$, so
$v_p(x+y) \geq a = \min(v_p(x), v_p(y))$. If $a < b$, the term $r/s$ has
$v_p = 0$ while $p^{b-a}\cdot u/v$ has $v_p \geq 1$, so the sum has $v_p = 0$
and equality holds.
\end{proof}

%------------------------------------------------------------------
\section{The \texorpdfstring{$p$}{p}-adic Absolute Value and the Ultrametric Inequality}
\label{sec:p-adic-abs}
\index{p-adic absolute value@$p$-adic absolute value}
\index{ultrametric inequality}
%------------------------------------------------------------------

\begin{definition}[$p$-adic absolute value]
\label{def:p-adic-abs}
The \emph{$p$-adic absolute value} on $\Q$ is defined by
\[
  \abs{x}_p =
  \begin{cases}
    p^{-v_p(x)} & \text{if } x \neq 0,\\
    0            & \text{if } x = 0.
  \end{cases}
\]
\end{definition}

\begin{example}[Sizes reversed]
\label{ex:p-adic-sizes}
In the $5$-adic world, $\abs{5^3}_5 = 5^{-3} = 1/125$ is very small, while
$\abs{1/5^3}_5 = 5^3 = 125$ is very large. High divisibility by~$p$ makes a
number \emph{small}, the opposite of the Archimedean intuition.
\end{example}

\begin{theorem}[Ultrametric inequality]
\label{thm:ultrametric}
\index{ultrametric inequality!proof}
For all $x, y \in \Q$,
\[
  \abs{x + y}_p \leq \max\bigl(\abs{x}_p,\, \abs{y}_p\bigr).
\]
Moreover, equality holds whenever $\abs{x}_p \neq \abs{y}_p$.
\end{theorem}

\begin{proof}
If $x = 0$ or $y = 0$ the result is immediate. Otherwise, by
Proposition~\ref{prop:vp-properties}(2),
\[
  v_p(x+y) \geq \min\bigl(v_p(x),\, v_p(y)\bigr).
\]
Applying $p^{-(\cdot)}$ (a decreasing function) to both sides reverses the
inequality:
\[
  \abs{x+y}_p = p^{-v_p(x+y)} \leq p^{-\min(v_p(x),\, v_p(y))}
  = \max\bigl(p^{-v_p(x)},\, p^{-v_p(y)}\bigr)
  = \max\bigl(\abs{x}_p,\, \abs{y}_p\bigr).
\]
The equality statement follows from the corresponding equality in
Proposition~\ref{prop:vp-properties}(2).
\end{proof}

\begin{remark}[Consequences of ultrametricity]
\label{rem:ultrametric-consequences}
\index{ultrametric inequality!consequences}
The ultrametric inequality is far stronger than the usual triangle inequality
$\abs{x+y}\leq\abs{x}+\abs{y}$. Among its striking consequences:
\begin{itemize}
  \item Every triangle in $\Q_p$ is isosceles: if $\abs{x}_p \neq \abs{y}_p$
    then $\abs{x+y}_p = \max(\abs{x}_p, \abs{y}_p)$.
  \item Every point of an open ball is its centre.
  \item Two open balls are either disjoint or one contains the other.
\end{itemize}
\end{remark}

%------------------------------------------------------------------
\section{The \texorpdfstring{$p$}{p}-adic Integers and \texorpdfstring{$p$}{p}-adic Numbers}
\label{sec:Zp-Qp}
\index{p-adic integers@$p$-adic integers $\Z_p$}
\index{p-adic numbers@$p$-adic numbers $\Q_p$}
%------------------------------------------------------------------

\begin{definition}[$p$-adic integers and $p$-adic numbers]
\label{def:Zp-Qp}
The \emph{$p$-adic numbers} $\Q_p$ are the completion of $\Q$ with respect to
the $p$-adic absolute value $\abs{\cdot}_p$. The \emph{$p$-adic integers} are
the closed unit ball:
\[
  \Z_p = \bigl\{ x \in \Q_p : \abs{x}_p \leq 1 \bigr\}
       = \bigl\{ x \in \Q_p : v_p(x) \geq 0 \bigr\}.
\]
\end{definition}

Every element $\alpha \in \Z_p$ admits a unique representation as a
\emph{$p$-adic expansion}\index{p-adic expansion@$p$-adic expansion}:
\[
  \alpha = a_0 + a_1\, p + a_2\, p^2 + a_3\, p^3 + \cdots,
  \quad a_i \in \{0, 1, \ldots, p-1\}.
\]
This converges in $\abs{\cdot}_p$ because $\abs{a_n p^n}_p \leq p^{-n} \to 0$.
A general element of $\Q_p$ has the form
$\alpha = \sum_{n=k}^{\infty} a_n\, p^n$ for some $k \in \Z$ (possibly negative).

\begin{example}[The element $-1$ in $\Z_p$]
\label{ex:minus-one-p-adic}
In $\Z_p$, we have
\[
  -1 = (p-1) + (p-1)\,p + (p-1)\,p^2 + \cdots = \sum_{n=0}^{\infty}(p-1)\,p^n,
\]
since the partial sums are $\sum_{n=0}^{N}(p-1)\,p^n = p^{N+1} - 1$, and
$\abs{p^{N+1}-1 - (-1)}_p = \abs{p^{N+1}}_p = p^{-(N+1)} \to 0$.
\end{example}

\begin{proposition}[Structure of $\Z_p$]
\label{prop:Zp-structure}
\index{p-adic integers@$p$-adic integers $\Z_p$!structure}
\begin{enumerate}
  \item $\Z_p$ is a local ring\index{local ring} with unique maximal ideal $p\Z_p$.
  \item The units of $\Z_p$ are $\Z_p^{\times} = \{ x \in \Z_p : \abs{x}_p = 1 \}
    = \{ x \in \Z_p : v_p(x) = 0 \}$.
  \item $\Z_p$ is compact in the $p$-adic topology.
  \item $\Q_p = \Z_p[1/p]$, i.e., every $p$-adic number has the form $p^{-k}\alpha$
    for some $\alpha \in \Z_p$ and $k \geq 0$.
\end{enumerate}
\end{proposition}

%------------------------------------------------------------------
\section{Hensel's Lemma}
\label{sec:hensel}
\index{Hensel's lemma|(}
%------------------------------------------------------------------

Hensel's lemma is the $p$-adic analogue of Newton's method: a simple root
modulo~$p$ can be ``lifted'' to a root in $\Z_p$.

\begin{theorem}[Hensel's lemma, simple case]
\label{thm:hensel}
Let $f(x) \in \Z_p[x]$ be a polynomial. Suppose there exists $a_0 \in \Z_p$
such that
\[
  \abs{f(a_0)}_p < \abs{f'(a_0)}_p^2.
\]
Then there exists a unique $a \in \Z_p$ with $f(a) = 0$ and
$\abs{a - a_0}_p \leq \abs{f(a_0)}_p / \abs{f'(a_0)}_p < \abs{f'(a_0)}_p$.
\end{theorem}

\begin{proof}
We construct a sequence $(a_n)$ by the $p$-adic Newton iteration:
\[
  a_{n+1} = a_n - \frac{f(a_n)}{f'(a_n)}.
\]
Set $\varepsilon = \abs{f(a_0)}_p / \abs{f'(a_0)}_p^2 < 1$ and
$\delta = \abs{f(a_0)}_p / \abs{f'(a_0)}_p$.

\medskip\noindent
\textbf{Step 1: The derivative stays bounded.}
We claim $\abs{f'(a_n)}_p = \abs{f'(a_0)}_p$ for all~$n$. Indeed, by
Taylor expansion,
\[
  f'(a_1) = f'(a_0) + (a_1 - a_0)\, g(a_0, a_1)
\]
for some polynomial $g$. Since $\abs{a_1 - a_0}_p = \delta < \abs{f'(a_0)}_p$
and $\abs{g(\ldots)}_p \leq 1$, the ultrametric inequality gives
$\abs{f'(a_1)}_p = \abs{f'(a_0)}_p$. Inductively this holds for all~$n$.

\medskip\noindent
\textbf{Step 2: Rapid convergence.}
By Taylor's theorem for polynomials,
\[
  f(a_{n+1}) = f(a_n) + f'(a_n)(a_{n+1}-a_n) + (a_{n+1}-a_n)^2\, h
\]
for some $h \in \Z_p$. Since $a_{n+1} - a_n = -f(a_n)/f'(a_n)$, the first two
terms cancel, giving
\[
  \abs{f(a_{n+1})}_p \leq \abs{a_{n+1}-a_n}_p^2 = \frac{\abs{f(a_n)}_p^2}{\abs{f'(a_0)}_p^2}.
\]
This yields the bound $\abs{f(a_{n+1})}_p \leq \varepsilon \cdot \abs{f(a_n)}_p$,
so $\abs{f(a_n)}_p \to 0$ at least as fast as $\varepsilon^{2^n - 1}$.

\medskip\noindent
\textbf{Step 3: Convergence and uniqueness.}
The sequence $(a_n)$ is Cauchy in $\Q_p$: for $m > n$,
$\abs{a_m - a_n}_p \leq \max_{k=n}^{m-1}\abs{a_{k+1}-a_k}_p \to 0$.
By completeness, $a_n \to a \in \Z_p$, and continuity of~$f$ gives $f(a) = 0$.
Uniqueness in the stated ball follows from the ultrametric inequality: if
$f(a) = f(b) = 0$ with both in the ball, then
$0 = f(a) - f(b) = (a-b)\bigl(f'(a) + (a-b)h\bigr)$; since $\abs{a-b}_p < \abs{f'(a)}_p$,
the factor in brackets is a unit, forcing $a = b$.
\end{proof}

\begin{corollary}[Practical lifting criterion]
\label{cor:hensel-practical}
If $f(x) \in \Z[x]$ and $a_0 \in \Z$ satisfies $f(a_0) \equiv 0 \pmod{p}$
with $f'(a_0) \not\equiv 0 \pmod{p}$, then there exists a unique
$a \in \Z_p$ with $f(a) = 0$ and $a \equiv a_0 \pmod{p}$.
\end{corollary}

\index{Hensel's lemma|)}

%------------------------------------------------------------------
\section{Applications of Hensel's Lemma}
\label{sec:hensel-apps}
%------------------------------------------------------------------

\begin{example}[Solving $x^2 = -1$ in $\Q_5$]
\label{ex:sqrt-minus1-Q5}
\index{p-adic numbers@$p$-adic numbers!square root of $-1$}
Consider $f(x) = x^2 + 1$. We look for a root modulo~$5$: trying $a_0 = 2$
gives $f(2) = 5 \equiv 0 \pmod{5}$, and $f'(2) = 4 \not\equiv 0 \pmod{5}$.
By Corollary~\ref{cor:hensel-practical}, there exists $i \in \Z_5$ with
$i^2 = -1$ and $i \equiv 2 \pmod{5}$.

We can compute the first few digits. Newton's iteration gives:
\begin{align*}
  a_0 &= 2,\\
  a_1 &= a_0 - \frac{f(a_0)}{f'(a_0)} = 2 - \frac{5}{4}
       = 2 - 5\cdot 4^{-1} \equiv 2 - 5 \cdot 4 \equiv 2 + 5\cdot 1
       = 7 \pmod{25},\\
  a_2 &\equiv 7 - \frac{50}{14} \equiv 7 - 50 \cdot 14^{-1} \pmod{125}.
\end{align*}
Continuing, one obtains the $5$-adic expansion
$i = 2 + 1\cdot 5 + 2\cdot 5^2 + 1\cdot 5^3 + \cdots \in \Z_5$.
Note that $-i$ is also a root, with expansion starting $3 + 3\cdot 5 + 2\cdot 5^2 + \cdots$
\end{example}

\begin{remark}[When does $\sqrt{-1}$ exist in $\Q_p$?]
\label{rem:sqrt-minus1-Qp}
A square root of $-1$ exists in $\Q_p$ if and only if $p \equiv 1 \pmod{4}$
or $p = 2$ (but the case $p=2$ requires care: $\sqrt{-1}$ does not exist
in $\Q_2$). This is closely related to our earlier study of quadratic residues
(Chapter~6).
\end{remark}

%------------------------------------------------------------------
\section{The Local-Global Principle}
\label{sec:local-global}
\index{local-global principle}
\index{Hasse--Minkowski theorem}
%------------------------------------------------------------------

One of the most powerful ideas in number theory is that solving an equation
over $\Q$ can be reduced to solving it over each ``local'' completion
$\Q_p$ (for every prime~$p$) and over $\R$.

\begin{theorem}[Hasse--Minkowski]
\label{thm:hasse-minkowski}
\index{Hasse--Minkowski theorem}
A quadratic form $Q(\mathbf{x}) = \sum_{i,j} a_{ij}\, x_i x_j$ with
$a_{ij} \in \Q$ represents zero nontrivially over~$\Q$ (i.e., there exists
$\mathbf{x} \neq \mathbf{0}$ in $\Q^n$ with $Q(\mathbf{x})=0$) if and only if
it represents zero nontrivially over~$\R$ and over $\Q_p$ for every prime~$p$.
\end{theorem}

\begin{remark}[Limitations of the local-global principle]
\label{rem:local-global-limitations}
The Hasse--Minkowski theorem holds for quadratic forms but \emph{fails} in
general. A famous counterexample due to Selmer\index{Selmer} is the cubic curve
$3x^3 + 4y^3 + 5z^3 = 0$, which has nontrivial solutions in $\R$ and in every
$\Q_p$, but no nontrivial solution in $\Q$. Understanding when and why the
local-global principle fails is a central theme in modern arithmetic geometry,
involving the Brauer--Manin obstruction\index{Brauer--Manin obstruction}.
\end{remark}

%------------------------------------------------------------------
\section{Ostrowski's Theorem}
\label{sec:ostrowski}
\index{Ostrowski's theorem}
%------------------------------------------------------------------

The $p$-adic absolute values, together with the usual absolute value, exhaust
all possibilities for measuring size on $\Q$.

\begin{theorem}[Ostrowski's theorem]
\label{thm:ostrowski}
Every nontrivial absolute value on $\Q$ is equivalent to either:
\begin{enumerate}
  \item the usual (Archimedean) absolute value $\abs{\cdot}_\infty$, or
  \item the $p$-adic absolute value $\abs{\cdot}_p$ for some prime~$p$.
\end{enumerate}
Here, two absolute values are \emph{equivalent}\index{equivalent absolute values}
if they induce the same topology on $\Q$.
\end{theorem}

\begin{remark}[Product formula]
\label{rem:product-formula}
\index{product formula}
For every $x \in \Q^*$,
\[
  \abs{x}_\infty \cdot \prod_{p \text{ prime}} \abs{x}_p = 1.
\]
This elegant identity encodes the idea that the ``places'' of $\Q$ --- one for
each prime and one Archimedean --- together capture all arithmetic information.
\end{remark}

%------------------------------------------------------------------
\section{Visualising \texorpdfstring{$p$}{p}-adic Structure}
\label{sec:p-adic-tikz}
%------------------------------------------------------------------

The $p$-adic integers have a tree-like structure: truncating the $p$-adic
expansion to $n$ digits gives a map $\Z_p \to \Z/p^n\Z$, and the preimages
of each residue class form a ball of radius $p^{-n}$.

\begin{figure}[ht]
\centering
\begin{tikzpicture}[
  level distance=1.4cm,
  level 1/.style={sibling distance=4cm},
  level 2/.style={sibling distance=1.3cm},
  every node/.style={circle, draw, inner sep=2pt, minimum size=6mm, font=\small},
  edge from parent/.style={draw, -}
]
\node {$\Z_3$}
  child { node {$0$}
    child { node {$0$} }
    child { node {$3$} }
    child { node {$6$} }
  }
  child { node {$1$}
    child { node {$1$} }
    child { node {$4$} }
    child { node {$7$} }
  }
  child { node {$2$}
    child { node {$2$} }
    child { node {$5$} }
    child { node {$8$} }
  };
\end{tikzpicture}
\caption{The tree structure of $\Z_3$: the first two levels correspond to
residues modulo~$3$ and modulo~$9$. Each path from root to leaf determines
the first digits of a $3$-adic expansion.}
\label{fig:3-adic-tree}
\end{figure}

\begin{figure}[ht]
\centering
\begin{tikzpicture}
  % Outer disk: Z_3
  \fill[blue!8] (0,0) circle (3cm);
  \draw[blue!60, thick] (0,0) circle (3cm);
  \node[blue!70] at (0, 3.35) {$\Z_3$: $\abs{x}_3 \leq 1$};

  % Three sub-disks for residues mod 3
  \fill[red!10] (0, 1.2) circle (0.95cm);
  \draw[red!60, thick] (0, 1.2) circle (0.95cm);
  \node[red!70, font=\small] at (0, 1.2) {$\equiv 0\ (3)$};

  \fill[green!10] (-1.1, -0.8) circle (0.95cm);
  \draw[green!60!black, thick] (-1.1, -0.8) circle (0.95cm);
  \node[green!60!black, font=\small] at (-1.1, -0.8) {$\equiv 1\ (3)$};

  \fill[orange!10] (1.1, -0.8) circle (0.95cm);
  \draw[orange!70, thick] (1.1, -0.8) circle (0.95cm);
  \node[orange!70, font=\small] at (1.1, -0.8) {$\equiv 2\ (3)$};

  % Annotation
  \node[font=\footnotesize, text width=3.5cm, align=center] at (0, -3.6)
    {In $\Z_3$, balls of radius $1/3$\\partition the unit ball into 3 clopen disks.};
\end{tikzpicture}
\caption{Disk structure of $\Z_3$. The unit ball decomposes into three disjoint
balls of radius $1/3$, each consisting of elements with a given residue
modulo~$3$.}
\label{fig:3-adic-disks}
\end{figure}

%------------------------------------------------------------------
\section{Exercises}
\label{sec:p-adic-exercises}
%------------------------------------------------------------------

\begin{exercise}
\label{exer:vp-factorial}
\index{Legendre's formula}
Prove \emph{Legendre's formula}: for any prime $p$ and positive integer $n$,
\[
  v_p(n!) = \sum_{k=1}^{\infty} \floor{\frac{n}{p^k}}.
\]
Use this to compute $v_2(100!)$ and $v_5(100!)$.
\end{exercise}

\begin{exercise}
\label{exer:ultrametric-isosceles}
Let $(K, \abs{\cdot})$ be a field with an ultrametric absolute value. Prove
that in any triangle with side lengths $a$, $b$, $c$ (i.e., $a = \abs{x-y}$,
etc.), the two largest sides have equal length.
\end{exercise}

\begin{exercise}
\label{exer:hensel-sqrt2}
\index{Hensel's lemma!application}
Determine for which primes $p \leq 13$ the equation $x^2 = 2$ has a solution
in $\Z_p$. For $p=7$, find the first three digits of the $7$-adic expansion
of $\sqrt{2}$.
\end{exercise}

\begin{exercise}
\label{exer:p-adic-series}
Show that the series $\sum_{n=0}^{\infty} n!\,$ converges in $\Q_p$ for every
prime~$p$. What is its $p$-adic absolute value?
\end{exercise}

\begin{exercise}
\label{exer:Zp-compact}
Prove that $\Z_p$ is compact by showing it is a closed and totally bounded
subset of the complete metric space $\Q_p$.
\end{exercise}

\begin{exercise}
\label{exer:product-formula-verify}
Verify the product formula $\abs{x}_\infty \prod_p \abs{x}_p = 1$ for
$x = 75/14$.
\end{exercise}

\begin{exercise}
\label{exer:minus-one-expansion}
Write out the $7$-adic expansion of $-1$ (first five digits) and verify your
answer by computing the partial sums.
\end{exercise}

\begin{exercise}[Challenging]
\label{exer:Zp-not-field}
Show that $\Z_p$ is not a field but that $\Q_p$ is. Identify the non-units of
$\Z_p$ explicitly.
\end{exercise}

%------------------------------------------------------------------
\section{Chapter Summary}
\label{sec:p-adic-summary}
%------------------------------------------------------------------

\begin{itemize}
  \item The \textbf{$p$-adic valuation} $v_p$ measures divisibility by~$p$;
    the $p$-adic absolute value $\abs{x}_p = p^{-v_p(x)}$ makes highly divisible
    numbers \emph{small}.
  \item The \textbf{ultrametric inequality}
    $\abs{x+y}_p \leq \max(\abs{x}_p, \abs{y}_p)$ leads to exotic but useful
    topology: every point is a centre, and all triangles are isosceles.
  \item The completion $\Q_p$ contains $\Z_p$ (the unit ball), whose elements
    have $p$-adic expansions analogous to power series.
  \item \textbf{Hensel's lemma} lifts simple roots modulo~$p$ to exact roots
    in~$\Z_p$, providing a powerful tool for solving polynomial equations.
  \item \textbf{Ostrowski's theorem} classifies all absolute values on $\Q$:
    the usual one and the $p$-adic ones.
  \item The \textbf{Hasse--Minkowski theorem} exemplifies the local-global
    principle for quadratic forms, connecting $\Q_p$-solutions for all~$p$ and
    $\R$-solutions to $\Q$-solutions.
\end{itemize}

\index{p-adic numbers@$p$-adic numbers|)}


%%%%%%%%%%%%%%%%%%%%%%%%%%%%%%%%%%%%%%%%%%%%%%%%%%%%%%%%%%%%%
%% CHAPTER 10: Preview of Analytic Number Theory
%%%%%%%%%%%%%%%%%%%%%%%%%%%%%%%%%%%%%%%%%%%%%%%%%%%%%%%%%%%%%
\chapter{Preview of Analytic Number Theory}
\label{ch:analytic-nt}
\index{analytic number theory|(}

\begin{quote}
\emph{``The primes are the atoms of arithmetic, and the zeta function is the
  spectrograph that reveals their hidden structure.''}
\end{quote}

\noindent
Throughout this course, we have studied the primes using algebraic and
combinatorial methods. In this final chapter, we catch a glimpse of
\emph{analytic number theory}\index{analytic number theory}, where the tools of
calculus and complex analysis are brought to bear on the distribution of primes.
The central object is the Riemann zeta function, and the central result is the
Prime Number Theorem.

%------------------------------------------------------------------
\section{The Riemann Zeta Function}
\label{sec:zeta}
\index{Riemann zeta function|(}
\index{zeta function|see{Riemann zeta function}}
%------------------------------------------------------------------

\begin{definition}[Riemann zeta function]
\label{def:zeta}
For a complex variable $s$ with $\operatorname{Re}(s) > 1$, the
\emph{Riemann zeta function}\index{Riemann zeta function!definition} is
\[
  \zeta(s) = \sum_{n=1}^{\infty} \frac{1}{n^s}.
\]
\end{definition}

\begin{proposition}[Convergence]
\label{prop:zeta-convergence}
\index{Riemann zeta function!convergence}
The series $\sum_{n=1}^{\infty} n^{-s}$ converges absolutely for
$\operatorname{Re}(s) > 1$ and uniformly on every half-plane
$\operatorname{Re}(s) \geq 1 + \delta$ with $\delta > 0$.
\end{proposition}

\begin{proof}
For $s = \sigma + it$ with $\sigma > 1$, we have
$\abs{n^{-s}} = n^{-\sigma}$, and $\sum n^{-\sigma}$ converges by the
integral test since $\int_1^\infty x^{-\sigma}\,dx = 1/(\sigma-1) < \infty$.
Uniform convergence on $\sigma \geq 1+\delta$ follows from the Weierstrass
$M$-test with $M_n = n^{-(1+\delta)}$.
\end{proof}

%------------------------------------------------------------------
\section{Euler Product}
\label{sec:euler-product}
\index{Euler product|(}
%------------------------------------------------------------------

The deep connection between $\zeta(s)$ and the primes is encoded in the
following identity, discovered by Euler in 1737 for real $s > 1$ and extended
by Riemann to the complex half-plane.

\begin{theorem}[Euler product formula]
\label{thm:euler-product}
\index{Euler product!proof}
For $\operatorname{Re}(s) > 1$,
\[
  \zeta(s) = \prod_{p\ \mathrm{prime}} \frac{1}{1 - p^{-s}}.
\]
\end{theorem}

\begin{proof}
We give the classical sieve argument. For $\operatorname{Re}(s) > 1$,
each factor $(1 - p^{-s})^{-1} = \sum_{k=0}^{\infty} p^{-ks}$ converges
absolutely.

Let $p_1 = 2, p_2 = 3, p_3 = 5, \ldots$ be the primes in order. Consider
the finite product:
\[
  P_N(s) = \prod_{j=1}^{N} \frac{1}{1 - p_j^{-s}}
  = \prod_{j=1}^{N} \sum_{k_j=0}^{\infty} p_j^{-k_j s}.
\]
Expanding the product, we obtain a sum over all $n$ of the form
$n = p_1^{k_1} p_2^{k_2} \cdots p_N^{k_N}$, i.e., all positive integers whose
prime factors are among $\{p_1, \ldots, p_N\}$:
\[
  P_N(s) = \sum_{\substack{n \geq 1 \\ p \mid n \Rightarrow p \leq p_N}}
  \frac{1}{n^s}.
\]
By the fundamental theorem of arithmetic, each such $n$ appears exactly once.

Now we estimate the difference:
\[
  \abs{\zeta(s) - P_N(s)}
  = \left\lvert \sum_{\substack{n \geq 1 \\ \exists\, p > p_N,\; p \mid n}}
    \frac{1}{n^s} \right\rvert
  \leq \sum_{n > p_N} \frac{1}{n^{\sigma}}.
\]
Since $\sigma > 1$, this tail of a convergent series tends to~$0$ as
$N \to \infty$. Hence $P_N(s) \to \zeta(s)$, establishing the product formula.
\end{proof}

\begin{corollary}[$\zeta(s) \neq 0$ for $\operatorname{Re}(s)>1$]
\label{cor:zeta-nonzero}
\index{Riemann zeta function!nonvanishing}
Since each factor $(1-p^{-s})^{-1}$ is nonzero for $\operatorname{Re}(s) > 1$
and the product converges absolutely, we have $\zeta(s) \neq 0$ in this
half-plane.
\end{corollary}

\begin{corollary}[Infinitude of primes, analytic proof]
\label{cor:infinitude-analytic}
\index{primes!infinitude!analytic proof}
Since $\zeta(s) \to \infty$ as $s \to 1^+$ (the harmonic series diverges), the
Euler product $\prod_p (1-p^{-1})^{-1}$ must be an infinite product, hence there
are infinitely many primes.
\end{corollary}

\index{Euler product|)}

%------------------------------------------------------------------
\section{Special Value: \texorpdfstring{$\zeta(2) = \pi^2/6$}{zeta(2) = pi squared over 6}}
\label{sec:zeta-2}
\index{Riemann zeta function!special values}
\index{Basel problem}
%------------------------------------------------------------------

The evaluation of $\zeta(2)$ was a famous open problem (the \emph{Basel
problem}\index{Basel problem}) solved by Euler in 1735.

\begin{theorem}[Basel problem]
\label{thm:zeta-2}
$\displaystyle\zeta(2) = \sum_{n=1}^{\infty} \frac{1}{n^2} = \frac{\pi^2}{6}.$
\end{theorem}

\begin{proof}[Sketch of proof]
Consider the function $f(x) = x^2$ on $[-\pi, \pi]$. Its Fourier
series\index{Fourier series} is
\[
  x^2 = \frac{\pi^2}{3} + 4\sum_{n=1}^{\infty} \frac{(-1)^n}{n^2}\cos(nx).
\]
Setting $x = \pi$, and using $\cos(n\pi) = (-1)^n$, we obtain:
\[
  \pi^2 = \frac{\pi^2}{3} + 4\sum_{n=1}^{\infty} \frac{1}{n^2},
\]
which gives $\sum_{n=1}^{\infty} n^{-2} = \pi^2/6$.

An alternative elegant proof uses the Weierstrass product for
$\sin(\pi x)/(\pi x)$:
\[
  \frac{\sin(\pi x)}{\pi x} = \prod_{n=1}^{\infty}\left(1 - \frac{x^2}{n^2}\right).
\]
Expanding both sides and comparing the coefficient of~$x^2$ yields the result.
\end{proof}

%------------------------------------------------------------------
\section{Dirichlet \texorpdfstring{$L$}{L}-Functions}
\label{sec:dirichlet-L}
\index{Dirichlet L-function@Dirichlet $L$-function|(}
\index{Dirichlet character}
%------------------------------------------------------------------

To study primes in arithmetic progressions, Dirichlet generalised the zeta
function using \emph{characters}\index{Dirichlet character}.

\begin{definition}[Dirichlet character]
\label{def:dirichlet-character}
A \emph{Dirichlet character modulo~$q$} is a function $\chi\colon \Z \to \C$
satisfying:
\begin{enumerate}
  \item $\chi(n+q) = \chi(n)$ for all $n$ (periodicity);
  \item $\chi(mn) = \chi(m)\chi(n)$ for all $m,n$ (complete multiplicativity);
  \item $\chi(n) = 0$ if and only if $\gcd(n,q) > 1$.
\end{enumerate}
The \emph{principal character}\index{Dirichlet character!principal} $\chi_0$ has
$\chi_0(n)=1$ whenever $\gcd(n,q)=1$.
\end{definition}

\begin{definition}[Dirichlet $L$-function]
\label{def:dirichlet-L}
For a Dirichlet character $\chi$ and $\operatorname{Re}(s)>1$,
\[
  L(s, \chi) = \sum_{n=1}^{\infty} \frac{\chi(n)}{n^s}
  = \prod_{p\ \mathrm{prime}} \frac{1}{1 - \chi(p)\,p^{-s}}.
\]
\end{definition}

\begin{remark}[Relation to the zeta function]
\label{rem:L-to-zeta}
For the principal character $\chi_0$ modulo~$q$,
\[
  L(s, \chi_0) = \zeta(s) \prod_{p \mid q}(1 - p^{-s}),
\]
so $L(s, \chi_0)$ is essentially $\zeta(s)$ with finitely many Euler factors
removed.
\end{remark}

\index{Dirichlet L-function@Dirichlet $L$-function|)}

%------------------------------------------------------------------
\section{Dirichlet's Theorem on Primes in Arithmetic Progressions}
\label{sec:dirichlet-theorem}
\index{Dirichlet's theorem on primes in arithmetic progressions}
%------------------------------------------------------------------

\begin{theorem}[Dirichlet, 1837]
\label{thm:dirichlet-primes-AP}
If $a$ and $q$ are coprime positive integers, then the arithmetic progression
\[
  a,\quad a + q,\quad a + 2q,\quad a + 3q,\quad \ldots
\]
contains infinitely many primes. More precisely,
\[
  \sum_{\substack{p \leq x \\ p \equiv a \pmod{q}}} \frac{1}{p}
  \sim \frac{1}{\varphi(q)} \ln\ln x
  \quad\text{as } x \to \infty,
\]
where $\varphi$ is Euler's totient function.
\end{theorem}

\begin{remark}[Key idea of the proof]
\label{rem:dirichlet-proof-idea}
The proof uses the orthogonality of Dirichlet characters to isolate primes
in a given residue class: for $\gcd(a,q) = 1$,
\[
  \frac{1}{\varphi(q)} \sum_{\chi \bmod q} \overline{\chi(a)}\, \chi(n)
  = \begin{cases} 1 & \text{if } n \equiv a \pmod{q},\\ 0 & \text{otherwise}.\end{cases}
\]
The critical analytic input is that $L(1, \chi) \neq 0$ for all nonprincipal
characters~$\chi$. The proof of non-vanishing is the deepest part of the
argument, especially for real characters.
\end{remark}

%------------------------------------------------------------------
\section{The Prime Number Theorem}
\label{sec:PNT}
\index{Prime Number Theorem|(}
%------------------------------------------------------------------

\begin{definition}[Prime-counting function]
\label{def:pi-x}
\index{prime-counting function}
For $x > 0$, let $\pi(x) = \#\{p \leq x : p \text{ is prime}\}$.
\end{definition}

The Prime Number Theorem, conjectured independently by Legendre and Gauss in
the 1790s, gives the asymptotic law governing the distribution of primes.

\begin{theorem}[Prime Number Theorem]
\label{thm:PNT}
\index{Prime Number Theorem!statement}
As $x \to \infty$,
\[
  \pi(x) \sim \frac{x}{\ln x},
\]
meaning $\lim_{x\to\infty} \pi(x) / (x/\ln x) = 1$.
\end{theorem}

The PNT was proved independently by Hadamard\index{Hadamard, Jacques} and de la
Vall\'{e}e-Poussin\index{de la Vall\'ee-Poussin, Charles} in 1896. Both proofs
relied on showing that $\zeta(1+it) \neq 0$ for all real~$t \neq 0$ --- that is,
the zeta function has no zeros on the line $\operatorname{Re}(s) = 1$.

\begin{remark}[A better approximation: $\operatorname{Li}(x)$]
\label{rem:Li-approximation}
\index{logarithmic integral}
A significantly better approximation to $\pi(x)$ is the
\emph{logarithmic integral}:
\[
  \operatorname{Li}(x) = \int_2^x \frac{dt}{\ln t}.
\]
Gauss conjectured that $\pi(x) \sim \operatorname{Li}(x)$, and indeed the
PNT can be stated as $\pi(x) = \operatorname{Li}(x) + o(x/\ln x)$. The
function $\operatorname{Li}(x)$ typically overestimates $\pi(x)$ for ``small''
$x$ (up to about $10^{316}$), but Littlewood\index{Littlewood, J.\,E.} showed
that $\pi(x) - \operatorname{Li}(x)$ changes sign infinitely often.
\end{remark}

\begin{remark}[Historical significance]
\label{rem:PNT-history}
\index{Prime Number Theorem!history}
The PNT was one of the crowning achievements of 19th-century mathematics.
It demonstrated that deep questions about the discrete world of primes could
be answered using the continuous methods of complex analysis. In 1949--50,
Erd\H{o}s\index{Erdos, Paul@Erd\H{o}s, Paul} and Selberg\index{Selberg, Atle}
gave ``elementary'' (i.e., not using complex analysis) proofs, though these
are far from simple.
\end{remark}

%------------------------------------------------------------------
\section{Chebyshev Functions}
\label{sec:chebyshev}
\index{Chebyshev functions|(}
%------------------------------------------------------------------

Chebyshev\index{Chebyshev, Pafnuty} introduced two functions that are technically
more convenient than $\pi(x)$ for studying the distribution of primes.

\begin{definition}[Chebyshev functions]
\label{def:chebyshev}
\begin{align*}
  \theta(x) &= \sum_{p \leq x} \ln p, \\
  \psi(x)   &= \sum_{p^k \leq x} \ln p = \sum_{n \leq x} \Lambda(n),
\end{align*}
where $\Lambda(n)$\index{von Mangoldt function} is the \emph{von Mangoldt
function}: $\Lambda(n) = \ln p$ if $n = p^k$ for some prime $p$ and $k \geq 1$,
and $\Lambda(n) = 0$ otherwise.
\end{definition}

\begin{proposition}[Equivalence with PNT]
\label{prop:chebyshev-PNT}
The following are equivalent:
\begin{enumerate}
  \item $\pi(x) \sim x/\ln x$;
  \item $\theta(x) \sim x$;
  \item $\psi(x) \sim x$.
\end{enumerate}
\end{proposition}

\begin{remark}[Why $\psi(x)$ is preferred]
\label{rem:psi-preferred}
The function $\psi(x)$ has a beautiful explicit formula connecting it to the
zeros of $\zeta(s)$:
\[
  \psi(x) = x - \sum_{\rho} \frac{x^\rho}{\rho} - \ln(2\pi) - \frac{1}{2}\ln(1-x^{-2}),
\]
where the sum is over the nontrivial zeros $\rho$ of $\zeta(s)$. This formula
makes precise how the distribution of primes is governed by the zeros of the
zeta function.
\end{remark}

\index{Chebyshev functions|)}

%------------------------------------------------------------------
\section{The Riemann Hypothesis}
\label{sec:RH}
\index{Riemann Hypothesis|(}
%------------------------------------------------------------------

The zeta function can be analytically continued\index{analytic continuation}
to all of $\C \setminus \{1\}$ (with a simple pole at $s=1$). It satisfies
the \emph{functional equation}\index{Riemann zeta function!functional equation}
\[
  \pi^{-s/2}\,\Gamma\!\left(\frac{s}{2}\right)\zeta(s)
  = \pi^{-(1-s)/2}\,\Gamma\!\left(\frac{1-s}{2}\right)\zeta(1-s).
\]
The zeros at $s = -2, -4, -6, \ldots$ are called \emph{trivial zeros}%
\index{Riemann zeta function!trivial zeros}. All other zeros lie in the
\emph{critical strip}\index{critical strip} $0 \leq \operatorname{Re}(s) \leq 1$.

\begin{definition}[The Riemann Hypothesis]
\label{def:RH}
The \emph{Riemann Hypothesis}\index{Riemann Hypothesis!statement} (RH) asserts
that every nontrivial zero of $\zeta(s)$ has real part equal to $1/2$. That is,
all nontrivial zeros lie on the \emph{critical line}\index{critical line}
$\operatorname{Re}(s) = 1/2$.
\end{definition}

\begin{remark}[Consequences of the Riemann Hypothesis]
\label{rem:RH-consequences}
\index{Riemann Hypothesis!consequences}
If RH is true, the error in the Prime Number Theorem becomes remarkably sharp:
\[
  \pi(x) = \operatorname{Li}(x) + O\bigl(\sqrt{x}\,\ln x\bigr).
\]
Among the many other consequences:
\begin{itemize}
  \item Sharp bounds on gaps between consecutive primes.
  \item The Lindelöf hypothesis\index{Lindelöf hypothesis} on the growth of
    $\zeta(1/2+it)$.
  \item Improved estimates in the distribution of primes in arithmetic
    progressions.
  \item Implications for the distribution of eigenvalues of random
    matrices\index{random matrix theory}.
\end{itemize}
\end{remark}

\begin{remark}[The Millennium Prize]
\label{rem:millennium}
\index{Riemann Hypothesis!Millennium Prize}
The Riemann Hypothesis is one of the seven Millennium Prize
Problems\index{Millennium Prize Problems} selected by the Clay Mathematics
Institute in 2000, each carrying a prize of \$1{,}000{,}000. As of this
writing, the RH remains unresolved, despite enormous computational
evidence\index{Riemann Hypothesis!numerical evidence} (over $10^{13}$ nontrivial
zeros have been verified to lie on the critical line) and more than 160 years
of effort by the world's best mathematicians since Riemann's 1859 memoir.
\end{remark}

%------------------------------------------------------------------
\section{Visualising Prime Distribution and the Critical Strip}
\label{sec:analytic-tikz}
%------------------------------------------------------------------

\begin{figure}[ht]
\centering
\begin{tikzpicture}
\begin{axis}[
  width=12cm, height=7.5cm,
  xlabel={$x$},
  ylabel={},
  xmin=2, xmax=100,
  ymin=0, ymax=30,
  legend pos=north west,
  legend style={font=\small},
  grid=major,
  grid style={gray!30},
  samples=50,
  domain=2:100,
  thick
]
% pi(x) as a step function using a few key data points
\addplot[blue, very thick, const plot mark left] coordinates {
  (2,1) (3,2) (5,3) (7,4) (11,5) (13,6) (17,7) (19,8) (23,9) (29,10)
  (31,11) (37,12) (41,13) (43,14) (47,15) (53,16) (59,17) (61,18) (67,19)
  (71,20) (73,21) (79,22) (83,23) (89,24) (97,25) (100,25)
};
\addlegendentry{$\pi(x)$}

% x / ln(x)
\addplot[red, dashed, thick] {x / ln(x)};
\addlegendentry{$x/\ln x$}

% Li(x) approximation via 1/ln(t) (using the antiderivative approximation)
% Li(x) ~ x/ln(x) + x/(ln x)^2 + ...
% We use a reasonable approximation
\addplot[green!60!black, densely dotted, thick] {x/ln(x) + x/(ln(x))^2};
\addlegendentry{$\operatorname{Li}(x)$ (approx.)}

\end{axis}
\end{tikzpicture}
\caption{Comparison of the prime-counting function $\pi(x)$ (blue step function)
with the estimates $x/\ln x$ (red, dashed) and $\operatorname{Li}(x)$ (green, dotted).
The logarithmic integral provides a visibly superior approximation.}
\label{fig:pi-x-comparison}
\end{figure}

\begin{figure}[ht]
\centering
\begin{tikzpicture}[scale=1.1]
  % Critical strip
  \fill[blue!8] (0,-3.5) rectangle (1,3.5);
  \draw[blue!50, thick] (0,-3.5) -- (0,3.5);
  \draw[blue!50, thick] (1,-3.5) -- (1,3.5);
  \node[blue!60, font=\footnotesize, rotate=90] at (-0.3, 0) {$\operatorname{Re}(s)=0$};
  \node[blue!60, font=\footnotesize, rotate=90] at (1.3, 0) {$\operatorname{Re}(s)=1$};

  % Critical line
  \draw[red, very thick] (0.5,-3.5) -- (0.5,3.5);
  \node[red, font=\footnotesize, rotate=90] at (0.7, -2.8) {$\operatorname{Re}(s)=\tfrac{1}{2}$};

  % Some zeros on the critical line
  \foreach \y in {-2.94, -2.13, -1.52, -0.93, -0.51, 0.51, 0.93, 1.52, 2.13, 2.94} {
    \fill[red] (0.5, \y) circle (2pt);
  }

  % Pole at s=1
  \node[draw, circle, inner sep=1.5pt, thick, orange!80!black] at (1, 0) {};
  \node[orange!80!black, font=\footnotesize, anchor=west] at (1.15, 0.15) {pole at $s=1$};

  % Axes
  \draw[->, thick] (-1.5,0) -- (2.5,0) node[right] {$\operatorname{Re}(s)$};
  \draw[->, thick] (0.5,-3.5) -- (0.5, 3.8) node[above] {$\operatorname{Im}(s)$};

  % Trivial zeros
  \foreach \x in {-0.5, -1.0} {
    \fill[gray] (\x, 0) circle (2pt);
  }
  \node[gray, font=\footnotesize, anchor=north] at (-0.75, -0.15) {trivial zeros};

  % Labels
  \node[blue!70, font=\small] at (0.5, 3.5) [above] {\textbf{Critical strip}};

  % Bracket for critical strip
  \draw[decorate, decoration={brace, amplitude=5pt, raise=2pt}, blue!60]
    (0, -3.5) -- (1, -3.5) node[midway, below=8pt, font=\footnotesize, blue!70]
    {$0 \leq \operatorname{Re}(s) \leq 1$};
\end{tikzpicture}
\caption{The critical strip and the critical line of the Riemann zeta function.
The nontrivial zeros (red dots) are shown on the critical line
$\operatorname{Re}(s) = 1/2$, as predicted by the Riemann Hypothesis.
Trivial zeros (grey dots) lie at $s = -2, -4, \ldots$ on the real axis.}
\label{fig:critical-strip}
\end{figure}

%------------------------------------------------------------------
\section{Open Problems in Number Theory}
\label{sec:open-problems}
\index{open problems}
%------------------------------------------------------------------

We conclude this course by listing several famous unsolved problems, each of
which has resisted the efforts of mathematicians for decades or centuries.

\begin{definition}[Goldbach's conjecture]
\label{def:goldbach}
\index{Goldbach's conjecture}
\emph{Goldbach's conjecture} (1742) asserts that every even integer greater
than~$2$ can be expressed as the sum of two primes.
\end{definition}

Despite extensive computational verification (up to at least $4\times 10^{18}$),
Goldbach's conjecture remains unproven. The best partial result is due to
Helfgott\index{Helfgott, Harald} (2013), who proved the \emph{weak} (or ternary)
Goldbach conjecture\index{Goldbach's conjecture!weak}: every odd integer greater
than~$5$ is the sum of three primes.

\begin{definition}[Twin prime conjecture]
\label{def:twin-primes}
\index{twin prime conjecture}
The \emph{twin prime conjecture} asserts that there are infinitely many pairs
of primes $(p, p+2)$.
\end{definition}

A major breakthrough was achieved by Zhang\index{Zhang, Yitang} (2013), who
proved that there are infinitely many pairs of primes differing by at most
$70{,}000{,}000$. This bound has been reduced to~$246$ through the collaborative
Polymath project\index{Polymath project}, but the gap to~$2$ remains.

\begin{remark}[Other notable open problems]
\label{rem:other-open}
\begin{itemize}
  \item \textbf{Riemann Hypothesis}\index{Riemann Hypothesis}: discussed in
    Section~\ref{sec:RH}; the most important open problem in mathematics.
  \item \textbf{Birch and Swinnerton-Dyer conjecture}%
    \index{Birch and Swinnerton-Dyer conjecture}: relates the rank of an
    elliptic curve to the order of vanishing of its $L$-function at $s=1$.
    Another Millennium Prize problem.
  \item \textbf{Legendre's conjecture}\index{Legendre's conjecture}: there is
    always a prime between $n^2$ and $(n+1)^2$.
  \item \textbf{Are there infinitely many Mersenne primes?}%
    \index{Mersenne primes}: primes of the form $2^p - 1$. As of 2024, 51
    Mersenne primes are known.
  \item \textbf{Collatz conjecture}\index{Collatz conjecture}: for the map
    $n \mapsto n/2$ (if even) or $3n+1$ (if odd), does every positive integer
    eventually reach~$1$?
\end{itemize}
\end{remark}

%------------------------------------------------------------------
\section{Exercises}
\label{sec:analytic-exercises}
%------------------------------------------------------------------

\begin{exercise}
\label{exer:euler-product-compute}
Verify the Euler product by computing $\prod_{p \leq 7}(1-p^{-2})^{-1}$ and
comparing it with $\sum_{n=1}^{N} n^{-2}$ for a suitable~$N$. How close is
your answer to $\pi^2/6$?
\end{exercise}

\begin{exercise}
\label{exer:zeta-even}
\index{Riemann zeta function!even values}
Using the identity $\zeta(2) = \pi^2/6$, show that $\zeta(4) = \pi^4/90$.
\emph{Hint:} use the Fourier series of $x^4$ on $[-\pi,\pi]$, or the identity
$\zeta(4) = (1/3)\bigl(\zeta(2)^2 + \zeta(2)\bigr) - \text{correction}$.
Alternatively, use the product formula for $\sin(\pi x)/(\pi x)$ and compare
the $x^4$ coefficient.
\end{exercise}

\begin{exercise}
\label{exer:dirichlet-3-mod4}
Use Dirichlet's theorem to show that there are infinitely many primes of the
form $4k+3$. Can you give an elementary proof (without $L$-functions)?
\end{exercise}

\begin{exercise}
\label{exer:PNT-consequences}
Assuming the PNT, show that the $n$-th prime $p_n$ satisfies
$p_n \sim n \ln n$ as $n \to \infty$.
\end{exercise}

\begin{exercise}
\label{exer:mertens-theorem}
\index{Mertens' theorem}
\emph{Mertens' theorem} states that $\sum_{p\leq x}1/p = \ln\ln x + M + o(1)$,
where $M \approx 0.2615$ is the Meissel--Mertens constant. Use this to show
that the product $\prod_{p\leq x}(1-1/p) \sim e^{-\gamma}/\ln x$, where
$\gamma$ is the Euler--Mascheroni constant.
\end{exercise}

\begin{exercise}
\label{exer:chebyshev-bounds}
\index{Chebyshev functions!bounds}
Prove that for all $x \geq 2$,
\[
  \frac{1}{6}\,\frac{x}{\ln x} \leq \pi(x) \leq 6\,\frac{x}{\ln x}.
\]
\emph{Hint:} consider $\binom{2n}{n}$ and use the fact that every prime
$n < p \leq 2n$ divides $\binom{2n}{n}$.
\end{exercise}

\begin{exercise}
\label{exer:von-mangoldt}
\index{von Mangoldt function}
Show that for all $n \geq 1$, $\sum_{d \mid n}\Lambda(d) = \ln n$.
Use Möbius inversion to deduce $\Lambda(n) = -\sum_{d\mid n}\mu(d)\ln d$.
\end{exercise}

\begin{exercise}
\label{exer:zeta-1-nonzero}
\index{Riemann zeta function!nonvanishing on $\operatorname{Re}(s)=1$}
Show that $\zeta(s)$ has a simple pole at $s=1$ with residue~$1$:
$\lim_{s\to 1}(s-1)\zeta(s) = 1$. \emph{Hint:} compare $\sum n^{-s}$ with
$\int_1^\infty x^{-s}\,dx$.
\end{exercise}

\begin{exercise}
\label{exer:goldbach-verify}
\index{Goldbach's conjecture}
Verify Goldbach's conjecture for all even integers from~$4$ to~$100$.
\end{exercise}

\begin{exercise}[Challenging]
\label{exer:dirichlet-class-number}
\index{class number formula}
Look up Dirichlet's class number formula and explain how $L(1,\chi)$ for a
quadratic character $\chi$ relates to the class number of a quadratic number
field. What role does the non-vanishing $L(1,\chi) \neq 0$ play?
\end{exercise}

%------------------------------------------------------------------
\section{Chapter Summary}
\label{sec:analytic-summary}
%------------------------------------------------------------------

\begin{itemize}
  \item The \textbf{Riemann zeta function}
    $\zeta(s) = \sum n^{-s} = \prod_p (1-p^{-s})^{-1}$
    encodes the distribution of primes via its Euler product.
  \item The \textbf{Euler product} is proved by expanding geometric series and
    using unique factorisation; it gives an analytic proof that there are
    infinitely many primes.
  \item \textbf{Dirichlet $L$-functions} $L(s,\chi)$ generalise the zeta
    function and underpin Dirichlet's theorem that every coprime arithmetic
    progression contains infinitely many primes.
  \item The \textbf{Prime Number Theorem} states $\pi(x) \sim x/\ln x$; the
    better approximation $\operatorname{Li}(x) = \int_2^x dt/\ln t$ captures
    finer behaviour.
  \item The \textbf{Chebyshev functions} $\theta(x)$ and $\psi(x)$ provide
    technically convenient equivalents of the PNT, with $\psi(x)$ directly
    connected to the zeros of $\zeta(s)$.
  \item The \textbf{Riemann Hypothesis} --- that all nontrivial zeros of $\zeta$
    lie on $\operatorname{Re}(s) = 1/2$ --- remains the most important open
    problem in mathematics, with profound implications for prime distribution.
  \item Famous open problems including \textbf{Goldbach's conjecture},
    the \textbf{twin prime conjecture}, and the \textbf{RH} illustrate that
    number theory remains a vibrant and active field of research.
\end{itemize}

\index{analytic number theory|)}

%%%%%%%%%%%%%%%%%%%%%%%%%%%%%%%%%%%%%%%%%%%%%%%%%%%%%%%%%%%%%
%% END MATTER
%%%%%%%%%%%%%%%%%%%%%%%%%%%%%%%%%%%%%%%%%%%%%%%%%%%%%%%%%%%%%

\backmatter

%------------------------------------------------------------------
% BIBLIOGRAPHY
%------------------------------------------------------------------
\begin{thebibliography}{99}
\label{sec:bibliography}

\bibitem{apostol}
\index{Apostol, Tom M.}
T.\,M.\ Apostol,
\emph{Introduction to Analytic Number Theory},
Undergraduate Texts in Mathematics,
Springer-Verlag, New York, 1976.

\bibitem{davenport}
\index{Davenport, Harold}
H.\ Davenport,
\emph{Multiplicative Number Theory},
3rd ed., revised by H.\,L.\ Montgomery,
Graduate Texts in Mathematics~74,
Springer-Verlag, New York, 2000.

\bibitem{hardy-wright}
\index{Hardy, G. H.}\index{Wright, E. M.}
G.\,H.\ Hardy and E.\,M.\ Wright,
\emph{An Introduction to the Theory of Numbers},
6th ed., revised by D.\,R.\ Heath-Brown and J.\,H.\ Silverman,
Oxford University Press, Oxford, 2008.

\bibitem{ireland-rosen}
\index{Ireland, Kenneth}\index{Rosen, Michael}
K.\ Ireland and M.\ Rosen,
\emph{A Classical Introduction to Modern Number Theory},
2nd ed., Graduate Texts in Mathematics~84,
Springer-Verlag, New York, 1990.

\bibitem{koblitz}
\index{Koblitz, Neal}
N.\ Koblitz,
\emph{$p$-adic Numbers, $p$-adic Analysis, and Zeta-Functions},
2nd ed., Graduate Texts in Mathematics~58,
Springer-Verlag, New York, 1984.

\bibitem{nathanson}
\index{Nathanson, Melvyn B.}
M.\,B.\ Nathanson,
\emph{Elementary Methods in Number Theory},
Graduate Texts in Mathematics~195,
Springer-Verlag, New York, 2000.

\bibitem{nzm}
\index{Niven, Ivan}\index{Zuckerman, Herbert S.}\index{Montgomery, Hugh L.}
I.\ Niven, H.\,S.\ Zuckerman, and H.\,L.\ Montgomery,
\emph{An Introduction to the Theory of Numbers},
5th ed., John Wiley \& Sons, New York, 1991.

\bibitem{serre}
\index{Serre, Jean-Pierre}
J.-P.\ Serre,
\emph{A Course in Arithmetic},
Graduate Texts in Mathematics~7,
Springer-Verlag, New York, 1973.

\end{thebibliography}

%------------------------------------------------------------------
% INDEX
%------------------------------------------------------------------
\printindex

\end{document}
